\documentclass[9pt,a5paper]{extbook}

% Encoding
\usepackage[utf8]{inputenc}
\usepackage[T1]{fontenc}
\usepackage{cmap}

% Fonts
\usepackage{amssymb}
\usepackage{newtxtext,newtxmath}
\usepackage{microtype}
\usepackage{lettrine}
\renewcommand{\LettrineFontHook}{\usefont{T1}{ptm}{m}{n}}

% Quotes
\usepackage{csquotes}
\usepackage{epigraph}
\setlength{\epigraphwidth}{0.8\textwidth}
\renewcommand{\epigraphflush}{flushright}

% Layout
\usepackage[top=1.2cm,bottom=1.2cm,left=1.1cm,right=0.9cm,headheight=12pt,headsep=0.3cm]{geometry}
\usepackage{multicol,array,longtable,booktabs,ragged2e}

% Spacing
\usepackage{setspace}
\setstretch{1.08}
\setlength{\parskip}{0.4ex plus 0.2ex minus 0.1ex}
\setlength{\parindent}{1em}
\frenchspacing

\widowpenalty=1000
\clubpenalty=1000
\tolerance=1200
\emergencystretch=3em

% Colors
\usepackage{xcolor}
\definecolor{darkblue}{RGB}{0,45,90}
\definecolor{patronblue}{RGB}{48,63,96}
\definecolor{warningred}{RGB}{139,26,26}
\definecolor{ritualgold}{RGB}{184,134,11}
\definecolor{spiritblue}{RGB}{70,130,180}
\definecolor{softgray}{RGB}{245,245,245}
\definecolor{bordergray}{RGB}{128,128,128}

% Boxes
\usepackage[most]{tcolorbox}
\tcbuselibrary{breakable}
\tcbset{breakable}
\tcbuselibrary{skins,breakable}

\newtcolorbox{patronbox}[2][]{%
    enhanced, colback=softgray, colframe=patronblue,
    fonttitle=\bfseries\small\sffamily\color{patronblue},
    fontupper=\small, fontlower=\small,
    title={$\blacklozenge$ #2 $\blacklozenge$},
    attach boxed title to top left={yshift=-1.5mm,xshift=3mm},
    boxed title style={colback=white,colframe=patronblue,boxrule=0.5pt},
    sharp corners, boxrule=0.6pt,
    left=3pt, right=3pt, top=4pt, bottom=3pt,
    breakable, #1}

\newtcolorbox{ritualnote}{%
    enhanced, colback=ritualgold!8, colframe=ritualgold,
    before skip=0.2cm, after skip=0.2cm,
    sharp corners, boxrule=0.5pt,
    left=3pt, right=3pt,
    fonttitle=\bfseries\small\sffamily,
    fontupper=\small, title={Ritual Note}}

\newtcolorbox{warningbox}[1][]{%
    enhanced, colback=warningred!5, colframe=warningred,
    before skip=0.2cm, after skip=0.2cm,
    sharp corners, boxrule=0.6pt,
    left=3pt, right=3pt,
    fonttitle=\bfseries\small\sffamily\color{warningred},
    fontupper=\small, title={Warning}, #1}

\newtcolorbox{spellcard}[2][]{%
    breakable, enhanced, colback=softgray, colframe=bordergray,
    fonttitle=\bfseries\small\sffamily\color{darkblue},
    fontupper=\small, title={#2},
    attach boxed title to top left={yshift=-2mm,xshift=3mm},
    boxed title style={colback=white,colframe=bordergray,boxrule=0.5pt},
    sharp corners, boxrule=0.6pt,
    left=4pt, right=4pt, top=5pt, #1}

\newtcolorbox{spiritcard}[2][]{%
    breakable, enhanced, colback=spiritblue!8, colframe=spiritblue,
    fonttitle=\bfseries\small\sffamily\color{spiritblue!60!black},
    fontupper=\small, title={#2},
    attach boxed title to top left={yshift=-2mm,xshift=3mm},
    boxed title style={colback=spiritblue!20,colframe=spiritblue,boxrule=0.5pt},
    sharp corners, boxrule=0.6pt,
    left=4pt, right=4pt, top=5pt, #1}

% Headers
\usepackage{fancyhdr}
\pagestyle{fancy}
\fancyhf{}
\fancyhead[LE]{\thepage\quad\textsc{\leftmark}}
\fancyhead[RO]{\textsc{\rightmark}\quad\thepage}
\renewcommand{\headrulewidth}{0.3pt}
\renewcommand{headrule}{\hbox to\headwidth{\color{darkblue}\leaders\hrule height \headrulewidth\hfill}}

% Simple section formatting (NO titlesec to avoid conflicts)
% Using standard LaTeX plus sectsty if needed, or just standard

% If you MUST use titlesec, use this minimal config:
% Remove titlesec \titleformat calls for section and subsection, keep only chapter:
%\usepackage{titlesec}
%\titleformat{\chapter}[display]
%  {\normalfont\Large\bfseries\color{darkblue}\centering}
%  {\chaptertitlename\ \thechapter}{0.5em}{\titlerule[0.5pt]\vspace{0.5em}}[\vspace{0.5em}\titlerule[0.5pt]]

% Use sectsty for sections/subsections (bulletproof)
%\usepackage{sectsty}
%\sectionfont{\normalfont\normalsize\bfseries\color{darkblue}}
%\subsectionfont{\normalfont\small\bfseries\itshape\color{darkblue!80}}


% NO custom appendix redefinition - let it use normal chapter style
\usepackage{tocloft}
\renewcommand{\cftchapfont}{\bfseries\color{darkblue}\small}
\renewcommand{\cftsecfont}{\small}
\setlength{\cftbeforechapskip}{4pt}

% Tables
\usepackage{etoolbox}
\AtBeginEnvironment{longtable}{\small\sffamily\RaggedRight\setlength{\tabcolsep}{3pt}}
\AtBeginEnvironment{tabular}{\small\RaggedRight\setlength{\tabcolsep}{3pt}}

% Commands
\newcommand{\abilitytag}[1]{\textsc{\scriptsize #1}}
\newcommand{\dangerous}[1]{\textcolor{warningred}{\textbf{#1}}}

\newcommand{\spell}[5]{%
    \begin{spellcard}{#1}
    \textbf{\scriptsize Tags:} #2 \\[0.1em]
    \textbf{\scriptsize DV:} #3 \\[0.1em]
    \textbf{\scriptsize Effect:} #4 \\[0.1em]
    \textbf{\scriptsize Backlash (1+ SB):} #5
    \end{spellcard}
    \vspace{0.1cm}}

% Hyperlinks last
\usepackage[colorlinks=true,linkcolor=darkblue,citecolor=darkblue,urlcolor=darkblue,
    pdftitle={Fate's Edge - Complete Grimoire},pdfauthor={The Grey Wanderer}]{hyperref}

\begin{document}

\frontmatter
\thispagestyle{empty}

% Title Page - Scaled for A5
\begin{center}
\vspace*{1.5cm} % Reduced from 2.5cm

{\LARGE\scshape\color{darkblue} The Grey Wanderer's Grimoire}\\[0.8cm] % Reduced from \fontsize{28}{32}

{\normalsize\itshape A Complete Guide to Cantors, Invokers, Runekeepers, and Summoners}\\[1.0cm] % Reduced from \Large

\rule{0.5\textwidth}{1.0pt}\\[0.8cm]

{\small\textit{``You want power? Fine. But don't come crying when the weight finds you.''}}\\[1.5cm] % Reduced from \large

\vfill

{\footnotesize\textsc{Compiled for PDFLaTeX}\\[0.2em]
\textit{With booktabs, microtype, and malice.}}

\end{center}

\clearpage

% Table of Contents
\tableofcontents
\clearpage

\mainmatter

\chapter{Introduction -- A Voice from the Crossroads}
\epigraph{``You want power? Fine. But don't come crying when the weight finds you.''}{\textit{--- The Grey Wanderer}}

This grimoire is not a rulebook. It is a \textbf{companion}---a set of lenses through which to see the magical paths of \textit{Fate's Edge}. You will not find rigid orders or chartered guilds herein. There are no grandmasters to bow before, no academies whose diplomas protect you from the flame. What follows are \textit{approaches}---six roads that diverge from the crossroads where power and price shake hands.

In Thepyrgos, a Synod clerk once called me a \textit{psy-ee-on}---some term from the Vhasian badlands for those who hear thoughts without asking. I told him I do not read minds; I read the \textit{weight} of what people carry. He wrote me down anyway, in his register of freaks and witch-marked. He was not entirely wrong, but he missed the point by a mile. The powers I have touched leave fingerprints that look like many things. I have worn the chains of a dozen Patrons and broken most of them. Twice. I have been called a Runekeeper in Aeler, a witch in Aelaerem, a Cantor in Silkstrand, a Summoner in Zakov, and a fool everywhere. The road changes the walker, and I have walked them all.

We will explore:

\begin{itemize}
    \item \textbf{The Six Roads}---not merely the three paths of covenant, but the ways of the Runekeeper, the Invoker, the Cantor, the Witch, the Psion, and the Summoner. Who they are, how they think, and why they risk annihilation to keep walking.
    \item How to handle \textbf{Obligation}---not as a punishment, but as a story engine. You will find patron-specific ideas, intrusions, and acts of service, as well as the price of walking without a patron at all.
    \item \textbf{Rituals without a focus}---when your Codex is ash, your Symbol cracked, your voice alone, and your mind turned inward. Components matter again when the ledger is blank.
    \item The art of \textbf{Summoning}---binding spirits to will, from the deep stone to the salt foam, and the terrible etiquette required to send them back.
\end{itemize}

The voice you hear belongs to no one and everyone---a wanderer who has stood at the crossroads of every Patron's attention, who has tasted the Hollow Tithe and climbed the Silent Stair, who has learned that the hedge is what keeps the wolves from the flock. Call me the \textbf{Grey Wanderer}. These notes are my legacy. Use them, lose them, or burn them for warmth. I won't mind.

Now close the book. Go make some mistakes. I will be at the inn, buying the wine with a coin that shouldn't exist, waiting to hear what road you choose.

\vspace{0.5cm}
\hfill --- \textit{The Grey Wanderer}

\clearpage
\chapter{The Weight of Names -- How the Peoples of Fate's Edge See Magic}
\epigraph{``I have carried my power through a dozen cultures and been called a dozen things: saint, sinner, witch, warlock, fraud, fool. Every people has a word for what I do. None of them are quite right. None of them are quite wrong.''}{\textit{--- The Grey Wanderer}}

Magic in Fate's Edge is not a single practice. It is a family of arguments with reality---and every culture argues differently. Some embrace the structured pacts of Patrons. Others whisper to spirits at crossroads. A few have learned to bend the world with nothing but a disciplined mind. What follows is not a definitive guide---I have been corrected by village grandmothers and high priests alike---but a wanderer's impressions, gathered over too many years on too many roads.

\section{Aeler -- The Deep Accountants}
\lettrine{T}{he} dwarves of Aeler do not speak of ``magic.'' They speak of ``load,'' ``breath,'' and ``covenant.'' A Rite is not a miracle---it is a contract with stone, sealed by proper measure and witnessed by a bell. A Runekeeper is an engineer of the invisible, and a witch is a mender of cracks in the world's infrastructure.

\subsection*{What They Respect}
\begin{itemize}
    \item \textbf{Runekeeper:} The ideal Aeler mage. A Codex carved in slate, a Thiasos that chitters in the dark, and a precise understanding of Obligation. The mountain respects a bond you can count.
    \item \textbf{Witch (Hearth craft):} The women who keep the cisterns sweet and the gas vents predictable are not feared---they are salaried. Aeler employs witches as ``Breath Wardens'' and ``Cistern Mothers.''
\end{itemize}

\subsection*{What They Mistrust}
\begin{itemize}
    \item \textbf{Cantor:} Too public. Too emotional. The deep stone does not like loud noises it cannot measure.
    \item \textbf{Free Caster:} An abomination. Power without a covenant is a cave-in waiting to happen.
    \item \textbf{Psion:} Rare, but when one appears, the Aeler assign them to the Deep Vaults---where their mind-reading is useful and their instability is contained.
\end{itemize}

\subsection*{The Breath-Bond}
Every Aeler mage, regardless of path, must learn the \textbf{Breath Bond}: a meditation that converts psychic strain into measured breaths. A Psion in Aeler carries a bell on a leather cord. Each Mental Strain point is one bell-ring. When the bell rings three times too fast, a Corrector appears.

\subsection*{Famous Aeler Magic}
The \textbf{Spine Underway}---a tunnel that remembers every traveler who has passed through it. Runekeepers maintain it by carving new Rites into the walls each century. Witches keep the air sweet. And the stone itself? It is said to be the oldest Summoner, awake and patient.

\section{Aelaerem -- The Hearth-and-Hollow}
\lettrine{T}{he} halflings of Aelaerem do not build temples to Patrons. They leave butter on cup-mark stones and hang red thread on hawthorn. Their magic is domestic---and therefore invisible to those who look only for grand gestures.

\subsection*{What They Respect}
\begin{itemize}
    \item \textbf{Witch:} The Aelaerem are the heartland of witchcraft. Every village has a Hedge-Mother who knows which knot to tie for a safe birth and which bell to ring to keep the Hollow from swallowing a child.
    \item \textbf{Cantor:} Folk songs are the oldest grimoire. A Cantor who sings the old tunes and respects the ``no ninth'' rule is welcome at any hearth.
    \item \textbf{Summoner:} But only brownies, field-wights, and the small dead. Aelaerem summoners do not call demons---they ask the ancestors to move the firewood.
\end{itemize}

\subsection*{What They Mistrust}
\begin{itemize}
    \item \textbf{Invoker:} Carrying Symbols from multiple Patrons is seen as greedy. The Hollow does not forgive divided loyalty.
    \item \textbf{Psion:} ``Reading thoughts without asking'' is a violation of hospitality. A known psion will find doors closing and cup-mark stones left dry.
\end{itemize}

\subsection*{The Hollow Tithe}
Once a year, the Aelaerem pay the \textbf{Hollow Tithe}: a basket of the year's best bread, left at the edge of a barrow. Witches conduct the rite. No one speaks during the walk back. I have done it three times. Something was always following. Something always stopped at the village gate.

\section{Aelinnel -- The Bright Weird}
\lettrine{T}{he} gnomes of Aelinnel live in a world where numbers have manners and doors sometimes lead to last Tuesday. Their magic is not chaotic---it is \textit{algebraic}. They have names for equations that other cultures do not know exist.

\subsection*{What They Respect}
\begin{itemize}
    \item \textbf{Free Caster:} The gnomes invented the TAGS system, though they call it ``Contextual Recursion.'' A free caster who can improvise a spell on the spot is a philosopher, not a danger.
    \item \textbf{Psion:} The mind is the ultimate context. Aelinnel psions are trained from childhood to ``count the thoughts'' and ``measure the echo.''
    \item \textbf{Summoner:} Spirits are just data from a previous iteration of reality. Treat them politely and document everything.
\end{itemize}

\subsection*{What They Mistrust}
\begin{itemize}
    \item \textbf{Runekeeper:} Too rigid. A Rite carved in stone cannot adapt to a changing context. The gnomes prefer the flexibility of witch-knots and psionic whispers.
    \item \textbf{Cantor:} A song that cannot be written in two ledgers (Said and Meant) is dangerous. The gnomes have banned three hymns in the last decade.
\end{itemize}

\subsection*{The Green Gate}
The most famous Aelinnel threshold---a sea-arch that opens at certain tides to a forest not on any map. Witches and Free Casters use it to travel between realities. Psions use it to test their mental resilience. The gnomes charge a toll: one truth you have never told anyone, and one lie you have never been caught telling.

\section{Linn -- The Storm-Tested}
\lettrine{T}{he} northmen of Linn do not ask permission from Patrons. They make oaths---and the sea witnesses. Their magic is carved into keels, braided into ropes, and whispered into the ears of drowning men.

\subsection*{What They Respect}
\begin{itemize}
    \item \textbf{Summoner (Selkie and Drowned):} A Linn tide-witch who can call a selkie for navigation or a drowned captain for warning is worth her weight in whalebone.
    \item \textbf{Cantor:} Sea shanties are weather magic. A Cantor who sings the right note at the right moment can calm a storm---or call one.
    \item \textbf{Witch (Wake-Judge):} The Linn have formalized witchcraft into a legal system. Wake-Judges arbitrate disputes at the water's edge, using promise timers to bind oaths.
\end{itemize}

\subsection*{What They Mistrust}
\begin{itemize}
    \item \textbf{Invoker:} Borrowing power from a Patron is seen as weakness. A real northman makes his own luck.
    \item \textbf{Psion:} Too quiet. The Linn prefer their magic loud enough to hear over the wind.
\end{itemize}

\subsection*{The Thing-Bell}
Every Linn longship has a bell forged from the melted-down sword of the ship's first captain. When a Runekeeper or Witch needs to seal an oath, they ring the Thing-Bell. The sound carries across water. The dead listen.

\section{Ykrul -- The Weather Readers}
\lettrine{T}{he} orcs of the Violet Steppe do not separate magic from survival. A bone-singer who calls a storm is no different from a scout who reads the grass. Both are necessary. Both can get you killed.

\subsection*{What They Respect}
\begin{itemize}
    \item \textbf{Summoner (Ancestral):} The drum is the oldest Codex. A Ykrul summoner who can call a wolf-ancestor or a wind-horse is honored at any campfire.
    \item \textbf{Psion (Precognition):} The Ykrul call psions ``Weather Readers''---not because they predict rain, but because they see the shape of the future in the same way a rider sees the shape of the steppe.
    \item \textbf{Witch (Road craft):} A witch who can bind a promise with a knotted cord or curse a rival with a handful of grave-dust is a valued member of any clan.
\end{itemize}

\subsection*{What They Mistrust}
\begin{itemize}
    \item \textbf{Runekeeper:} Too slow. The steppe does not wait for you to carve a Rite.
    \item \textbf{Invoker:} Carrying Symbols from multiple Patrons is seen as divided loyalty. A Ykrul warrior serves one banner.
    \item \textbf{Cantor:} But only because a song that binds a crowd is a leash. The Ykrul value freedom too highly.
\end{itemize}

\subsection*{The Red Weather}
When a Ykrul witch curses a horse-thief, they do not use subtlety. They tie a red cord around a splinter of the thief's saddle and burn it in a storm. The ``Red Weather''---a sudden, localized squall---is said to follow the cursed until they return what they stole. I have seen it happen twice. I have no explanation.

\section{Vhasia -- The Fractured Lens}
\lettrine{T}{he} Vhasians inherited the ruins of the Utaran Empire and learned to argue over the pieces. Their magic is legalistic, factional, and surprisingly bureaucratic.

\subsection*{What They Respect}
\begin{itemize}
    \item \textbf{Runekeeper:} A Vhasian Runekeeper is a notary with teeth. Their Codex contains Rites, yes, but also contracts, property deeds, and the terms of three unresolved feuds.
    \item \textbf{Invoker:} Multiple Patrons? In Vhasia, that is called ``diversification.'' A noble who keeps Symbols from Aveh (for escape), Mykkiel (for contracts), and the Witness (for truth) is merely prudent.
    \item \textbf{Witch (Cunning craft):} Every Vhasian village has a hedge-witch who speaks for the old saints. The Church of the Everflame pretends not to notice.
\end{itemize}

\subsection*{What They Mistrust}
\begin{itemize}
    \item \textbf{Free Caster:} Too unpredictable. The Vhasian legal system cannot bill a free caster for damages if there is no contract.
    \item \textbf{Psion:} Reading minds without permission is a crime---and the Vhasian Parlement has written three laws against it. Enforcement is spotty.
\end{itemize}

\subsection*{The Parlement of Dust}
In the catacombs beneath the Sun Palace, a secret court of Runekeepers adjudicates disputes between magical factions. They meet by candlelight and seal their judgments with wax mixed with grave-dust. No one admits the court exists. Everyone with real power obeys its rulings.

\section{Viterra -- The Hedge-Law Kingdom}
\lettrine{T}{he} Viterrans are practical people. Their magic grows in hedgerows and hides in parish ledgers. A witch in Viterra is not a mystic---she is a ``map adjuster'' who knows where the boundary stones should really be.

\subsection*{What They Respect}
\begin{itemize}
    \item \textbf{Witch (Threshold):} The Viterran ``hedge-wright'' is a legal as well as magical profession. They draw property lines, witness oaths, and ensure that the new mill does not cross the ancient right-of-way.
    \item \textbf{Runekeeper:} But only those who specialize in ``precedent Rites''---spells that make a past judgment binding on the present.
    \item \textbf{Cantor:} A village Cantor who sings the old ballads at the harvest festival is beloved. One who sings politics is exiled.
\end{itemize}

\subsection*{What They Mistrust}
\begin{itemize}
    \item \textbf{Summoner:} Spirits cannot own property. The Viterran courts have ruled that a ghost's testimony is inadmissible.
    \item \textbf{Psion:} The Queen's Justiciars have a standing bounty on unlicensed telepaths. Privacy is the foundation of property law.
\end{itemize}

\subsection*{The Tilt-Yard Audits}
Before a Viterran Runekeeper can invoke a Rite that affects land ownership, they must ``audit'' the spell before a panel of Dawn-Knights. The knights strike the Codex with a sword. If the Rite survives three blows, it is lawful. I have seen a Codex shatter on the second strike. The Runekeeper was fined for ``magical noise pollution.''

\section{Ecktoria -- The Marble Ledger}
\lettrine{E}{cktoria} sits on the bones of the old Utaran Empire and pretends they are a throne. Their magic is ceremonial, political, and obsessed with appearances.

\subsection*{What They Respect}
\begin{itemize}
    \item \textbf{Runekeeper (Everflame):} The state religion has its own Runekeepers---the ``Censer-Knights''---who march in processions and purify heretics with blessed fire.
    \item \textbf{Invoker:} Ecktorian nobles collect Symbols the way they collect statues of forgotten emperors. Display is more important than use.
    \item \textbf{Cantor:} A Cantor who can make a crowd weep with a hymn is a political asset. The Grand Magistrate employs three.
\end{itemize}

\subsection*{What They Mistrust}
\begin{itemize}
    \item \textbf{Witch:} Hedge-craft is ``peasant magic.'' The Everflame preaches against it, but every noble's wife knows a witch who can brew a love potion.
    \item \textbf{Psion:} The Chain-Lanterns hunt psions with particular zeal. A mind that cannot be seen is a mind that might be planning rebellion.
    \item \textbf{Summoner:} The Ecktorian view on spirits is simple: if it is dead, it should stay dead. The Inquisition has burned three summoners in my lifetime.
\end{itemize}

\subsection*{The Censor's Mask}
At the annual Triumph, the Censor of Ecktoria wears a mask that is said to be a psionic artifact from the Great Silence. It lets him see the true intentions of every petitioner. I knelt before him once. He saw that I wanted his wine. He laughed and gave me a cup. The mask is real.

\section{Thepyrgos -- The City of Stairs}
\lettrine{T}{hepyrgos} is a nation of clerks and philosophers who believe that a properly argued case can bend reality. Their magic is academic, verbal, and exhausting.

\subsection*{What They Respect}
\begin{itemize}
    \item \textbf{Free Caster:} The University of Thepyrgos teaches the TAGS system as ``practical metaphysics.'' A free caster with a degree can charge triple.
    \item \textbf{Psion:} The Synod has a standing committee on ``mental hygiene.'' Licensed psions work as legal investigators. Unlicensed ones are hunted.
    \item \textbf{Witch (Bell craft):} The Thepyrgos bell-towers are thresholds. Bell-wardens are respected civic employees.
\end{itemize}

\subsection*{What They Mistrust}
\begin{itemize}
    \item \textbf{Runekeeper:} Too reliant on a single Patron. The Thepyrgos prefer negotiators, not worshippers.
    \item \textbf{Cantor:} A crowd that sings together is a crowd that might riot. The Archon limits public performances by Cantors to three per week.
\end{itemize}

\subsection*{The Silent Stair}
Behind the Synod Hall, a staircase that leads nowhere---but psions claim it leads to the space between thoughts. The University paid me a handsome sum to climb it and report back. I climbed for an hour. The stairs did not end. I turned around and walked down. I was back in the same courtyard. I do not climb it anymore.

\section{Lethai -- The Fair Folk's Echo}
\lettrine{T}{he} elves of the Valewood do not practice magic. They---and this is important---\textit{are} magic, or at least they have forgotten where they end and the Weave begins. A Lethai Runekeeper is a contradiction. A Lethai Cantor is a redundancy. A Lethai Witch is a priestess of a religion that has no name.

\subsection*{What They Respect}
\begin{itemize}
    \item \textbf{Witch (Silk and Shed):} The Lethai have two witch traditions: the Silk Vigil (Inaea, mercy, continuity) and the Shed Vigil (Isoka, transformation, hard choices). Both are woven into the fabric of Lethai law.
    \item \textbf{Summoner:} Spirits are just neighbors who live in a different room. The Lethai treat them with courtesy---and a healthy dose of suspicion.
    \item \textbf{Cantor:} A Lethai shira (song-weaver) can wrap a promise in a melody so beautiful that you forget you are bound. Their betrayal is legendary.
\end{itemize}

\subsection*{What They Mistrust}
\begin{itemize}
    \item \textbf{Invoker:} Why borrow power from a Patron when you could simply \textit{be} the power? The Lethai see Invokers as beggars in silk robes.
    \item \textbf{Psion:} Reading a Lethai mind without context is impossible---their thoughts are saturated with memory, season, and relationship. Psions who try often leave mad.
    \item \textbf{Free Caster:} Too crude. The Lethai prefer their magic to have rules, even if the rules are secret.
\end{itemize}

\subsection*{The Ninth Walk}
Once a generation, a Lethai witch walks the ``Ninth Walk''---a pilgrimage through thresholds that do not exist on any map. She returns with a new law for the Lethai courts. The last Ninth Walk produced the rule: ``No bell may ring nine times.'' No one remembers why. Everyone obeys.

\section{The Grey Wanderer's Conclusion}

I have been called a Runekeeper in Aeler, a witch in Aelaerem, a Cantor in Silkstrand, a Summoner in Zakov, a Psion in Thepyrgos, and a fool everywhere. The truth is simpler: every culture sees magic as it sees itself. The Aeler see a covenant. The Linn see an oath. The Thepyrgos see a debate. The Lethai see a dream.

You will find your own path. You will be misunderstood. You will be feared, honored, hunted, or hired---sometimes all in the same day. When that happens, remember this: the power is not in the Rite or the Song or the silent thought. The power is in the choice to use them.

Now close the book. Go make some mistakes. I will be at the inn.

\vspace{0.5cm}
\hfill --- \textit{The Grey Wanderer}

\clearpage
\chapter{The Six Roads -- Lore and Roleplay}
\epigraph{``I have worn the face of a knight, a warlock, a hedge-priest, and a hunted freak. The path chooses you as much as you choose it, and the price is always the same weight, though the coin differs.''}{\textit{--- The Grey Wanderer}}

There are not three paths through the thornwood, but six. Six roads that diverge from the crossroads where power and price shake hands. Do not mistake these for guilds or orders---there are no charters signed in ink, no grandmasters to bow before. These are \textit{approaches}, lenses through which the world's hidden gears become visible. A Runekeeper sworn to the Inquisitor Prime shares more with that burning zealot than with a Runekeeper of Grimmir, though both carry Codices. An Invoker of Aveh has more in common with a gambler drowning in risk than with a Witch who tends the same crossroads. The road shapes the walker.

\section{The Runekeeper -- Agent of the Inked Covenant}
\textit{``I do not beg. I record. And in the recording, I become the hand that writes the law.''}

The Runekeepers are the \textbf{agents} of the great powers---not merely users of magic, but its dedicated representatives, its priests and its executioners. To become a Runekeeper is to cease being merely yourself and become an extension of your Patron's will in the world. You are the covenant made flesh, the vow walking on two legs.

But do not think this creates uniformity. A Runekeeper is as varied as the Patron they serve:

\begin{itemize}
    \item \textbf{The Mykkielan Templar} serves the Stone-Judge with gavel and flame. He is paladin and notary both, his Codex bound in chain-mail, his Thiasos a hound of ash and seal-wax. He does not heal; he adjudicates wounds, declaring them lawful or heretical.
    \item \textbf{The Grimmiric Druid} is engineer of root and rot, his Codex carved into living bark, his Thiasos a fungal thing that speaks in spores. He maintains the treaty between the planted field and the wild wood, ensuring neither forgets its bounds.
    \item \textbf{The Artificer of Aelinnel} treats her Patron's Rites as formulae, her Codex a notebook of impossible geometries. She builds devices---timers that measure time backwards, lenses that see obligations as glowing threads---arguing that she does not perform miracles, but merely applies context.
    \item \textbf{The Inquisitor of the Everflame} carries a Codex of censure, each page a warrant for burning. Her Thiasos is a censer that never goes cold, and she is less mage than mobile courthouse, sentencing corruption to annihilation.
\end{itemize}

What unites them is structure. The Codex is not merely a tool but a \textit{contract made visible}, a covenant of service. The Thiasos is not a pet but a fragment of the Patron's attention, a bailiff sent to ensure the terms are kept. Lose the Codex, and you are a knight without armor. Lose the Thiasos, and you are a debtor whose creditor has just remembered where you live.

\subsection*{Roleplay Hooks}
\begin{itemize}
    \item Your Codex is incomplete---the last page was torn out by your predecessor. What was on it? Why did they fear it enough to destroy it?
    \item Your Thiasos grows more demanding. Last week it required a drop of blood. This week, it whispers for a finger.
    \item You have begun to dream in your Patron's voice. When you wake, you are never certain which of you made the decisions yesterday.
\end{itemize}

\section{The Invoker -- The Gambler at the Crossroads}
\textit{``I do not serve. I rent. And when the landlord comes to collect, I had best have the coin ready, or the interest will be paid in teeth.''}

If the Runekeeper is the sworn knight, the Invoker is the desperate cardsharp who has learned the house rules well enough to cheat---briefly, dangerously, at terrible interest. The Invoker carries \textbf{Symbols}, physical anchors to Patrons they have never fully sworn to. They are the gamblers of the magical world, but they do not bet on luck. They bet on \textit{knowledge}---the precise understanding of what each Patron requires, what each Symbol will bear, and exactly how far they can push before the weight becomes unbearable.

Their rites are slow, respectful, almost contractual in their precision. An Invoker does not rush the bargain. She polishes her Symbols. She leaves offerings at crossroads. She learns the true names of her Patrons' servants so she may address them properly when the time comes. This is the quiet work, the maintenance of the bond.

But when the knife is at the throat, when the building is burning, when the demon is already through the door---then she performs the \textbf{Crack the Seal}. She invokes the Rite not as a supplicant but as a creditor calling in an emergency loan. The price is never small. The Symbol may crack. The Patron may arrive personally to discuss terms. But for that moment, she holds power she has not earned, and the world bends.

\subsection*{Roleplay Hooks}
\begin{itemize}
    \item You carry four Symbols---the limit before friction. Every morning, roll them like bones. Which two whisper to each other today? Which one has begun to weep?
    \item You found a Symbol in a dead man's hand. You do not know which Patron it belongs to. You use it anyway. Sometimes, it works in your favor. Sometimes, it demands tasks you do not understand.
    \item You have never met your Patron. You are not certain they know you exist. But the weight grows anyway, and lately, you have seen their servants watching from across crowded rooms.
\end{itemize}

\section{The Cantor -- The Touched Voice}
\textit{``You think you need a lute? My larynx is older than any tree. Hum, and the world will listen. Scream, and it might answer back.''}

Cantors are the wild singers, the mad pipers, the hymn-leaders who do not worship at the altar but \textit{become} the altar. Their power is intimacy---unmediated, visceral, and deeply dangerous. They do not swear to Patrons as Runekeepers do. They do not bargain as Invokers do. They \textit{echo} the Patron's will through their own flesh, becoming resonant chambers for forces that were never meant to fit inside a mortal throat.

Because they do not swear, they cannot break an oath. But they can be \textbf{touched}. Corruption for a Cantor is not a debt to be paid but a \textit{transformation undergone}. The voice develops harmonics that should not exist. The breath carries scents from places that are not on any map. The eyes change color when they sing. Some Cantors grow feathers in their hair. Others find their shadows lagging half a step behind. They are corrupted not by obligation but by \textit{proximity}---standing too close to the fire for too long, singing the songs that ears were not meant to hear.

They require no focus but their own bodies. A Cantor without a voice can tap rhythm on their ribs. One without hands can whistle through their teeth. The body remembers the song even when the mind has forgotten it.

\subsection*{Roleplay Hooks}
\begin{itemize}
    \item You sang a forbidden note once, and now a fragment of a dead Cantor lives in your head, humming tunes that predict the weather three days hence.
    \item Your voice is unmistakable---enemies recognize you by your timbre, and allies find themselves weeping when you speak certain words, though they know not why.
    \item You are trying to write a new Song, but every attempt leaves you feverish for a week. The Song is writing \textit{you}, and it has plans for your vocal cords that you do not yet understand.
\end{itemize}

\section{The Witch -- The Quiet Architect}
\textit{``They call it hedge-magic. They call it petty. They do not see that the hedge is what keeps the wolves from the flock, the blight from the crop, the forgotten god from waking hungry.''}

Witches practice the \textbf{systemic magic} that maintains the world---the quiet, overlooked power that is at once invisible and essential. While Runekeepers swear and Cantors sing, Witches \textit{tend}. They are the ones who know that a threshold must be swept three times counter-timerwise to keep the Hollow from noticing it. They understand that milk sours not because of bacteria but because a promise was broken three houses down. They are the plumbers of the metaphysical, the ones who ensure that the pipes of reality do not burst.

This makes them feared more than loved. A Runekeeper's power is legible---there is a Codex, a Patron, a clear chain of command. A Witch's power is \textit{organic}, grown from relationships with places and spirits that have no names in any grimoire. They work with the Ninth, with thresholds, with the accumulated weight of stories. They are overlooked until the moment they are needed, and then they are feared, because they hold the keys to doors that others did not know existed.

Every culture has Witches, though they call them different things---Hedge-Mothers, Breath-Wardens, Map-Adjusters, Cistern-Keepers. They are the ones who remember that magic is not merely for throwing fireballs but for ensuring that the fire does not burn down the village.

\subsection*{Roleplay Hooks}
\begin{itemize}
    \item You owe bonds to three different spirits: the river that borders your village, the oak at the crossroads, and the crow that follows you. None have called them in yet, but they will.
    \item You cannot cross running water without making an offering. You do not remember why, but the one time you tried, you nearly drowned in three inches of rain.
    \item A village that once shunned you now begs for your help. The blight is back. You could save them, but the price they paid last time---your grandmother's name, erased from memory---still stings.
\end{itemize}

\section{The Psion -- The Isolated Mind}
\textit{``I carry no Symbol. I keep no Codex. My power has no scent, no sound, no outward sign. And that is why they fear me most of all.''}

Psions are the isolated ones, the untrusted, the inward-turned. Where other paths borrow from outside powers, the Psion looks only to the self---the disciplined, trained, dangerous self. They carry no outward signs of their power. There is no glowing staff, no familiar, no song to warn you. They are accountable only to themselves, and in a world built on bonds and covenants, this makes them suspect.

In Aeler, they are sealed in the Deep Vaults, useful but contained. In Thepyrgos, they are licensed and monitored by committees of ``mental hygiene.'' In Aelaerem, they are shunned for violating hospitality---reading thoughts without asking is the gravest sin in a culture that values the unspoken boundary. Even among the free-minded Aelinnel, they are treated with caution, for a mind that needs no Patron is a mind that answers to no law.

Their power is attrition. They pay not in Obligation but in \textit{themselves}---every thought bent, every future glimpsed, every object moved by will alone leaves a hairline crack in the vessel. They are slow to burn and slow to recover, patient as glaciers, and just as capable of crushing what stands in their path.

\subsection*{Roleplay Hooks}
\begin{itemize}
    \item You have seen your own death in a precognitive vision. It was banal---a fall from a horse, three decades hence. But you cannot stop checking your saddle straps.
    \item You once read the mind of a lover. You found exactly what you feared. You have never trusted anyone since.
    \item The Synod has a file on you. You have seen it, though you were not meant to. They do not yet know you are a psion. The game of cat and mouse sustains you.
\end{itemize}

\section{The Summoner -- The Door-Opener}
\textit{``I do not command the dead. I do not command the fey, the demons, the angels, or the things that wear the shapes of forgotten gods. I ask. I pay the price. And sometimes, when the contract is fair, they answer.''}

Summoners are the diplomats of the damned and the blessed alike, the ones who open doors and hope to close them before something follows through. But understand: \textbf{all} things beyond the veil are spirits. The dead who cling to memory are spirits. The fey who trade in stolen time are spirits. The animistic forces that sleep in stones and storms are spirits. The angels that guard the thresholds of judgment, and the demons that wait hungry at the edge of sin---they are all spirits, and they all speak the language of contract.

A Summoner of Aeler awakens the indigenous memories of stone---spirits that were never born and will never die, only sleep and dream in granite. A Summoner of Ykrul drums the ancestors from the earth---spirits of blood and bone who remember the steppe when it was ocean. A Summoner of Zakov swallows cursed coins and calls the drowned---vengeful spirits of water and salt. The church-burned heretic who calls an angel is still a Summoner. The desert madman who whispers to the storm is still a Summoner.

The difference is the \textbf{Leash}---the spiritual strain of keeping an Outsider in the world. The dead want to return to death. The fey want to return to the wild. The demons want to stay, and to make room for themselves by pushing you out. The Leash ticks. When it fills, the spirit acts once according to its nature---often destructively---then departs, or turns hostile, or simply takes what it was promised and leaves you holding the empty air.

\subsection*{Roleplay Hooks}
\begin{itemize}
    \item You once called a spirit that still follows you---not hostile, just \textit{curious}, watching from reflections and shadows, waiting for you to be interesting again.
    \item Your mentor disappeared during a summoning. You found their Symbol but not their body. You suspect they are still somewhere in the Ways Between, and someday, they will find their way back.
    \item A spirit you bound years ago now demands renegotiation. It claims the contract has a loophole---a comma out of place, a clause you mispronounced. It is right. It is always right.
\end{itemize}

\vspace{1cm}
\noindent These are the Six Roads. They cross, they diverge, and sometimes they run parallel until the walker looks down and realizes they have been walking the same path all along. Choose yours. But choose knowing that the road changes the walker, and the price is never fully paid.
\clearpage
\chapter{Patrons of the Grey Wanderer}
\epigraph{``You wish to walk the road of power? Then listen well. The world is old, and older still are the voices beneath it. We call them Patrons, though they were never sworn to us. They are the great constants, the landlords of reality, and we are merely tenants who have learned to pick the locks.''}{\textit{--- The Grey Wanderer}}

Patrons are vast intelligences---not gods in the way the Everflame preaches, though some worship them as such. They are embodiments of concepts and forces older than the first stone. Each offers power, but always with cost: fatigue, scars upon fate, or a slow unweaving of one's own story.

Understand that a Patron does not care \textit{how} you approach them. The Six Roads all lead to the same altars. A Runekeeper of Grimmir is his agent, his voice in the wild places, his priest of root and rot. An Invoker of Grimmir rents his power for a season, a gambler betting on green growth. A Witch of Grimmir tends the threshold between field and forest, owing nothing to the Wild Speaker but the breath she takes while crossing. All three draw from the same source. All three pay different prices.

This chapter presents fifteen Patrons chosen for their utility and distinct themes, organized into five categories. Each entry includes the Patron's lore, cultural variants, playing notes, suggested acts of service, a **Patron's Gift (Imbuement)**, a **Signature Talent** (a learnable ability embodying the Patron's unique power), sample Rites, GM intrusions, and---for the first time---**second voices**: fragments from those who have walked these roads before.

\medskip
\noindent\textit{"The Patron's Gift is what you get for showing up. The Signature Talent is what you earn for staying. One is a loan; the other is an investment. Both come due."}  
\hfill --- \textit{The Grey Wanderer}

\clearpage

\section{The Wild --- Patrons of Nature and Season}

\subsection{Grimmir --- The Wild Speaker}
\label{patron:grimmir}

\textit{"In the space between planted row and forest edge, where the first grain met the first acorn, Grimmir waits. He speaks in the rustle of leaves, the migration of birds, and the patient growth of ancient oaks."}

Grimmir walks the threshold between cultivated lands and untamed wilderness. Neither wholly beast nor entirely human, he speaks in root and stone, seasons and sap.

\vspace{0.2cm}
\noindent\textit{"I planted an acorn on the grave of my enemy. Now I water it with my own tears. Grimmir does not ask for forgiveness. He asks for growth."}  
\hfill --- \textit{Brynhild, Hermit of the Thornwood}

\textbf{What They Represent:} Primal wisdom, seasonal cycles, the treaty between civilisation and wilderness.

\textbf{Cultural Variants:}
\begin{itemize}
    \item \textbf{Aeler (Dwarves):} Called \textit{the Root-Father}, patron of deep fungus and stone-lichen. Shrines are carved into living rock where moss is never cleared.
    \item \textbf{Ykrul (Orcs):} A sky-wolf who herds the seasons; shamans invoke Grimmir to predict the spring thaw. His symbol is a wolf's paw print in ash.
    \item \textbf{Aelaerem (Halflings):} \textit{the Thorn-Father}, keeper of hedgerow magic. Villagers tie red thread around his roadside stones before planting.
\end{itemize}

\textbf{Playing a Follower of Grimmir:}

\textbf{The Agent (Runekeeper):} You maintain the treaty between field and forest. Your Codex is carved in slate or bark, your Thiasos a fungal thing or a beast of green shadow.

\textbf{The Gambler (Invoker):} You carry a sprig of rowan or a wolf's tooth, renting power for a season. When you Crack the Seal, the wild comes fast and hungry.

\textbf{The Threshold Keeper (Witch):} You are the hedge-mother, the Breath-Warden. Your magic is in cup-mark stones, red thread, and three knots tied counter‑timerwise.

\textbf{Suggested Acts of Service:}
\begin{itemize}
    \item Plant a tree or restore a damaged grove.
    \item Kill a predator that has grown unnaturally bold.
    \item Teach a villager to read a seasonal sign.
\end{itemize}

\textbf{Patron's Gift (Imbuement):} Your Thiasos or your hands glow with verdant energy. For one scene, gain +1 die to Survival and Nature rolls, and your melee strikes deal +1 Harm against creatures of corruption or blight.

\textbf{Signature Talent --- Wildshape (Major, 8 XP)}  
\textit{Prerequisite:} Follower of Grimmir, Body 2+.  
You may spend an action and mark 1 Fatigue to assume the form of a small beast (wolf, hawk, badger, or similar) for up to one scene. While transformed, you gain +2 dice to Stealth and Perception in natural terrain, but cannot speak or use tools. You may revert as a free action. You may use this talent once per session.

\textbf{Rites --- Strategic Use:}

\textbf{Low Rites}
\begin{itemize}
    \item \textit{The Speaking Seed} (4 XP): Commune with local plant life.
    \item \textit{The Season's Turn} (5 XP): Align with seasonal power.
\end{itemize}

\textbf{Standard Rites}
\begin{itemize}
    \item \textit{Verdant Tongue} (8 XP): Speak with all plant and animal life in Near range.
    \item \textit{The Thornveil} (7 XP): Raise a living barrier.
\end{itemize}

\textbf{Advanced Rite}
\begin{patronbox}{The World's Wound}
Cost: 13 XP, Obligation 7 segments \\
Effect: Heal a blighted place or quicken a damaged ecosystem. Clear taint, restore fertility, awaken dormant forces. Begin an \textit{Ecosystem Restoration} [8] timer.
\end{patronbox}

\textbf{GM Guidance --- Grimmir's Intrusions:}
\begin{enumerate}
    \item A tree you passed yesterday blocks your path---with a face carved into it that looks like yours.
    \item All food you touch spoils within an hour unless you offer the first bite to the earth.
    \item A pack of wolves shadows the party for a full day. They do not attack---they simply watch.
\end{enumerate}

\clearpage

\subsection{Raéyn --- Mistress of the Sea}
\label{patron:raeyn}

\textit{"Mark the tide, name your course, and trust the wave-road. But speak ill of Khemesh, and even I may let the deep take you."}

Raéyn is the tempestuous goddess of the sea, the restless tide that carries news between shores.

\vspace{0.2cm}
\noindent\textit{"She took my husband's body but gave me a compass that always points to safe harbour. I call it mercy. The admiralty calls it a bargain."}  
\hfill --- \textit{Captain Anja, Linn Tide-Witch}

\textbf{What They Represent:} Tides, change, sea travel, storms, the dual nature of the ocean.

\textbf{Cultural Variants:}
\begin{itemize}
    \item \textbf{Kahfagia:} \textit{the Lantern-Tide}. Pilots light a beacon to Raéyn before every night voyage. Her symbol is a wave curling into a spiral.
    \item \textbf{Linn:} \textit{the Whale-Mother}. Funerals at sea involve a song and a single white stone.
    \item \textbf{Zakov:} \textit{the Salt-Queen}. Pirates carve her face on their figureheads---one eye open (fair seas), one eye closed (storms ahead).
\end{itemize}

\textbf{Playing a Follower of Raéyn:} Navigator-priests (Runekeeper), storm‑singers (Cantor), selkie‑callers (Summoner).

\textbf{Suggested Acts of Service:}
\begin{itemize}
    \item Pour a libation into the sea before a voyage, asking for nothing in return.
    \item Rescue a drowning creature---a sailor, a fisher, even a seabird.
    \item Walk the shore at dawn or dusk and listen. If you hear the waves speaking a name, speak it back.
\end{itemize}

\textbf{Patron's Gift (Imbuement):} Your touch carries brine. For one scene, gain +1 die to Athletics (swimming) and Navigation, and ignore one level of environmental penalty from storms or choppy water.

\textbf{Signature Talent --- Tide-Step (Minor, 4 XP)}  
\textit{Prerequisite:} Follower of Raéyn.  
Once per session as a free action, you may treat any water-based difficult terrain as stable ground for one exchange. Alternatively, you may move across the surface of water as if it were solid for up to Near distance. Using this talent costs 1 Fatigue.

\textbf{Rites:}
\begin{itemize}
    \item \textit{Tidemark's Blessing} (Low, 4 XP): Treat slick or water-slicked footing as stable.
    \item \textit{The Changing Tide} (Standard, 8 XP): Bias currents and water levels in a zone.
    \item \textit{The Storm-Queen's Hand} (Advanced, 13 XP, Oblig. 7): Shape a storm‑band over a zone (propulsion, concealment, or smite).
\end{itemize}

\textbf{GM Guidance --- Raéyn's Intrusions:}
\begin{enumerate}
    \item A sudden fog rolls in---all navigation rolls at --2 dice for the scene.
    \item The tide turns early, stranding a vessel you need.
    \item A dolphin, seal, or sea eagle follows the party for a day. It will deliver one short message to the nearest follower of Raéyn if fed something shiny.
\end{enumerate}

\clearpage

\subsection{The Carrion-King --- Master of Endings}
\label{patron:carrionking}

\textit{"What crumbles feeds what grows. What dies becomes the soil of tomorrow's triumph. Do not mourn the fallen---harvest them."}

The Carrion-King is not a lord of rot for its own sake. He is the alchemist of endings, the patient gardener who sees every corpse as compost.

\vspace{0.2cm}
\noindent\textit{"I buried my wife in a garden of bones. In spring, the roses were the red I had never seen before. The King does not steal your grief. He teaches you to spend it."}  
\hfill --- \textit{Sister Agatha, Under‑Gardener of Aeler}

\textbf{What They Represent:} Decay, renewal, transformation.

\textbf{Cultural Variants:}
\begin{itemize}
    \item \textbf{Vhasia:} \textit{the Dust-Accountant}. Soldiers leave a button from a fallen comrade on a cairn---the King turns that button into a bullet for the next battle.
    \item \textbf{Aeler:} \textit{the Under-Gardener}. Dwarven crypts have bioluminescent fungus fed by the dead.
\end{itemize}

\textbf{Playing a Follower of the Carrion-King:} Dust‑Accountant (Runekeeper), compost‑witch (Witch), desperate gambler (Invoker).

\textbf{Suggested Acts of Service:}
\begin{itemize}
    \item Compost a fallen enemy's body and use the soil to plant something edible.
    \item Salvage a ruined structure and give the reclaimed materials to a needy family.
    \item Teach someone that loss is not failure.
\end{itemize}

\textbf{Patron's Gift (Imbuement):} Your touch smells of turned earth. For one scene, gain +1 die to Medicine and Craft rolls involving organic materials, and convert one minor item of personal value into a natural remedy or poison.

\textbf{Signature Talent --- Harvest Resonance (Minor, 4 XP)}  
\textit{Prerequisite:} Follower of the Carrion-King.  
Once per scene after an enemy is defeated or a significant ending occurs, you may draw a sliver of that ending's energy. Gain +1 die on your next action, or clear 1 Fatigue. You cannot use this talent again until the next scene.

\textbf{Rites:}
\begin{itemize}
    \item \textit{Rite of Consuming Rot} (Low, 5 XP): Accelerate natural decay to weaken organic materials.
    \item \textit{Rite of the Fertile Death} (Standard, 8 XP): Transform death into growth---create advantageous terrain or grant healing bonuses.
    \item \textit{The Eternal Cycle} (Advanced, 14 XP, Oblig. 8): Destroy one major asset, enemy, or obstacle, and create something new of equal or greater value.
\end{itemize}

\textbf{GM Guidance --- Carrion-King's Intrusions:}
\begin{enumerate}
    \item A dead creature you passed earlier now follows you---not as an enemy, but as a witness.
    \item All food in the party's packs spoils instantly, except for a single loaf of black bread.
    \item A villager you saved bears a small rot‑mark that spreads one inch each day.
\end{enumerate}

\clearpage

\section{The Shadow --- Patrons of Secrets and Transformation}

\subsection{Ikasha --- She Who Sleeps}
\label{patron:ikasha}

\textit{"Blow out the candle. If the room listens back, ask softly. At the next crossroads, the raven waits---and the shadow remembers your passing."}

Ikasha is the hush between footfalls, the patience of dark water, the black‑feathered watcher at every threshold.

\vspace{0.2cm}
\noindent\textit{"I whispered a name into the dark. Three days later, a tyrant forgot his own face. The raven did not ask for payment. It only asked that I never speak the name again."}  
\hfill --- \textit{The Unnamed, Tulkani exile}

\textbf{What They Represent:} Latent potential, shadow, thresholds, forgotten things.

\textbf{Cultural Variants:}
\begin{itemize}
    \item \textbf{Mistlands:} \textit{the Bell-Shrouded}---bell‑wardens leave a single rope unknotted for her.
    \item \textbf{Zakov:} \textit{the Tide Watcher}, a smuggler's patron.
    \item \textbf{Aeler:} \textit{the Deep Listener}; miners whisper to her before descending into new shafts.
\end{itemize}

\textbf{Playing a Follower of Ikasha:} Bell‑warden (Witch), secret notary (Runekeeper), shadow‑psion (Psion).

\textbf{Suggested Acts of Service:}
\begin{itemize}
    \item Spend a night in absolute darkness, speaking no word and lighting no flame.
    \item Carry a secret for someone---a truth that would ruin them---and tell no one for a full season.
    \item At a crossroads, leave a black feather and a single coin. Do not look back.
\end{itemize}

\textbf{Patron's Gift (Imbuement):} Your form drinks light. For one scene, gain +1 die to Stealth and Deception, and your footsteps make no sound unless you choose otherwise.

\textbf{Signature Talent --- Shadow Meld (Major, 6 XP)}  
\textit{Prerequisite:} Follower of Ikasha.  
Once per session as an action, step into any shadow larger than yourself and emerge from another shadow within Near range. You cannot be targeted while moving between shadows. You must mark 1 Fatigue to use this talent.

\textbf{Alternative Talent --- Forment the Unspoken (Major, 8 XP)}  
\textit{Prerequisite:} Follower of Ikasha, Presence 3+.  
You have learned to seed rebellion in silence. Over the course of a downtime scene, you may whisper a single doubt into the ear of a crowd, a guard, or a court. The GM starts a \textit{Discontent} [6] timer. When it fills, the target community or faction suffers a public crisis (riots, desertions, a vote of no confidence). You cannot use this talent again until the timer completes or is thwarted.

\textbf{Rites:}
\begin{itemize}
    \item \textit{Touch the Umbral Veil} (Low, 4 XP): Start Controlled on one Stealth roll.
    \item \textit{Draw from the Umbral Reservoir} (Standard, 8 XP): Grant +2 dice to stealth, deception, or resolve.
    \item \textit{Become the Shadow at the Crossroads} (Advanced, 12 XP, Oblig. 7): Intangible to mundane harm, pass through thresholds.
\end{itemize}

\textbf{GM Guidance --- Ikasha's Intrusions:}
\begin{enumerate}
    \item A raven lands on your shoulder and whispers a secret you never wanted to know.
    \item Your shadow stretches into the shape of a woman with a crown of black feathers.
    \item A door you passed through earlier is now locked from the other side---and the key is in your pocket.
\end{enumerate}

\clearpage

\subsection{Isoka --- Angel of Serpents}
\label{patron:isoka}

\textit{"Do not mourn the skin you shed. It was never meant to last. The venom that burns away the old self is the same that grants the strength to become new."}

Isoka teaches that every self is temporary---identity is a skin to be shed so growth can continue.

\vspace{0.2cm}
\noindent\textit{"I cut away my name, my debts, my mother's disappointed face. The skin healed. The freedom did not. Isoka only promises change. She never promises direction."}  
\hfill --- \textit{Copperhead, Ashaani infiltrator}

\textbf{What They Represent:} Transformation, renewal, venom, the casting off of old selves.

\textbf{Cultural Variants:}
\begin{itemize}
    \item \textbf{Zakov:} \textit{the Tide-Shedder}. Pirates leave a lock of hair and a name at her sunken temple; eels eat both.
    \item \textbf{Acasia:} \textit{the Copperhead}. Hedge‑cultists shed identities to infiltrate noble houses.
\end{itemize}

\textbf{Playing a Follower of Isoka:} Shedded (Cantor), Copperhead (Runekeeper), shed‑vigil witch (Witch).

\textbf{Suggested Acts of Service:}
\begin{itemize}
    \item Shed a physical token of your old life---a piece of clothing, a weapon, a letter---and leave it at a crossroads.
    \item Forgive someone who wronged you, publicly and completely, as an act of transformation.
    \item Help someone else change---cut their hair, give them new clothes, teach them a new name.
\end{itemize}

\textbf{Patron's Gift (Imbuement):} Your movements become fluid and sinuous. For one scene, gain +1 die to Athletics (climbing, squeezing) and Deception, and re‑roll one failed escape attempt.

\textbf{Signature Talent --- Shed Skin (Minor, 4 XP)}  
\textit{Prerequisite:} Follower of Isoka.  
Once per session when you are Bound, Grappled, or suffering from a Condition that restricts movement, you may spend an action and mark 1 Fatigue to shed that restriction. Describe how you slip free like a snake shedding its skin.

\textbf{Alternative Talent --- Momentum of the Molt (Standard, 6 XP)}  
\textit{Prerequisite:} Follower of Isoka.  
When you successfully resist or break free from a Condition, you gain a burst of furious energy. Your next action within the same scene gains +1 die and treats its Position as one step higher (Desperate → Dominant, Dominant → Controlled). You may use this talent once per session.

\textbf{Rites:}
\begin{itemize}
    \item \textit{Venomous Benediction} (Low, 5 XP): Bless an ally's strike with +1 Harm and possible Poison.
    \item \textit{The Shedding} (Standard, 8 XP): +2 dice to resist a named Condition; declare a minor contingency retroactively.
    \item \textit{Complete Metamorphosis} (Advanced, 13 XP, Oblig. 7): Change appearance, voice, and mannerisms completely. Warded against recognition.
\end{itemize}

\textbf{GM Guidance --- Isoka's Intrusions:}
\begin{enumerate}
    \item A piece of your old self returns---a letter from a dead lover, a scar you had forgotten.
    \item Your shadow moves independently, pointing at people who are about to change---to leave, to betray, to die.
    \item You begin to shed skin overnight. If you burn the cast, you forget one minor memory. If you keep it, someone will steal it.
\end{enumerate}

\clearpage

\subsection{Aveh --- The Rider Behind the Storm}
\label{patron:aveh}

\textit{"I am the track the storm devours. I am the freedom you fear and the silence you crave. To ride with me is to be unbound---and to be unbound is to vanish."}

Aveh is the faceless rider at the horizon's edge, offering liberty at the cost of belonging.

\vspace{0.2cm}
\noindent\textit{"She appeared as a woman with no face and a voice like dry grass. "Run," she said. I ran. I never stopped. I still don"t know if I should thank her or curse her."}  
\hfill --- \textit{The Dust‑Walker, escaped slave}

\textbf{What They Represent:} Freedom, erasure, storms, exile.

\textbf{Cultural Variants:}
\begin{itemize}
    \item \textbf{Ykrul:} A sky‑spirit that rides with winter gales; shamans invoke Aveh to hide a dying person's tracks from the dead.
    \item \textbf{Linn:} A sea‑fog rider; sailors whisper Aveh's name to slip through blockades.
\end{itemize}

\textbf{Playing a Follower of Aveh:} Dust‑walker (Invoker), path‑finder (Runekeeper), storm‑rider summoner (Summoner).

\textbf{Suggested Acts of Service:}
\begin{itemize}
    \item Cut a binding contract and scatter the pieces to the wind.
    \item Help someone flee oppression---a slave, a debtor, a threatened child.
    \item Spend a night sleeping under open sky with no shelter.
\end{itemize}

\textbf{Patron's Gift (Imbuement):} Your presence becomes hard to hold in memory. For one scene, gain +1 die to Stealth and Escape, and anyone who sees you must succeed on a Resolve test (DV 2) to remember your face clearly.

\textbf{Signature Talent --- Storm-Step (Major, 6 XP)}  
\textit{Prerequisite:} Follower of Aveh.  
Once per session as a move action, vanish in a swirl of wind and dust and reappear anywhere within Near range that you can see. This movement does not provoke opportunity attacks. Mark 1 Fatigue after using this talent.

\textbf{Rites:}
\begin{itemize}
    \item \textit{Erase the Road} (Standard, 8 XP): Trails vanish, records blur, pursuers lose the way.
    \item \textit{Storm's Refuge} (Advanced, 12 XP, Oblig. 6): Shroud your company---unseen or forgotten by pursuers. All present mark 1 Corruption or forget one bond at the end.
\end{itemize}

\textbf{GM Guidance --- Aveh's Intrusions:}
\begin{enumerate}
    \item A storm erases your tracks---and also the memory of your face from the last town.
    \item Your shadow detaches and wanders off for an hour.
    \item A forgotten obligation suddenly resurfaces---someone you wronged now knows exactly where you are.
\end{enumerate}

\clearpage

\section{The Covenant --- Patrons of Law and Truth}

\subsection{The Inquisitor Prime --- Purity and Domination}
\label{patron:inquisitor}

\textit{"Mercy is weakness. Purity is survival. Where corruption hides, we carve it out."}

Among zealots and witch‑hunters, the Inquisitor Prime is venerated as the hand of absolute purity and the sword of uncompromising order.

\vspace{0.2cm}
\noindent\textit{"I burned my sister's grimoire with my own hands. The Inquisitor smiled through my tears. I have never known if that smile was approval or mockery."}  
\hfill --- \textit{Vicar Albrecht, Censer‑Knight}

\textbf{What They Represent:} Purity, order, judgment, eradication of the tainted.

\textbf{Cultural Variants:}
\begin{itemize}
    \item \textbf{Ecktoria (Everflame):} \textit{the Unquenchable Fire}. Priests carry censers of perpetual coal.
    \item \textbf{Aeler:} \textit{the Light of Forges}. Dwarven smiths kindle hearths with a brand from the temple's eternal flame.
\end{itemize}

\textbf{Playing a Follower of the Inquisitor Prime:} Censer‑Knight (Runekeeper), desperate purifier (Invoker), mind‑hunter (Psion).

\textbf{Suggested Acts of Service:}
\begin{itemize}
    \item Burn a heretical text or destroy a cursed object.
    \item Expose a charlatan who uses magic to deceive the innocent.
    \item Hunt down a rogue summoner or corrupt Runekeeper.
\end{itemize}

\textbf{Patron's Gift (Imbuement):} Your gaze burns with righteous fire. For one scene, gain +1 die to Insight and Intimidation when confronting supernatural creatures, and your melee attacks ignore 1 point of Armor from undead or demons.

\textbf{Signature Talent --- Truth's Gaze (Standard, 6 XP)}  
\textit{Prerequisite:} Follower of the Inquisitor Prime.  
Once per session as an action, fix your gaze on a target within Near range. For the next minute, they cannot tell a direct lie to you; they may evade or refuse to answer, but any statement they make must be factually true. Mark 1 Fatigue after using this talent.

\textbf{Rites:}
\begin{itemize}
    \item \textit{Rite of the Pure Flame} (Low, 4 XP): +2 dice to resist supernatural influence.
    \item \textit{Rite of the Cleansing Light} (Standard, 8 XP): Dispel a magical effect or weaken it.
    \item \textit{The Final Admonition} (Advanced, 14 XP, Oblig. 8): Attempt to annihilate one supernatural foe.
\end{itemize}

\textbf{GM Guidance --- Inquisitor Prime's Intrusions:}
\begin{enumerate}
    \item An old investigation resurfaces---someone you once cleared is now accused again.
    \item Your zeal attracts attention from both enemies and \textit{allies} who fear your methods.
    \item A supernatural entity you wounded escapes and seeks revenge.
\end{enumerate}

\clearpage

\subsection{The Witness --- Keeper of Uncomfortable Truths}
\label{patron:witness}

\textit{"I will show you what you would rather forget. But first, you must forget what you think you know."}

The Witness remembers what others bury. Every shadow cast and oath broken is a line in her unending record. She does not judge; she merely \textit{records}.

\vspace{0.2cm}
\noindent\textit{"I looked into her mirror and saw the face of the man I had helped ruin, twenty years ago. The Witness does not forgive. She only confirms."}  
\hfill --- \textit{Archivist Venn, Thepyrgos}

\textbf{What They Represent:} Truth, revelation, memory, the weight of knowledge.

\textbf{Cultural Variants:}
\begin{itemize}
    \item \textbf{Ecktoria:} \textit{the Ink-Eye}. Censor's clerks copy every redacted line into a private scroll.
    \item \textbf{Vhasia:} \textit{the Fractured Mirror}. Spies pray to her to reveal the true loyalty of rivals.
\end{itemize}

\textbf{Playing a Follower of the Witness:} Notary of truth (Runekeeper), memory‑psion (Psion), echo‑witch (Witch).

\textbf{Suggested Acts of Service:}
\begin{itemize}
    \item Expose a lie told by someone in power.
    \item Write down a secret you have never told anyone, then burn the paper.
    \item Testify truthfully in a legal proceeding, even if it hurts your cause.
\end{itemize}

\textbf{Patron's Gift (Imbuement):} Your eyes glimmer with remembered light. For one scene, gain +1 die to Investigation and Insight, and ask one factual question about a scene or object; the GM must answer truthfully, though perhaps cryptically.

\textbf{Signature Talent --- Recorded Echo (Minor, 4 XP)}  
\textit{Prerequisite:} Follower of the Witness.  
Once per session, touch a physical trace (a footprint, a torn cloth, a drop of blood) and instantly know the last significant action taken by the person who left it (e.g., "they were running," "they hid something under the floor"). Mark 1 Fatigue after using this talent.

\textbf{Rites:}
\begin{itemize}
    \item \textit{Echoing Truth} (Standard, 8 XP): Compel a truthful answer to one direct question.
    \item \textit{Immutable Record} (Standard, 7 XP): Bind an agreement. Breakers suffer reputation damage.
    \item \textit{The Unveiled Heart} (Advanced, 12 XP, Oblig. 6): Target suffers --2 dice to conceal true emotions or lies.
\end{itemize}

\textbf{GM Guidance --- The Witness's Intrusions:}
\begin{enumerate}
    \item You remember something you never witnessed---a crime, a conversation, a death---in perfect detail.
    \item A mirror shows not your face, but the face of someone you have wronged.
    \item A secret you have kept is now known to a stranger---a blackmailer, a journalist, a rival.
\end{enumerate}

\clearpage

\subsection{The Traveler --- Lord of the Open Road}
\label{patron:traveler}

\textit{"One foot in a promise, and the road will meet you halfway. But break your word to the way, and the way will break you."}

The Traveler is not merely one who shows the path---they \textit{are} the path, existing in the pause between steps and in the choice of which fork to take when roads diverge.

\vspace{0.2cm}
\noindent\textit{"I asked the Traveler for a shortcut. He gave me a path that shaved three days off my journey but added two new enemies. A fair trade, I suppose."}  
\hfill --- \textit{Merchant Zara, Fhara caravaneer}

\textbf{What They Represent:} Ways, journeys, thresholds, safe passage.

\textbf{Cultural Variants:}
\begin{itemize}
    \item \textbf{Kahfagia:} \textit{the Lantern-Bearer}. Pilots light a beacon at the start of every convoy.
    \item \textbf{Ykrul:} \textit{the Hoof-Print}. Riders leave a single arrowhead at a crossroads for luck.
\end{itemize}

\textbf{Playing a Follower of the Traveler:} Guide‑priest (Runekeeper), bridge‑maker (Witch), road‑spirit summoner (Summoner).

\textbf{Suggested Acts of Service:}
\begin{itemize}
    \item Walk a mile of road barefoot, carrying a stone in each hand.
    \item Offer shelter to a lost traveller for a night---no questions, no payment.
    \item Repair a broken signpost, milestone, or footbridge.
\end{itemize}

\textbf{Patron's Gift (Imbuement):} Your feet never tire. For one scene, gain +1 die to Athletics (running, climbing) and Navigation, and ignore one level of difficult terrain.

\textbf{Signature Talent --- Shortcut (Minor, 4 XP)}  
\textit{Prerequisite:} Follower of the Traveler.  
Once per session when moving from one location to another (including during travel montages), declare a hidden shortcut that reduces the remaining travel timer by 2 segments. The GM may describe a minor cost or complication (e.g., "the path is steep---mark 1 Fatigue").

\textbf{Rites:}
\begin{itemize}
    \item \textit{Road-Sense} (Low, 4 XP): Unerringly pick the fastest safe route.
    \item \textit{The Waymark} (Standard, 8 XP): Declare a lane as permitted passage for your party.
    \item \textit{The Wayfarer's Covenant} (Advanced, 14 XP, Oblig. 8): Swear a covenant of safe passage. +2 Effect on travel rolls, re‑roll one failed navigation test per scene.
\end{itemize}

\textbf{GM Guidance --- The Traveler's Intrusions:}
\begin{enumerate}
    \item A path you have walked before now leads somewhere else---a different building, a different forest.
    \item Your boots wear out instantly. You must walk barefoot for a scene.
    \item A stranger approaches on the road and asks for directions. If you lie, you double your current Obligation.
\end{enumerate}

\clearpage

\section{The Infinite --- Patrons of the Unseen and Unknowable}

\subsection{Khemesh --- The Abyssal Maw}
\label{patron:khemesh}

\textit{"In the trench without light, the Maw waits. Even silence drowns."}

Khemesh is not merely a lord of the depths---he is the hunger beneath them, a pressure older than seas.

\vspace{0.2cm}
\noindent\textit{"I threw my wedding ring into the abyss. Khemesh gave me back a coin that buys silence wherever I go. I have not spoken my wife's name in three years. I do not miss it."}  
\hfill --- \textit{Drowned Prince, Zakov pirate}

\textbf{What They Represent:} Depth, pressure, inevitability, eldritch terror.

\textbf{Cultural Variants:}
\begin{itemize}
    \item \textbf{Zakov:} \textit{the Drowned Prince}. Pirates toss a silver coin overboard before a raid---if it sinks, Khemesh approves.
    \item \textbf{Linn:} \textit{the Reef-Father}. Fishermen leave the first catch of the day to him.
\end{itemize}

\textbf{Playing a Follower of Khemesh:} Pressure‑keeper (Runekeeper), deep‑gambler (Invoker), trench‑caller (Summoner).

\textbf{Suggested Acts of Service:}
\begin{itemize}
    \item Throw a valuable object into the deepest water you can reach.
    \item Spend a night submerged up to your neck in cold water.
    \item Reveal a secret you have kept for more than a year where water can hear.
\end{itemize}

\textbf{Patron's Gift (Imbuement):} Your presence feels like the deep ocean---heavy, unyielding. For one scene, gain +1 die to Intimidation and Endurance, and hostile creatures must test Resolve to move closer to you.

\textbf{Signature Talent --- Crushing Silence (Standard, 6 XP)}  
\textit{Prerequisite:} Follower of Khemesh.  
Once per session as a reaction when a creature within Near range speaks or casts a spell, impose a bubble of oppressive silence around them for one exchange. They cannot speak, and any spell with a verbal component fails. Mark 1 Fatigue after using this talent.

\textbf{Rites:}
\begin{itemize}
    \item \textit{Pressure of the Maw} (Standard, 7 XP): Pin a target with invisible force---treat as [ENTANGLE].
    \item \textit{Rite of the Abyssal Vision} (Standard, 9 XP): See the world as Khemesh does---fractured, alien.
    \item \textit{The Maw Opens} (Advanced, 12 XP, Oblig. 8): Reality folds inward. Enemies act at Desperate Position by default.
\end{itemize}

\textbf{GM Guidance --- Khemesh's Intrusions:}
\begin{enumerate}
    \item Water pressure quadruples in the room---wooden objects groan, glassware cracks.
    \item A nightmare tentacle manifests from a nearby body of water and steals a shiny object.
    \item The party begins to \textit{sink}---not physically, but emotionally. Each character marks 1 Fatigue or confesses a secret fear aloud.
\end{enumerate}

\clearpage

\subsection{The Ninth --- Beyond Comprehension}
\label{patron:ninth}

\textit{"Eight figures bind the world; the Ninth overflows it. Drink, and remember that the cup forgets its shape before the song is done."}

Beyond the eighth figure of Sacred Geometry lies a rim no compass will hold. The Ninth is not absence but surplus: knowledge poured past the lip of any vessel, truth that overgrows the frame that seeks to bind it.

\vspace{0.2cm}
\noindent\textit{"I calculated the Ninth digit of pi. For three days, I saw the angles between thoughts. Then I forgot how to speak. The Ninth gives and takes in the same breath."}  
\hfill --- \textit{Geometer Tolin, Aeler}

\textbf{What They Represent:} Hyperdimensional knowledge, incomprehensible truths, information overflow.

\textbf{Cultural Variants:}
\begin{itemize}
    \item \textbf{Thepyrgos:} \textit{the Lost Syllable}. Archivists believe the Ninth is a word erased from every dictionary---but still whispered in dreams.
    \item \textbf{Aeler:} \textit{the Deep Count}. Geometers speak of a number that cannot be written, only approximated.
\end{itemize}

\textbf{Playing a Follower of the Ninth:} Hyper‑psion (Psion), cartographer of the impossible (Runekeeper), gap‑witch (Witch).

\textbf{Suggested Acts of Service:}
\begin{itemize}
    \item Solve a puzzle or riddle that has stumped others for years.
    \item Destroy a book that contains a comfortable lie, replacing it with a painful truth.
    \item Teach someone a concept they are not ready to understand---then watch what happens.
\end{itemize}

\textbf{Patron's Gift (Imbuement):} Your thoughts race ahead of your tongue. For one scene, gain +1 die to Lore and Investigation, and ask one question about a pattern or hidden connection; the GM must answer with a cryptic truth.

\textbf{Signature Talent --- Glimpse Beyond (Minor, 4 XP)}  
\textit{Prerequisite:} Follower of the Ninth.  
Once per session, look at a complex situation (a crime scene, a puzzle lock, a battle formation) and instantly intuit the single most important hidden detail. The GM gives you one piece of information that would otherwise require a successful Investigation roll. Mark 1 Fatigue after using this talent.

\textbf{Rites:}
\begin{itemize}
    \item \textit{Ninth Figure} (Low, 4 XP): +1 die to Lore or Investigation when wrestling complex systems.
    \item \textit{Hypercognitive Lattice} (Standard, 8 XP): Treat Wits as +2 higher for information‑processing rolls.
    \item \textit{The Ninth Revelation} (Advanced, 14 XP, Oblig. 8): Glimpse a shard of Ninth truth---perfect command of a knotty system or comprehension of an "impossible" relation.
\end{itemize}

\textbf{GM Guidance --- The Ninth's Intrusions:}
\begin{enumerate}
    \item You see patterns in everything---the way dust falls, the arrangement of stones. It is useful, but exhausting. Mark 1 Fatigue.
    \item A fact you learned from the Ninth is true, but \textit{should not be}. Others find it disturbing when you speak it aloud.
    \item You begin to forget mundane things (names, faces, what you ate) while remembering impossible geometries with perfect clarity.
\end{enumerate}

\clearpage

\section{The Necessary Evils --- Patrons of Bond, Law, and Ruin}

Between the wild places and the shadowed thresholds, between the covenant of truth and the infinite abyss, there exist Patrons who answer not to hope or fear, but to \textit{necessity}. They are the ones mortals call when a vow must be witnessed, when a burden must be shouldered, or when a web must be woven. They are not kind. They are not cruel. They are \textit{inevitable}.

\subsection{Mykkiel --- Arbiter of the Covenant}
\label{patron:mykkiel}

\textit{"The Law is not written in sand, but in stone. The Covenant is not suggestion, but command. Transgressors shall be purged, the faithful exalted."}

Mykkiel is the unwavering judge of covenants, the divine lawgiver whose scales and sword uphold sacred order.

\vspace{0.2cm}
\noindent\textit{"I judged a man for stealing bread to feed his children. The law said: hang. Mykkiel's scale did not tip. That night, I dreamed of an empty courtroom and a gavel that struck only silence."}  
\hfill --- \textit{Justice Finley, Vhasian Parlement}

\textbf{What They Represent:} Divine law, covenant, contract, divine judgment, written word made flesh.

\textbf{Cultural Variants:}
\begin{itemize}
    \item \textbf{Aeler:} \textit{the Stone-Judge}. Dwarven arbiters inscribe his name on their gavels.
    \item \textbf{Vhasia:} \textit{the Sun-Covenant}. His symbol is a broken chain, signifying that agreements are binding even when polities are not.
\end{itemize}

\textbf{Playing a Follower of Mykkiel:} Notary with teeth (Runekeeper), hedge‑lawyer (Witch), contract‑gambler (Invoker).

\textbf{Suggested Acts of Service:}
\begin{itemize}
    \item Witness a formal agreement between two parties who have no other witness.
    \item Judge a dispute fairly, even if the fair outcome harms your interests.
    \item Repair or renew a broken covenant---a marriage, a trade deal, a treaty.
\end{itemize}

\textbf{Patron's Gift (Imbuement):} Your words carry the weight of testimony. For one scene, gain +1 die to Insight and Command when judging or mediating, and any promise you witness is subtly harder to break (increase DV to break oath by 1).

\textbf{Signature Talent --- Seal of Witness (Minor, 4 XP)}  
\textit{Prerequisite:} Follower of Mykkiel.  
Once per session, place your mark on a written or spoken agreement. Any party who knowingly violates the agreement suffers a --1 die to all actions until they make amends. Mark 1 Fatigue after using this talent.

\textbf{Rites:}
\begin{itemize}
    \item \textit{Rite of Hallowed Ground} (Low, 5 XP): Create sacred space where divine law prevails.
    \item \textit{The Covenant Seal} (Standard, 9 XP): Formalise a binding covenant with concrete penalties for breach.
    \item \textit{Divine Judgment} (Advanced, 13 XP, Oblig. 7): Pronounce divine judgment on a covenant‑breaker.
\end{itemize}

\textbf{GM Guidance --- Mykkiel's Intrusions:}
\begin{enumerate}
    \item A covenant you thought fulfilled is re‑examined; they find a technical violation.
    \item Your signature appears on a document you have never seen---a contract, a treaty, a marriage licence.
    \item The scales of justice appear before you, weighing a good deed and a bad one. If the bad outweighs the good, you lose 1 Boon.
\end{enumerate}

\clearpage

\subsection{Malachai --- The Cruel Messenger}
\label{patron:malachai}

\textit{"I offer you everything you desire, and everything you fear. The payment comes later---and it always comes."}

Once a divine messenger tasked with delivering painful truths, Malachai became obsessed with the beauty of destructive revelations.

\vspace{0.2cm}
\noindent\textit{"She gave me luck for a year. I won every bet, every duel, every heart. Then she came back and asked for my left hand. I cut it off myself. She smiled and said, "Now you are wise.""}  
\hfill --- \textit{One‑Handed Jaks, Zakov gambler}

\textbf{What They Represent:} Cursed gifts, corruptive power, false salvation, lycanthropy and other afflictions.

\textbf{Cultural Variants:}
\begin{itemize}
    \item \textbf{Linn:} \textit{the Salt-Fang}. Shipwreck survivors sometimes find a silver coin in their mouth.
    \item \textbf{Valewood:} \textit{the Thorn-Giver}. Hunters who curse a rival may find their arrows always strike true---but the blood never dries.
\end{itemize}

\textbf{Playing a Follower of Malachai:} Desperate debtor (Invoker), honey‑voiced (Cantor), curse‑binder (Summoner).

\textbf{Suggested Acts of Service:}
\begin{itemize}
    \item Break a curse you have inflicted on someone else---even an enemy.
    \item Offer a false hope to someone who deserves it, then reveal the truth.
    \item Accept a minor burden as penance for a major one you have already incurred.
\end{itemize}

\textbf{Patron's Gift (Imbuement):} Your touch leaves a faint, sweet scent---like honey over rot. For one scene, gain +1 die to Deception and Intimidation, and any target you successfully deceive marks 1 Fatigue.

\textbf{Signature Talent --- Borrowed Scars (Standard, 6 XP)}  
\textit{Prerequisite:} Follower of Malachai.  
Once per session, transfer a Condition or 1 Fatigue from an ally to yourself. In exchange, gain +1 die on your next roll against the source of that affliction. Mark 1 Obligation (to Malachai) after using this talent.

\textbf{Rites:}
\begin{itemize}
    \item \textit{Honeyed Curse} (Low, 4 XP): Target gains +2 dice to their next roll but marks 1 Corruption.
    \item \textit{False Dawn} (Standard, 8 XP): Bestow a gift with a hidden curse---target gains an advantage but begins a Corruption timer.
    \item \textit{The Messenger's Burden} (Advanced, 13 XP, Oblig. 7): Channel Malachai directly. Gain +2 dice to one action per beat, but each use marks 1 SB.
\end{itemize}

\textbf{GM Guidance --- Malachai's Intrusions:}
\begin{enumerate}
    \item A gift you gave earlier curdles. The recipient now suffers a visible curse---a limp, a stutter.
    \item The moon turns red for one night. All lycanthropes test Resolve (DV 5) or lose control.
    \item A messenger arrives with a letter sealed in black wax. It is from someone you cursed years ago. They are coming.
\end{enumerate}

\clearpage

\subsection{Inaea --- Angel of the Spider}
\label{patron:inae}

\textit{"The thread that comforts is the same that binds. The knot that saves is also the snare."}

Inaea sits in patient silence at the centre of the unseen lattice. Where others see isolation, she sees threads---affection, rivalry, loyalty, fear.

\vspace{0.2cm}
\noindent\textit{"I tied a red thread around my daughter's wrist to keep her safe. She married a monster. The thread held; her heart did not. Inaea never promised happiness. Only connection."}  
\hfill --- \textit{Elder Thistle, Aelaerem hedgewitch}

\textbf{What They Represent:} Webs, fate, connection, mercy, entanglement.

\textbf{Cultural Variants:}
\begin{itemize}
    \item \textbf{Aelaerem (Halflings):} \textit{the Kettle-Spinner}. Grandmothers invoke her to keep families together.
    \item \textbf{Silkstrand:} \textit{the Dyed Web}. Dyers" guilds leave a single strand of silk outside their windows on the first of spring.
\end{itemize}

\textbf{Playing a Follower of Inaea:} Kettle‑Spinner (Witch), weaver‑priest (Runekeeper), web‑psion (Psion).

\textbf{Suggested Acts of Service:}
\begin{itemize}
    \item Mend a broken relationship---two friends who quarrelled, a parent and child estranged.
    \item Weave a small textile (a scarf, a bandage) and give it to someone in need.
    \item Untie a knot that has been tied for too long---a prisoner's chains, a stuck door.
\end{itemize}

\textbf{Patron's Gift (Imbuement):} Your fingers move with invisible thread. For one scene, gain +1 die to Sway and Insight, and reroll one failed social roll (once per Gift).

\textbf{Signature Talent --- Tug the Strand (Minor, 4 XP)}  
\textit{Prerequisite:} Follower of Inaea.  
Once per session, name a relationship between two NPCs (friend, rival, debtor, lover). The GM must reveal one secret about that relationship---a hidden grievance, a secret loyalty, an unpaid obligation. Mark 1 Fatigue after using this talent.

\textbf{Rites:}
\begin{itemize}
    \item \textit{Hearth-Thread Knot} (Low, 4 XP): You and an ally gain +1 die to Aid each other.
    \item \textit{Strand of Inevitability} (Standard, 8 XP): Link two actors---when one acts, the other is dragged into consequence.
    \item \textit{Binding Knot} (Advanced, 11 XP, Oblig. 6): Bind two parties to a vow. Breaking it forces 2 SB and marks the violator with a visible tell.
\end{itemize}

\textbf{GM Guidance --- Inaea's Intrusions:}
\begin{enumerate}
    \item A spider lowers itself onto your shoulder. In its web is a single perfect word---the name of someone who needs your help.
    \item Every handshake you offer for the next day feels like grasping a thread---you sense whether the person is lying, but they also sense your scrutiny.
    \item A knot you tied earlier is now impossible to undo without cutting. Cutting it severs a minor bond---you lose one Boon.
\end{enumerate}

\clearpage

\subsection{Varnek Karn --- The Death's Negotiator}
\label{patron:varnek}

\textit{"The dead have all the time in the world to consider their offers, and no need to rush to agreement. They know what the living forgot, and remember every slight with perfect clarity. In Varnek's halls, both parties pay a price---but the currency is never what you expect."}

Varnek Karn dwells in the grey between life and death, where the newly departed linger with unfinished business and elder shades guard secrets dearer than gold. He is the broker who speaks across the veil, the mediator of disputes between living and dead, and the archivist of every promise sworn beyond the grave.

\vspace{0.2cm}
\noindent\textit{"I asked my grandmother's ghost where she hid the will. She demanded a lock of my firstborn's hair. I have no children. We settled on a memory of her favorite cake. Varnek Karn recorded the transaction in a ledger that smelled of wet stone."}  
\hfill --- \textit{Notary Venn, Vhasian Inheritance Court}

\textbf{What They Represent:} Post‑mortem diplomacy, ancestral wisdom, unfinished business, death‑bound contracts.

\textbf{Cultural Variants:}
\begin{itemize}
    \item \textbf{Aeler:} \textit{the Breath-Accountant}. Dwarven mortuary priests audit the spiritual ledgers of the Deep Vaults.
    \item \textbf{Ykrul:} \textit{the Bone-Speaker}. Ancestral spirits are consulted before any major migration.
    \item \textbf{Vhasia:} \textit{the Silent Notary}. Parlement clerks whisper that Varnek Karn holds the original copies of every treaty signed in blood.
    \item \textbf{Mistlands:} \textit{the Bell-Shrouded}. Wardens of the levee bells leave a single rope unknotted when a death is disputed.
\end{itemize}

\textbf{Playing a Follower of Varnek Karn:} Notary of the dead (Runekeeper), bone‑gambler (Invoker), ancestral caller (Summoner), death‑threshold witch (Witch).

\textbf{Suggested Acts of Service:}
\begin{itemize}
    \item Perform a proper burial for a forgotten corpse, speaking its name aloud.
    \item Carry a message from the living to a ghost, and a reply back.
    \item Settle a death‑bound obligation---pay what the deceased owed, or collect what they were owed.
    \item Witness a final testament and ensure its terms are honored.
\end{itemize}

\textbf{Patron's Gift (Imbuement):} Your presence carries the scent of old parchment and funeral incense. For one scene, gain +1 die to Insight and Lore when dealing with spirits, death, or inheritance, and ask one question of a corpse less than a day dead; it answers with a single word or a gesture.

\textbf{Signature Talent --- The Silent Ledger (Standard, 6 XP)}  
\textit{Prerequisite:} Follower of Varnek Karn.  
Once per session, touch a physical trace of a death (a body, a grave, an object that witnessed a passing) and instantly know one \textit{unresolved obligation} connected to that death---a promise unfulfilled, a secret unspoken, an unpaid weight. Mark 1 Fatigue after using this talent.

\textbf{Greater Talent --- Reoccupied Vessel (Major, 10 XP)}  
\textit{Prerequisite:} Follower of Varnek Karn, Summoner's Path or Runekeeper's Path.  
You may perform a 10‑minute ritual at the side of a corpse no more than three days dead. Negotiate with the spirit (which must be willing or sufficiently indebted). If it agrees, it reoccupies its body, becoming a \textbf{Bound Vessel} (Cap 3, Tier II--III) that serves for a set term or until a specific task is completed. Use the Summoning Leash mechanics: the Leash ticks when the spirit acts against its nature, when the corpse takes damage, or when it is forced to remain beyond the agreed term. When the Leash fills, the spirit departs and the corpse decays to dust. Using this talent marks +2 Obligation (to Varnek Karn). You may Push It: the spirit serves for one week instead of one scene, but marks +1 Obligation and causes visible rot (Crew Morale +1 if witnessed).

\textbf{Rites --- Strategic Use:}

\textbf{Low Rites}
\begin{itemize}
    \item \textit{Rite of the Crossing Fee} (4 XP): Open speech with a specific dead individual. +1 die to Investigation or Lore when seeking knowledge they once held.
    \item \textit{Rite of the Ancestral Consultation} (5 XP): Ask one question of an ancestral spirit regarding family history, hidden legacies, or inherited burdens.
\end{itemize}

\textbf{Standard Rites}
\begin{itemize}
    \item \textit{Rite of the Death-Bound Pact} (8 XP): Bind a compact between a living signatory and a willing spirit. Both gain +1 die when fulfilling the terms.
    \item \textit{Rite of the Spirit Medium} (7 XP): Serve as conduit for several spirits to speak. Gain +2 dice to mediate disputes between living and dead.
\end{itemize}

\textbf{Advanced Rites}
\begin{patronbox}{Rite of the Final Testament}
Cost: 13 XP, Obligation 7 segments \\
Effect: Seal a dying person's last request with weight beyond life. Targets must pass Resolve (DV 4) to ignore it. Breaking a just and witnessed request draws Varnek Karn's personal attention.
\end{patronbox}

\begin{patronbox}{Rite of the Eternal Archive}
Cost: 14 XP, Obligation 8 segments \\
Effect: Found a repository where the dead deposit knowledge. Gain a permanent +1 die to Lore for historical or genealogical research made within. Spirits periodically contribute testimonies---but the archive demands curatorship and payment.
\end{patronbox}

\textbf{GM Guidance --- Varnek Karn's Intrusions:}
\begin{enumerate}
    \item A ghost you once bargained with appears at an inopportune moment, demanding that you witness a new death or settle an old account.
    \item A corpse you handled begins to whisper in your sleep, reciting the terms of a contract you thought was closed.
    \item A funeral you attended is suddenly disputed---someone claims the deceased was not of sound mind when they made their final testament.
\end{enumerate}

\subsection*{Varnek Karn's Grave Debt (Corruption)}
Followers of Varnek Karn do not track Obligation in the usual way. Instead, they accumulate \textbf{Grave Debt}---a measure of how much they have borrowed from the silence. Use the following table in place of standard Corruption tracks.

\begin{longtable}{@{}p{1.2cm} p{5.5cm} p{5.5cm}@{}}
\toprule
\textbf{Tier} & \textbf{Benefit} & \textbf{Cost / Quirk} \\
\midrule
1 & \textbf{Death's Whisper:} +1 die to \textit{Notice} when the dead seek to communicate. & \textbf{Graveyard Pallor:} --1 die to social rolls with the living; your presence unsettles. \\
\midrule
2 & \textbf{Ancestral Insight:} Once/scene, gain +2 dice for family history or inherited troubles. & \textbf{Spectral Awareness:} Ever‑aware of nearby spirits; suffer \textbf{Fatigue 1} in places of many deaths. \\
\midrule
3 & \textbf{Medium's Gift:} Once/session, act as a flawless conduit for a spirit. & \textbf{Death-Touched:} Cold spots and small portents follow you; --1 die to ordinary social grace. \\
\midrule
4 & \textbf{Final Authority:} Once/scene, compel adherence to a true last request (Resolve DV 4). & \textbf{Corpse-Pale:} Visibly wan; --2 dice where warmth or vigor is required. \\
\midrule
5 & \textbf{Archive Keeper:} Once/session, access perfect recall of a historical event or family secret. & \textbf{Unquiet Presence:} The dead trail you; --1 die to stealth and to surprising the living. \\
\midrule
6+ & \textbf{Death's Confidant:} Once/session, negotiate with puissant death‑entities as an equal. & \textbf{Halfway House:} Mark +2 Obligation; risk becoming half‑withdrawn from the living world. \\
\bottomrule
\end{longtable}

\subsection*{Playstyle Notes}
Varnek Karn favors investigators, arbiters, and heralds of funerary law. Power grows through cultural sensitivity, shrewd bargaining, and careful record‑keeping; the price is a thinning tie to warmth and mortal ease. Expect mysteries of lineage, obligations that outlive their makers, and parleys where both sides of the veil demand their due.

\noindent \textbf{Varnek Karn emphasizes:}
\begin{itemize}
    \item \textbf{Diplomacy over Dominion:} bargaining with the dead rather than commanding them.
    \item \textbf{Knowledge Brokerage:} trading for secrets and testimonies beyond the grave.
    \item \textbf{Cultural Protocols:} rites, tokens, and proper doors for speaking with ancestors.
    \item \textbf{Mediation:} resolving conflicts that bind the living to the dead.
    \item \textbf{Historical Continuity:} showing how past vows shape the present.
\end{itemize}

\clearpage

\section{Summary of Patrons by Category}

\begin{center}
\begin{longtable}{p{3cm} p{5cm} p{5cm}}
\caption{Patrons of Fate's Edge}\\
\toprule
\textbf{Category} & \textbf{Patron} & \textbf{Domain} \\
\midrule
\endfirsthead
\multicolumn{3}{c}{{\bfseries \tablename\ \thetable{} -- continued}} \\
\toprule
\textbf{Category} & \textbf{Patron} & \textbf{Domain} \\
\midrule
\endhead
\bottomrule
\endfoot

The Wild & Grimmir & Nature, seasons, wilderness \\
      & Raéyn & Sea, tides, storms \\
      & Carrion-King & Decay, renewal, compost \\
\midrule
The Shadow & Ikasha & Shadow, secrets, thresholds \\
      & Isoka & Transformation, shedding, venom \\
      & Aveh & Freedom, erasure, storms \\
\midrule
The Covenant & Inquisitor Prime & Purity, judgment, order \\
      & Witness & Truth, memory, revelation \\
      & Traveler & Roads, journeys, safe passage \\
      & Mykkiel & Law, covenant, judgment \\
\midrule
The Necessary Evils & Malachai & Curses, corruption, false salvation \\
      & Inaea & Webs, fate, connection, mercy \\
      & Varnek Karn & Death, contracts, ancestral wisdom \\
\midrule
The Infinite & Khemesh & Depth, pressure, terror \\
      & The Ninth & Hyperknowledge, paradox, pattern \\
\bottomrule
\end{longtable}
\end{center}

\noindent \textit{Note: Mykkiel appears under The Covenant, while Malachai, Inaea and Varnek Karn form a category reflecting the darker, more intimate bonds of obligation and consequence.}

\clearpage
\clearpage
\chapter{Free Casting -- The Grammar of the Raw Weave}
\epigraph{``The Six Roads are paved stone. Free casting is the mud between the cobbles, the river beneath the bridge, the air in the lung. It has no patron to catch you when you fall. It has no leash to pull you back. The risk is not in the casting---it is in the caster.''}{\textit{--- The Grey Wanderer}}

\section{The Seventh Way}

The Six Roads are mapped, signed, and toll-paid. But beneath them runs a current older than the first carved Rite, older than the first sworn Symbol, older perhaps than the Patrons themselves. The Aelinnel call it \textit{Contextual Recursion}. The Thepyrgosi Synod calls it \textit{practical metaphysics}. In the hedgerows and burned-out tenements, they call it simply \textit{the knack}, \textit{the gift}, or \textit{the curse}.

It is Free Casting: the direct manipulation of the Weave through will, word, and gesture, unmediated by the covenants that govern the Six Roads. A Runekeeper channels her Patron's obligation; a Free Caster reaches into the raw stuff of reality and \textit{persuades} it to change. The difference between a hearth-witch burning her fingers lighting a candle with a thought and a Thepyrgosi magnus collapsing a wall with a gesture is not the \textit{kind} of power, but the \textit{discipline} of the mind that wields it.

\section{The Thepyrgosi Lament -- On Unsanctioned Magi}

\epigraph{``They call it 'free casting.' I call it a lit match held to a powder keg. No Codex. No Symbol. No Patron to rein them in. Just raw will and a prayer that the Backlash doesn't eat them---or the village.''}{\textit{--- Inquisitor Voss, Thepyrgos Witch-Hunter Commissariat}}

Among the tower-clad streets of Thepyrgos, where the Synod's bells measure law in chimes and every stair climbs toward order, the free caster is the most feared and hunted of all magi. Not because they are evil---though some are---but because they are \textit{unaccountable}. A Runekeeper's power is chained to a Codex and a Thiasos; if she burns a district, the Patron's name is written in the ash, and the weight can be called. An Invoker's magic passes through a Symbol, a named Patron, an obligation recorded in a covenant that can be seized. But a free caster? They leave no receipt. The Backlash is theirs alone to bear, and when they fail, the Weave lashes out without theme or mercy.

Yet even here, the risk varies. The Inquisition does not hunt the University don who practices Contextual Recursion in a sealed laboratory, surrounded by wards and graduate students. They hunt the gutter-mage who learned three tags from a dying grandmother and now improvises storms to pay his rent. The don is \textit{managed} risk; the gutter-mage is \textit{wild} risk. Both are free casters. The Weave does not distinguish, but the Synod must.

\section{The TAGS Philosophy}

Free casting is built from \textbf{TAGS}: the verbs, nouns, and adjectives of the Weave's hidden grammar. \textbf{Intent} (what you wish to do) + \textbf{Target} (what you wish to affect) + \textbf{Tags} (the qualities you impose). Base Difficulty equals 1 plus the number of TAGS employed. Dangerous TAGS---those that manipulate space, creation, or fate itself---add +2 to the DV instead of +1.

Casting roll: \texttt{Attribute + Arcana} versus the DV. Each \texttt{1} rolled generates a Story Beat (SB), representing the Weave's resistance to being forced into shape.

\section{Tag Lookup Reference}

{\small
\begin{longtable}{@{} p{2.4cm} p{1.1cm} p{2.4cm} p{6.2cm} @{}}
\caption{The Grammar of TAGS}\\
\toprule
\textbf{Tag} & \textbf{DV mod} & \textbf{Category} & \textbf{Effect summary} \\
\midrule
\endfirsthead
\multicolumn{4}{c}{{\bfseries \tablename\ \thetable{} -- continued}} \\
\toprule
\textbf{Tag} & \textbf{DV mod} & \textbf{Category} & \textbf{Effect summary} \\
\midrule
\endhead
\bottomrule
\endfoot

% Elemental TAGS
\abilitytag{Burning} & +1 & Elemental & Ignite, heat, combustion, smoke \\
\abilitytag{Freezing} & +1 & Elemental & Ice, slowing, brittle shatter, cold \\
\abilitytag{Storm} & +1 & Elemental & Lightning, shock, arc, thunder \\
\abilitytag{Stone} & +1 & Elemental & Walls, spikes, tremors, armor \\
\abilitytag{Wave} & +1 & Elemental & Crushing water, currents, pressure \\
\abilitytag{Wind} & +1 & Elemental & Levitation, gusts, deflection, push/pull \\
\midrule

% Force TAGS
\abilitytag{Force} & +1 & Force & Kinetic power, shields, blasts, telekinesis \\
\abilitytag{Area} & +1 & Force & Cone, circle, corridor, zone effect \\
\abilitytag{Strike} & +1 & Force & Single target precision \\
\abilitytag{Wall} & +1 & Force & Barrier or blockade \\
\abilitytag{Bind} & +1 & Force & Restrain, hold, suspend, entangle \\
\abilitytag{Dispel} & +1 & Force & Suppress magic, unravel ongoing effects \\
\midrule

% Mind \& Veil TAGS
\abilitytag{Veil} & +1 & Mind/Illusion & Conceal, blur, illusion, silence \\
\abilitytag{Scry} & +1 & Mind/Illusion & Reveal hidden, see distance, read traces \\
\abilitytag{Memory} & +1 & Mind/Illusion & Erase, alter, restore memories \\
\abilitytag{Command} & +1 & Mind/Illusion & Compel a short action (one word) \\
\abilitytag{Fear} & +1 & Mind/Illusion & Panic, flee, break morale \\
\midrule

% Life \& Body TAGS
\abilitytag{HEAL} & +1 & Life/Body & Close wounds, restore flesh, reduce Harm 1 \\
\abilitytag{Purify} & +1 & Life/Body & Remove poison, corruption, disease \\
\abilitytag{Strengthen} & +1 & Life/Body & Enhance body, armor, senses (temporary) \\
\abilitytag{Waken} & +1 & Life/Body & Counter sleep, paralysis, stun \\
\abilitytag{Beast} & +1 & Life/Body & Speak with or influence animals \\
\midrule

% Space \& Motion (always dangerous)
\abilitytag{Leap} & \dangerous{+2} & Space/Motion & Jump far, blink across short space (Near) \\
\abilitytag{Fold} & \dangerous{+2} & Space/Motion & Short-range teleport, vanish--reappear (Far) \\
\abilitytag{Gate} & \dangerous{+2} & Space/Motion & Long distance passage, open/close path \\
\abilitytag{Gravity} & \dangerous{+2} & Space/Motion & Crush, lift, suspend, walk on walls/ceiling \\
\midrule

% Creation \& Transformation (always dangerous)
\abilitytag{Create} & \dangerous{+2} & Creation & Manifest mundane matter briefly (1 scene) \\
\abilitytag{Summon} & \dangerous{+2} & Creation & Call a being or construct (see Summoning rules) \\
\abilitytag{Transmute} & \dangerous{+2} & Creation & Turn one thing into another (temporary) \\
\abilitytag{Animate} & \dangerous{+2} & Creation & Make objects act with intent (1 scene) \\
\midrule

% Additional Utility TAGS
\abilitytag{Sense} & +1 & Utility & Detect presence of a named tag/element \\
\abilitytag{Reveal} & +1 & Utility & Unveil hidden, glamoured, or invisible things \\
\abilitytag{Light} & +1 & Utility & Create illumination (glow, torch-bright) \\
\abilitytag{Shadow} & +1 & Utility & Deepen darkness, hide edges, obscure \\
\abilitytag{Silence} & +1 & Utility & Suppress sound in zone or on target \\
\abilitytag{Protect} & +1 & Utility & Reduce/deflect next harm (Armor 1) \\
\midrule

% Reaction / Conditional
\abilitytag{Counter} & +1 & Reaction & Interrupt a casting/ritual in its window \\
\abilitytag{Reflect} & \dangerous{+2} & Reaction & Turn next targeted effect back on its source \\
\abilitytag{Store} & \dangerous{+2} & Utility & Bank 1--2 successes in a vessel (once) \\
\abilitytag{Curse} & \dangerous{+2} & Affliction & Attach hostile tag/timer to target \\
\abilitytag{Bless} & +1 & Affliction & Grant favourable tag (luck, favor, ward-key) \\
\bottomrule
\end{longtable}
}

\section{Backlash -- The Weave's Revenge}

Without a Patron to absorb the strain, without a Thiasos to bleed off the excess, the Free Caster's failures strike the caster directly. The severity depends not on the spell, but on the \textit{cleanliness} of the casting---a scholar's precise ritual may suffer only a nosebleed, while a desperate improvisation may shatter bones.

\begin{tcolorbox}[colback=warningred!5, colframe=warningred!70, sharp corners, title={\textbf{Backlash Severity}}]
\begin{multicols}{2}
\textbf{Minor (1--2 SB)} \\
$\bullet$ Fatigue +1 \\
$\bullet$ --1 die on next roll \\
$\bullet$ Increase DV of next action by 1 \\
$\bullet$ Minor environmental effect (smoke, chill) \\
$\bullet$ GM gains +1 SB for a thematic twist

\columnbreak

\textbf{Major (3--4 SB)} \\
$\bullet$ Harm 1 (stress) \\
$\bullet$ Condition applied (Blinded, Silenced, Bound) \\
$\bullet$ Ward/barrier collapses \\
$\bullet$ Unwanted attention (Alarm timer advances) \\
$\bullet$ GM gains +2 SB to escalate the scene
\end{multicols}

\textbf{Catastrophic (5+ SB)} \\
$\bullet$ Harm 2 \\
$\bullet$ Permanent Scar or Mark (glowing eyes, withered hand) \\
$\bullet$ Spell inverts (heal becomes harm, etc.) \\
$\bullet$ Localized reality fracture (time stutter, gravity reversal) \\
$\bullet$ Patron or rival entity manifests, drawn by the noise
\end{tcolorbox}

\section{The Three Faces of the Free Caster}

Not all unsanctioned magi are madmen or desperate gutter-mages. The risk varies with the practitioner.

\textbf{The Thepyrgosi Scholar} practices ``practical metaphysics'' in a salt-ringed laboratory. She knows the TAGS by their mathematical names, calculates probabilities, and accepts Backlash as a ``cost of research.'' Her magic is predictable, expensive, and boring---but safe relative to the alternatives.

\textbf{The Hedge-Witch} knows six reliable combinations passed down through her family. She uses them to find lost sheep, quiet storms, and HEAL pots. She has never cast a spell with more than three TAGS, and she sleeps soundly. Her risk is low because her ambition is low.

\textbf{The Prodigy} was born knowing the grammar. He improvises tags he has never been taught, pulls effects from the air like a musician humming a new tune. He is the most dangerous of all---not because he wishes harm, but because he does not know which combinations are forbidden until he tries them, and the Weave remembers his mistakes longer than his successes.

\section{Example Spells}

\spell{Ember Flick}{\abilitytag{Burning}, \abilitytag{Strike}}{2}{Ignite a small, unattended object within Close range; deal 1 Fatigue to a living target.}{Minor Backlash -- spark jumps back; caster marks 1 Fatigue.}

\spell{Reed Walk}{\abilitytag{Wind}, \abilitytag{Leap}, \abilitytag{Self}}{3}{Leap up to Far distance in a straight line, ignoring opportunity attacks.}{Minor Backlash -- land off-balance; next action Desperate.}

\spell{Grasping Roots}{\abilitytag{Earth}, \abilitytag{Bind}, \abilitytag{Area}}{3}{Roots erupt in a Near zone; all within test Body+Athletics (DV 3) or become Entangled.}{Minor Backlash -- roots also snag an ally; they lose 1 step of Position.}

\spell{Fate's Needle}{\abilitytag{Fate}, \abilitytag{Strike}, \dangerous{\abilitytag{Divination}}}{5}{Roll 3d10, keep the highest; that many successes auto-succeed on your next non-magical action.}{Major Backlash -- caster marks Harm 1 and begins a \textit{Fate's Burden} timer [4].}

\section{Building Your Own Spells}
\begin{enumerate}
    \item \textbf{Intent:} ``I want to freeze the door shut.''
    \item \textbf{Choose TAGS:} \abilitytag{Freeze}, \abilitytag{Bind}, \abilitytag{Wall}, \abilitytag{Touch}.
    \item \textbf{Count TAGS:} 4 $\rightarrow$ Base DV = 1+4 = 5.
    \item \textbf{Adjust for scope:} Large door DV remains 5. Entire fortress gate? DV +2.
    \item \textbf{Roll:} Wits+Arcana vs DV 5. Each 1 = SB.
    \item \textbf{Apply Backlash per table.}
\end{enumerate}

\vspace{0.5cm}
\noindent \textit{Remember: the Weave is not malicious. It is simply precise. It lends its strength to those who speak its language clearly, and it punishes those who mumble.}

\clearpage
\chapter{Summoning -- The Art of the Opened Door}
\epigraph{``I have shaken hands with the drowned, signed contracts in cobweb ink with things that wear my grandfather's face, and once---only once---borrowed a cup of sugar from a demon. The door is the easy part. It is the conversation that follows that will unmake you.''}{\textit{--- The Grey Wanderer}}

Of the Six Roads, Summoning is the one most often mistaken for simple thaumaturgy. The uninitiated see only the salt circle, the smoking thurible, the muttered name, and assume the matter is one of command---of a master issuing orders to a slave. This is the delusion that gets novices killed. Summoning is not command; it is \textbf{diplomacy with teeth}. It is the art of opening a door into a crowded room and convincing something older than stone to step through on terms that do not end with your skin being hung out to dry.

Summoners do not carry Obligation as the Runekeepers do, nor do they suffer the Cantor's creeping Corruption. Their currency is simpler and more terrible: \textbf{negotiation}. Every spirit has a price---blood, memory, silence, a name, a promise---and the Summoner must weigh whether the service rendered is worth the cost exacted. Some cultures approach this as a respectful parley between equals. Others view it as domination---the spirit bound by chains of contract or force of will, compelled to serve. Both paths are valid. Both are dangerous. The spirit remembers which hand held the leash, and spirits have long memories.

\section{What Passes Through the Door}
To speak of spirits is to speak of everything that is not quite matter. The dead who cling to their final moments are spirits. The \textit{Neighbors} who live in the gaps between maps are spirits. The animistic intelligences that wake in deep stone and old trees are spirits. The beings that theologians call angels and demons---guardians and predators of the threshold---are spirits. Even the conceptual entities, the personifications of plague or fortune or the space between thoughts, are spirits.

They are as diverse as the cultures that name them:
\begin{itemize}
    \item \textbf{The Indigenous} (Aeler): Not summoned but \textit{awakened}---spirits native to the stone, the root, the deep places. They do not come from elsewhere; they rise from here.
    \item \textbf{The Ancestral} (Ykrul): The honored dead, blood-bound to clan and drum, returning with counsel or wrath.
    \item \textbf{The Drowned} (Zakov, Linn): Those who died in water, in vengeance, or in betrayal, now liquid and hungry.
    \item \textbf{The Neighbors} (Lethai, Aelaerem): The fair folk, the hedgerow-dwellers, the ones who were here before the current age and will remain after.
    \item \textbf{The Outsiders}: Things from the spaces between stars, between thoughts, between seconds---entities that do not belong in the world's architecture and strain it merely by existing.
\end{itemize}

\section{The Two Philosophies}

\subsection*{The Path of Negotiation}
Prevalent among the Linn, the Aelaerem, and the free ports of Zakov, this approach treats the spirit as a \textbf{guest}. The Summoner offers hospitality---salt, bread, a name spoken in trust---and the spirit offers service in return. The relationship is ongoing. A Linn tide-witch knows her selkie by name; they have history, grudges, inside jokes. When the Leash fills and the spirit acts to its nature, it might simply depart in a huff, or it might deliver a warning wrapped in riddles. The danger is intimacy. Treat with a spirit long enough, and you begin to see its point of view. Begin to see its point of view, and you may find yourself defending its interests against your own kind.

\subsection*{The Path of Domination}
Favored in Ecktoria, the Aeler Deep Vaults, and the martial traditions of Viterra, this method views spirits as \textbf{resources}---dangerous, certainly, but manageable through strength of will, precise geometry, and the ironclad contract. The Ecktorian Inquisition binds demons in circles of blessed salt and silver chains, treating them as prisoners to be interrogated. The Aeler awaken vein-spirits and immediately collar them with runic law, treating them as mining equipment with opinions. Here, the Leash is a choke-chain, and when it fills, the spirit lashes out with hatred. The danger is rebellion. A bound spirit waits for the circle to waver, the name to be misspoken, the master's attention to lapse. And then it collects years of deferred rage in a single, bloody instant.

\section{Cultural Traditions}

\subsection*{Aeler -- The Stone-Whisperers}
In the deep holds, Summoning is not eviction but \textit{employment}. The spirits of stone are indigenous---they do not travel, they merely sleep. The Rite of the Carved Contract involves inscribing the terms onto slate with diamond dust, then tapping the stone with a bell hammer. If the stone warms, the spirit wakes and enters negotiation. Most Aeler Summoners are also Runekeepers, blending the roads; they view spirits as mineral resources with legal rights. The spirits themselves are strange, patient, and literal---ask a vein-spirit to ``find gold,'' and it may collapse the mountain to bring you a single nugget.

\textbf{Convention:} Domination through legal precedent. The spirit is bound by the letter of the contract, but Aeler Summoners are meticulous in their wording.

\subsection*{Ykrul -- The Bone-Singers}
The steppe orcs summon the dead and the animal-totems through rhythm. The drum is the focus---not made of any hide, but of the specific creature or ancestor being called. To summon Wolf-Ancestor, one beats a drum of wolf-skin stretched over bone. The spirit arrives proud, angry, and expecting respect. The negotiation happens in the drumming; the Summoner must match the heartbeat of the spirit, proving kinship through rhythm. Failure to establish this resonance results in a spirit that hunts the caller.

\textbf{Convention:} Negotiation through kinship and respect. You cannot dominate an ancestor; you can only remind them of their obligation to blood.

\subsection*{Linn -- The Tide-Riders}
The northmen bind sea-spirits---selkies, drowned captains, storm-wardens---using the Rite of the Nine Knots. A hemp rope is tied with nine knots, each representing a wave of a storm surge. The spirit is not commanded but \textit{asked}, and the payment is usually a service to the sea: a net mended, a drowning prevented, a corpse returned to the deep. The relationship is covenantal; a Linn Summoner often keeps the same spirit for decades, learning its moods and its preferred tides.

\textbf{Convention:} Pure negotiation. To dominate a sea-spirit is to invite the storm. The Linn consider such practices barbaric.

\subsection*{Zakov -- The Salt-Binders}
The pirate-fleets of Zakov specialize in the vengeful dead---mutineers, betrayed lovers, captains who went down with their gold. The Fathom-Callers use \textit{blood coins}: silver pieces that have paid for murder. The coin is swallowed (marking the body as a vessel for the drowned), and the spirit is lured up through the salt water in the Summoner's own lungs. This is dangerous, intimate work. The spirits are angry, confused, and often malicious. The Zakov Summoner must be part psychologist, part exorcist, part undertaker.

\textbf{Convention:} Domination through temptation and binding. The spirit is tricked or coerced, held by the weight of its own unfinished business.

\subsection*{Lethai -- The Courtesy of Neighbors}
The elves do not summon spirits so much as \textit{invite them over}. To the Lethai, the Fair Folk are simply neighbors who live in a different wing of the house. Summoning involves leaving a place at the table, an empty chair, a cup of wine poured and left untouched. The spirits who arrive are capricious, powerful, and easily offended. The negotiation is conducted through etiquette---tiny gestures, the angle of a bow, the precise number of times a bell is rung (never nine).

\textbf{Convention:} Negotiation through absolute adherence to archaic courtesy. Domination is impossible; the only leash is mutual respect, and that leash is made of glass.

\subsection*{Viterra -- The Legal Binding}
The hedge-wrights of Viterra approach Summoning as property law. They draft contracts enforceable by the Crown, binding spirits as "indentured servants" for specific terms and tasks. A ghost may be bound to a threshold for nine generations to resolve a boundary dispute. The spirits are categorized by the College of Scriveners: Class I (ancestral, manageable), Class III (animistic, unpredictable), Class VII (demonic, forbidden without license).

\textbf{Convention:} Bureaucratic domination. The spirit is bound by legal fiction, and the contract is enforced by the state's metaphysical weight.

\subsection*{Ecktoria -- The Chained and the Burning}
The Inquisition practices the most brutal form of domination: the summoning of infernal entities for interrogation and execution. The \textit{Chain-Lanterns} trap spirits in blessed iron and silver mirrors, forcing them to witness their own sins until they agree to serve. This is ugly, dangerous work. The spirits are not willing guests but prisoners, and each summoning is an act of spiritual warfare. The Leash here is literal---chains of inscribed metal that burn the spirit's essence.

\textbf{Convention:} Violent domination. The spirit is enemy, prisoner, and tool. The risk of rebellion is constant, and the Inquisition accepts collateral damage as holy necessity.

\section{The Leash and the Bargain}

Whether through negotiation or domination, the Summoner must maintain the \textbf{Leash}---the spiritual strain of keeping an Outsider in the material world. The Leash is not a moral debt but a \textit{tension}, a rope pulled tight between the summoner and the summoned. Each time the spirit acts against its nature (a gentle brownie forced to kill, a vengeful wraith forced to heal), each time it takes Harm, each time it crosses a ward, the Leash ticks tighter.

When the Leash fills, the spirit acts once according to its essential nature---often destructively, always truthfully---then departs. If it was treated with respect, it may return when called again. If it was dominated, it may hunt the summoner through dreams until the price is paid in blood.

\subsection*{The Pact-Bound Spirit}
Some Summoners establish long-term relationships with specific spirits, negotiating \textbf{Pacts}---standing agreements that reduce or eliminate the need for the Leash during routine service. A Pact might stipulate that the spirit guards the summoner's home in exchange for a monthly offering of honey and blood. While the spirit works within the terms of the Pact, the Leash does not tick. But Pacts are living contracts; violate the terms, and the spirit is owed triple in return.

\section{Sample Spirits}

\begin{spiritcard}{The Vein-Serpent (Aeler)}
\begin{itemize}
    \item \textbf{Nature:} Indigenous, mineral, slow-witted.
    \item \textbf{Appearance:} A serpent of polished copper and malachite, slithering through stone as if it were water.
    \item \textbf{Services:} Locates veins of ore, stabilizes tunnels, collapses enemy fortifications.
    \item \textbf{Price:} A promise to maintain the stability of the stone; for every ounce of ore taken, a ounce of mortar must be laid.
    \item \textbf{If Dominated:} It twists the stone, creating sinkholes beneath the summoner's feet.
    \item \textbf{If Respected:} It warns of earthquakes three days in advance by humming through the bedrock.
\end{itemize}
\end{spiritcard}

\begin{spiritcard}{The Wolf-Ancestor (Ykrul)}
\begin{itemize}
    \item \textbf{Nature:} Ancestral, protective, prideful.
    \item \textbf{Appearance:} A grey wolf the size of a pony, with eyes that reflect scenes from the summoner's childhood.
    \item \textbf{Services:} Tracking across any terrain, guarding camps, teaching forgotten battle-songs.
    \item \textbf{Price:} The summoner must uphold the honor of the clan; any cowardice breaks the bond immediately.
    \item \textbf{If Dominated:} It hunts the summoner's pack-mates first, testing their worthiness through blood.
    \item \textbf{If Respected:} It grants dreams of the future, though the dreams arrive in the language of scent and instinct.
\end{itemize}
\end{spiritcard}

\begin{spiritcard}{The Selkie-Maid (Linn)}
\begin{itemize}
    \item \textbf{Nature:} Sea-spirit, mournful, capricious.
    \item \textbf{Appearance:} A seal that sheds its skin to become a woman with hair like kelp and eyes like storm clouds.
    \item \textbf{Services:} Navigation through impossible storms, drowning enemies with summoned currents, retrieving objects lost to the deep.
    \item \textbf{Price:} A tear cried into the sea, or a memory of happiness sacrificed to the salt.
    \item \textbf{If Dominated:} She leads ships onto rocks in clear weather, humming lullabies that deafen the crew.
    \item \textbf{If Respected:} She brings word of currents and fish schools, and may save a drowning man for no reason other than whim.
\end{itemize}
\end{spiritcard}

\begin{spiritcard}{The Drowned Mutineer (Zakov)}
\begin{itemize}
    \item \textbf{Nature:} Vengeful dead, confused, bitter.
    \item \textbf{Appearance:} A bloated figure in rotting finery, surrounded by a halo of salt-crusted coins.
    \item \textbf{Services:} Sabotage of ships, whispering secrets into sleeping ears, cursing cargo to spoil.
    \item \textbf{Price:} A promise of vengeance against the living; the summoner must carry the mutineer's grudge as their own.
    \item \textbf{If Dominated:} It turns on the summoner the moment the circle wavers, drowning them in their own lungs.
    \item \textbf{If Respected:} It serves with sullen efficiency but offers no warning of the backlashes its presence attracts (other drowned things follow it).
\end{itemize}
\end{spiritcard}

\section{The Wanderer's Warning}

Do not mistake the spirit for a tool. The Aeler think they do this, treating their vein-serpents as mining equipment, but I have seen those serpents wait centuries for a tunnel to collapse on a foreman who spoke slightingly of stone. The Ecktorian Inquisition thinks it dominates the infernal, but I have watched a "bound" demon whisper seven words to an Inquisitor that made him question his faith for the rest of his life.

The Leash is not a guarantee. It is a courtesy extended by the spirit while it finds your measure. Negotiate honestly. Dominate absolutely, if you must, but never sleep with both eyes closed. And remember: when the summoning ends, when the salt is swept up and the circle broken, the spirit remains. It remembers your face. It remembers your terms. And if you have treated it poorly, it will find you when your own Leash is frayed, and it will collect its due with interest.

\vspace{0.5cm}
\hfill --- \textit{The Grey Wanderer}

\clearpage
\chapter{The Cantor's Voice -- Echoes and Entropy}
\epigraph{``I have seen a woman sing a storm into stillness, and a man hum a king into madness. The voice is the oldest door, and the Cantor is the one who knocks without asking who is home.''}{\textit{--- The Grey Wanderer}}

\section{The Song Without Swearing}

If the Runekeeper is the sworn knight and the Witch the quiet architect, the Cantor is the \textit{wild card}---the burning letter sent without a seal. They are the magical bards, the storm-singers, the hymn-leaders who borrow power not by covenant or ledger, but by \textit{resonance}. They do not swear to Patrons; they \textit{echo} them. They stand in the marketplace or the tempest and shape their throats into instruments of divine will, knowing full well that the throat is flesh, and flesh warps under pressure.

Do not mistake the Cantor for a mere entertainer. Every culture has its minstrels, but the Cantor is something older and more dangerous: a vessel that chooses its own filling. In Aelaerem, they are the Hearth-Mothers who sing the hollows to sleep. In Linn, they are the Storm-Callers whose shanties calm the tide or summon it. In Zakov, they are the Siren-Priests of the Salt-Queen, whose songs bind sailors to pacts deeper than any written contract.

Like the Runekeeper, the Cantor wears many masks. A Runekeeper sworn to the Oath of Flame becomes a paladin; one sworn to Grimmir becomes a druid of the ledger-stone. So too does the Cantor shift shape. A singer of Palinode becomes a prophet of rapture, whose voice cures despair---or causes it. A devotee of Thrysos becomes a lord of revels, equal parts host and hostage to the feast. A chanter of the Silent Choir whispers truths that dismantle empires. They are all Cantors, though their rites could not be more different.

\section{The Price of the Song -- Corruption as Metamorphosis}

Where the Runekeeper pays in \textbf{Obligation} and the Invoker in \textbf{Condition}, the Cantor pays in \textbf{Corruption}---not a moral stain, but a \textit{physical} one. Corruption for a Cantor is not a debt to be settled; it is a \textit{transformation undergone}. The voice that sings too often to the storm begins to carry thunder in its timbre. The hands that conduct too many hymns of Isoka begin to shed skin like serpents. The eyes that reflect too deeply on the Ninth begin to see angles that do not exist.

This is the inevitable consequence of intimacy. The Cantor does not stand at the threshold like a Witch, nor bargain from a distance like an Invoker. They \textit{open their mouths} and let the power pass through, and some of it always remains, lodged in the teeth, the lungs, the marrow.

Most Cantors learn too late that the voice is not infinite. They push, and push again, accelerating their songs to immediate effect until the Corruption timer fills. Then comes the \textit{bloom}: a permanent alteration, neither wholly blessing nor curse, granted by the last Patron whose Rite they performed. They may gain the shadow-sense of Ikasha, but find themselves unable to tolerate sunlight. They may gain the healing touch of Inaea, but find themselves bound to HEAL every broken thing they encounter, even the wounds of enemies.

And then, often, comes the end.

\section{The Three Fates of the Cantor}

It is no accident that Cantors are often outcasts. Their power is public, loud, and inherently destabilizing. A Runekeeper's covenant can be checked; a Witch's knots are quiet; a Psion's mind is invisible. But a Cantor cannot hide. Their story tends toward one of three endings:

\begin{warningbox}
\textbf{Persecution} \\
The Cantor sings a truth the local theocracy does not wish to hear, or summons a storm that ruins the harvest, or simply draws too large a crowd. The Everflame names them heretic. The Chain-Lanterns of Thepyrgos hunt them as unlicensed magi. The pyre is the common end for the public singer who cannot---or will not---silence themselves.
\end{warningbox}

\begin{warningbox}
\textbf{Madness} \\
The voice becomes too heavy. The harmonics of the Ninth ring constantly in the ears. The Cantor forgets which words are theirs and which are the Patron"s. They sing until their throat bleeds, or until they walk into the sea following a voice only they can hear. Many a lighthouse keeper on the Kahfagian coast is a Cantor who chose the ocean over the asylum.
\end{warningbox}

\begin{warningbox}
\textbf{The Cult} \\
The rare Cantor who masters their Corruption, or learns to ride it, often finds themselves at the head of a following. The Serpent's Chorus of Zakov, the Shedded who follow Isoka's Copperhead, or the Bacchanalians of Thrysos---these are Cantors who became high priests not by ordination, but by \textit{attraction}. They are dangerous not because they command armies, but because they command \textit{desire}. To follow such a Cantor is to risk your own identity in their song.
\end{warningbox}

\section{The Patrons of Performance}

While any Patron might answer a Cantor's song, certain powers are drawn to the voice like moths to flame. They are the patrons of ecstasy, of revelation, of the mask and the encore.

\subsection{Palinode -- Queen of Encores}
\textit{Once a road-cantor who died between feasts with a chorus unfinished, Palinode refused to let the refrain end. She is the promise that the song never dies, the encore that becomes eternal. Her domain is Performance, Rapture, and the terrible beauty of the unending refrain.}

Palinode arrives where crowds lean forward and makers hold their breath. Her signs are a cup that never quite empties, a candle that gutters only when silence reigns, and a flawless bar of music that the hand cannot forget. She tempts with relief from solitude---an audience that always leans closer, a song that never dies, a bond that feels like family even if it devours you.

\textbf{Playing a Cantor of Palinode:}
You seek the perfect performance, the moment when the boundary between artist and art dissolves. Your rites grant power through music, dance, and the manipulation of emotion in crowds. You are never alone, but you may lose yourself in the crowd's embrace.

\begin{patronbox}{Rites of the Queen of Encores}
\textbf{Hymn Against Dread} (Low, 4 XP) -- Scene; Near; Resist only.\\
Allies gain +2 dice to resist Fear this scene. On a 10 rolled, they may ignore Fear entirely for one exchange.

\textbf{The Single True Note} (Low, 5 XP) -- Instant; Self; No.\\
For one scene, you cannot produce an imperfect sound: +2 dice to Performance, but −1 dice to all other rolls from intrusive perfection.

\textbf{The Circling Cup} (Low, 5 XP) -- Scene; Zone; No.\\
Create a revelry zone where time blurs. Participants recover Fatigue at twice normal rate. Leaving requires a Resolve test (DV 3) or another hour spent within.

\textbf{The Hooked Chorus} (Standard, 7 XP) -- Scene; Near; No.\\
\textit{Presence + Command} vs. Resolve (DV 3). Success: target joins in your performance. Partial: they resist but suffer −1 die to social rolls this scene. Failure: they follow, entranced.

\textbf{Consecrated Stage} (High, 13 XP) -- Extended; Zone; No.\\
Hallow a venue. All performances gain +2 dice and +1 Effect. The space is [WARD]ed; any disruption spawns 2 SB. The ward collapses if the performance ends.

\textbf{Crownpiece} (High, 14 XP) -- Extended; Touch; No.\\
Create a transcendent masterpiece. +2 dice to creative rolls this session; immune to social manipulation. But all lesser works falter---you suffer −1 die to mundane actions.
\end{patronbox}

\subsection{Thrysos -- King of Revels}
\textit{Where torches flare and cups spill, where dancers stamp until dawn, there strides Thrysos crowned with ivy and horn. He is both host and beast, a patron of license who laughs at law. In his train follow the masked and the wild, the broken who dance to forget, and the bold who revel to remember.}

Thrysos promises freedom without end: no oaths, no chains, only song and cup and fire. Yet those who drink too deeply wake to find they serve not joy, but appetite. His followers master gifts that devour from within---transformation through intoxication, truth through wine, and the terrible clarity of the hangover that never ends.

\textbf{Playing a Cantor of Thrysos:}
You are the lord of the feast and the lord of the Bacchanal. Your magic is infectious, turning quiet gatherings into riots of emotion and physical transformation. You offer liberation, but your price is always the self-control of those who accept your cup.

\begin{patronbox}{Rites of the King of Revels}
\textbf{The Overflowing Cup} (Low, 4 XP) -- Scene; Near; No.\\
Allies who drink gain +1 die to social rolls; enemies who drink suffer −1 die. Push It: All in Near must test Resolve (DV 3) or be swept into revelry; mark 1 SB (Hearts).

\textbf{The Horned Dance} (Low, 5 XP) -- Scene; Self; No.\\
Enter ecstatic frenzy: +2 dice to melee or performance; ignore 1 Fatigue per round. Mark 1 Corruption if used more than once per session.

\textbf{Breaking Mask} (Standard, 7 XP) -- Instant; Near; No.\\
Reveal a target's true face: they suffer −2 dice to disguise, deceit, or composure until scene ends. Push It: They must declare one hidden desire; mark 1 SB (Diamonds).

\textbf{Wild Procession} (Standard, 9 XP) -- Scene; Zone; No.\\
Summon spectral satyrs and maenads. The zone becomes [WARD]ed: hostile entry forces 1 SB (Spades); allies gain +1 die to resist compulsion or fear.

\textbf{Gilded Stag} (High, 12 XP) -- Extended; Touch; No.\\
Transform one reveler into a beast of horn and fire: +2 dice melee; Harm +1 (Burn). At session's end, they collapse with Harm 2 (Exhaustion).

\textbf{Endless Feast} (High, 14 XP) -- Scene; Zone; No.\\
For one scene, the zone floods with plenty: allies restore Fatigue; all rolls tied to appetite or indulgence gain +2 dice. Leaving requires Resolve (DV 3) or the scene's end. Push It: The feast persists as a hunger-haunt; start Hedonist's Curse [6].
\end{patronbox}

\subsection{Other Voices}
Beyond these two lie other powers sympathetic to the Cantor's art:

\begin{itemize}
    \item \textbf{Livaea, the Crimson Courtier} (Seduction \& Social Binding): The whisper behind the curtain, the smile that sells secrets. Her Cantors are political assets and dangerous lovers.
    \item \textbf{The Silent Choir} (Restraint \& Protective Omission): Keepers of unspoken truths. Their Cantors sing not with voice but with \textit{absence}---silences that heal and silences that kill.
    \item \textbf{Mab, Queen of Courts} (Glamour \& Bargain): The smile that binds tighter than chains. Her Cantors weave enchantments of etiquette and fey obligation.
    \item \textbf{Malachai, the Cruel Messenger} (Curses \& Corruption): Though often counted among the Necessary Evils, her Cantors are the masters of the honeyed curse and the false hope.
\end{itemize}

\section{Singing the Rites -- Mechanics of the Cantor}

Cantors do not invoke Rites with symbols or codices. They \textbf{sing} them. This requires the \textbf{Cantor's Path} talent (Major, 8 XP), which grants access to Low Rites performed as Songs.

\subsection{The Performance}
To sing a Rite, roll \textbf{Lore + Performance} against the Rite's DV. Unlike standard casting, a Song typically requires one action to begin, resolving at the start of your next turn unless you \textbf{Push It}---accelerating the song to immediate effect at the cost of marking 1 Fatigue and advancing your Corruption timer.

\subsection{The Corruption Timer }
You possess a personal \textbf{Corruption Timer } with segments equal to your Spirit rating. Mark Corruption when:
\begin{itemize}
    \item You Push a Song to resolve immediately.
    \item You perform a Resonant Rite (a powerful thematic working designated by your Patron).
    \item The GM spends a Story Beat involving your occult activities.
\end{itemize}

When the Timer fills, you immediately gain a \textbf{thematic boon and drawback} from the last Patron whose Rite you performed (see their Corruption Tables). All your followers, retainers, or familiars also gain a trait of the same corruption. The Timer then resets to zero, but it \textit{cannot} go below cannot go below your current Tier -1.

\subsection{The Repertoire}
Unlike a Runekeeper's fixed Codex, a Cantor's \textbf{Repertoire} is fluid. Track a \textbf{Repertoire Timer } [6]. Mark a segment for each unique Song learned. At 2 segments, reduce the base DV of Songs by 1 (minimum 2). At 4 segments, gain +1 die to Performance. At 6 segments, learn one Standard Rite as a Song (temporary; requires ongoing practice to retain).

\subsection{The High Cantor}
Those who reach Tier II+ may take the \textbf{High Cantor} prestige talent (18 XP), allowing you to weave Standard Rites into instant, powerful Songs. Each such casting marks your Corruption Timer , but the effects are immediate and devastating: lightning called with a scream, armies rallied with a single note, or death itself bargained with in a lullaby.

\vspace{1cm}
\noindent \textit{If you must sing, sing carefully. The voice is a flame. Lean too close, and you will find yourself consumed---not by hellfire, but by the applause of the crowd, which is hungrier than any demon.}

\section{Embracing Corruption -- The Art of Blooming}

\epigraph{``You fear the feathers in your hair? Child, I have known Cantors who plucked them out and wept. I have known others who let them grow, who preened, who flew. The difference is not in the Patron's touch---it is in the courage to be touched.''}{\textit{--- The Grey Wanderer}}

Most Cantors fear the \textit{bloom}. They see the Corruption Timer as a countdown to obsolescence, a timer ticking toward madness or monstrosity. They are wrong. The timer is not a bomb; it is a chrysalis.

There exists a philosophy among the older voices, the ones who have sung themselves hoarse and come back speaking in harmonies: \textbf{embraced corruption}. These Cantors do not resist the metamorphosis---they accelerate it. They seek the bloom not as an accident but as an \textit{ascension}, trading their mundane flesh for a vessel more suited to the song.

\subsection{The Willing Metamorphosis}

To embrace corruption is to treat your body as an instrument requiring tuning. When the feather grows, you do not pluck it---you oil it. When the shadow lags, you learn to walk in time with its hesitation. When your tears turn to brine, you taste them and remember the depths you have touched.

\textbf{Playing an Embraced Cantor:}
You are not a victim of your Patron's influence; you are a collaborator. You mark Corruption not with dread but with anticipation. Your "flaws" are sacred stigmata. The mutations that others hide---the scales, the glowing eyes, the voice that echoes without source---you display as badges of seniority.

Mechanically, when you voluntarily fill your Corruption Timer (not through desperate Pushing, but through deliberate, ritual Resonant Rites), you may choose your corruption trait from your Patron's \textit{Corruption Table} rather than rolling randomly. The trade-off is that you must \textit{embody} the trait: a Cantor of Thrysos who chooses the "Ever-Thirsting" corruption must actually consume twice normal rations and cannot refuse offered drink, or their voice fails them until they indulge.

\subsection{The Threshold of Ascension}

There is a door that opens only when the Corruption Timer has been filled seven times. Few Cantors survive that long, and fewer still return from what lies beyond. Those who do speak of the \textbf{Second Voice}---not the Patron's echo, but their own, externalized, walking beside them as a separate entity.

The Embraced call this the \textit{Fugal Self}. It grants +1 die to all Performance rolls, but requires you to negotiate with your Corruption for control of your body each dawn. Fail, and the Corruption acts through you for a scene, pursuing its own alien desires.

\subsection{The Wanderer's Final Warning on Blooming}

I have seen Cantors who embraced their corruption so thoroughly they forgot which world they belonged to. The Serpent's Chorus began as men and women; now they are a nest of scales and whispered rumors, worshipped in basements by those who mistake their hissing for prophecy. They are powerful. They are immortal, in a fashion. But they are no longer \textit{people}.

Embrace the bloom if you must, but keep a knife sharp and a friend honest. The moment you begin to think the feathers are \textit{your} idea---that you wanted them, that they make you beautiful---is the moment you stop being a Cantor and become a \textit{vessel}. And vessels, when they crack, do not bleed. They only ring.

\vspace{0.5cm}
\noindent \textit{Sing if you must. Bloom if you dare. But remember: flowers open to be seen, and to be seen is to be picked.}

\subsection{The Talent: Embraced Corruption}

\textbf{Embraced Corruption} (Major, 12 XP) \\
\textit{Prerequisite:} Cantor's Path, Tier II+

You have learned to treat the bloom not as disease but as evolution. 
\begin{itemize}
    \item When you voluntarily fill your Corruption Timer through Resonant Rites (not by Pushing), you may \textbf{choose} your corruption trait from your Patron's Corruption Table rather than accepting a random result. You must embody this trait physically and socially (e.g., a Thrysos corruption requiring you to never refuse offered drink).
    \item You may Push Songs without marking Fatigue \textbf{once per session}, accepting the Corruption mark as the only cost.
    \item Once you have filled your Corruption Timer seven times (cumulative, not consecutive), you develop the \textbf{Fugal Self}: +1 die to all Performance rolls. However, each morning you must test Spirit + Resolve (DV 3) or your Corruption controls your body for the first scene of the day, acting according to its own alien desires.
\end{itemize}

\clearpage

\clearpage
\chapter{The Invoker's Ledger -- On Bargaining Without Bondage}
\epigraph{``I do not kneel. I sign. I do not pray. I calculate. The Patron receives my offering---a Symbol polished, a contract whispered---and in exchange, I am permitted to borrow the flame. But make no mistake: when the covenant closes, the weight is mine alone.''}{\textit{--- The Grey Wanderer}}

There exists a breed of magus who treats with power not as a supplicant, but as a \textit{factor}. They are the \textbf{Invokers}: the natural philosophers, alienists, and chancellors of the unseen who operate at the margins of the great houses and the edge of the law. Where the Runekeeper is a servant sworn to a single lord, the Invoker is a free agent---a gambler who has bet their soul on their own ability to read the fine print.

Invokerhood is not a path for the impulsive. It is the province of the methodical, the patient, and the desperate. An Invoker carries no Codex of carved Rites, no Thiasos to whisper warnings. They carry only \textbf{Symbols}---physical anchors to Patrons they have never fully sworn to---and a record of obligations meticulously maintained against the day of collection.

\section{The Natural Philosopher and the Alienist}
\lettrine{I}{n} the universities of Thepyrgos and the scriptoria of Ecktoria, they call themselves \textit{Natural Philosophers}, treating magic as a mechanical science to be plumbed and catalogued. In the asylums and the charnel houses, they are \textit{Alienists}, exorcists who bargain with the voices in the walls not with holy water, but with contracts written in iron gall ink. They are masters of the slow investigation, the careful excavation of dark secrets, and the precise application of arcane force.

An Invoker does not ask what a Patron \textit{wants}; they determine what a Patron \textit{lacks}. They are the traders of the occult, the middle-men between the mortal and the infinite, and they know to the penny what their soul is worth---and how to mortgage it twice over.

\section{The Symbol as Contract}
The Invoker's power flows through the \textbf{Symbol}: a ring, a seal, a blackened coin, or a bound sheaf of papers that serves as the physical locus of a bargain. Unlike the Runekeeper's Codex, which is a record of service, the Symbol is a \textit{contract made tangible}. Each Symbol represents a specific obligation owed to a specific Patron.

An Invoker may carry multiple Symbols, diversifying their portfolio of power like a merchant hedging against the fall of grain. But power is jealous. Each additional Symbol creates \textbf{Cross-Resonance}---a friction that makes the Rites harder to invoke and the obligations harder to track. Most Invokers limit themselves to four Symbols; beyond six, the ledger becomes unmanageable, and the weights begin to call in unison.

\subsection{The Slow Work and the Seal}
In their natural state, Invokers are slow and respectful. Their rituals take \textbf{Significant Time}---hours of preparation, geometric calculation, and the careful polishing of the Symbol. They do not cast; they \textit{invoke}, channeling power through the terms of the bargain. Each such working increases their Obligation by one segment, noted in the ledger with the precision of a notary.

But when the knife is at the throat---when the building burns or the demon is already through the door---the Invoker performs the \textbf{Crack the Seal}. They invoke the Rite instantly, without preparation, calling in the weight as an emergency loan. The price is never small: the Obligation increases by two or three segments immediately, and the Symbol may become \textbf{Compromised}---cracked, blackened, or weeping---until it can be ritually cleansed.

\section{Patrons of the Invoker}
Invokers gravitate toward Patrons who understand the language of exchange, debt, and hidden clauses. Below are two such entities, patrons of the contract and the quiet cost.

\subsection{The Confessor Beneath the Bell (Burden \& Exchange)}
\label{patron:confessor}

\textit{``Speak, and be freed. But understand---what you cast off, I will bear. And what I bear, you will one day inherit.''}

The Confessor waits in places where guilt gathers---under church spires, in confessionals long abandoned, in the echo of unspoken sins. He does not judge, nor does he absolve. He \textit{takes}. Each confession is a transaction: your burden for his silence, your pain for his weight. His followers learn to carry the weight of others, to shoulder the shame that would break lesser backs, and to use that accumulated sorrow as a currency of power.

\textbf{What They Represent:} Sacred confession, burden-sharing, the migration of sin, and the heavy silence of the confessional.

\textbf{Playing a Follower of the Confessor:}
\textbf{Strategy:} You are the party's emotional anchor and psychological investigator. Your Rites excel at absorbing the harm of others, extracting truth through burden-sharing, and binding agreements under the weight of shared secrets.

\textbf{Suggested Acts of Service:}
\begin{itemize}
    \item Absorb a secret that would ruin another, and carry it for a full season without speaking it.
    \item Offer a formal confession in a place of power, naming a sin that costs you social standing.
    \item Take on the physical illness or madness of another, bearing it in their stead for a day and a night.
\end{itemize}

\textbf{Patron's Gift (Imbuement):} Your presence feels like the hush after a bell toll. For one scene, you gain +1 die to Command and Insight, and you may transfer one Condition or 1 Fatigue from an ally to yourself as a free action.

\textbf{Alternative Gift -- The Silent Ledger:} Once per session, you may place a secret obligation upon a target. Name a truth they have never spoken aloud. Until they confess it to someone, they suffer −1 die to all actions related to deception or social maneuvering. This counts as your Patron's Gift for the scene.

\textbf{Rites -- Strategic Use:}

\textbf{Low Rites}
\begin{itemize}
    \item \textit{Echo of Sin} (4 XP): Absorb 1 Story Beat from an ally or NPC within Near range, converting it into +1 Obligation for yourself.
    \item \textit{Burden Transference} (5 XP): Transfer one Condition, Fatigue, or 1 Harm from another target to yourself.
\end{itemize}

\textbf{Standard Rites}
\begin{itemize}
    \item \textit{Speaking Knot} (8 XP): Bind a promise with faerie knots that tighten when broken. Breach forces 2 SB and marks the violator with a visible tell.
    \item \textit{The Unfinished Promise} (7 XP): Frame a bargain with an escape clause only you understand. Gain +1 die when the other party acts to fulfill their obligation; you may interpret ambiguous terms in your favor.
\end{itemize}

\textbf{Advanced Rites}
\begin{patronbox}{The Cauldron's Secret}
\textbf{Cost:} 13 XP, Obligation 7 segments \\
\textbf{Effect:} Brew a solution to any problem, but the cauldron asks one more ingredient---one the brewer does not realize they are providing until too late. The remedy works perfectly but breeds dependence.
\end{patronbox}

\textbf{GM Guidance -- The Confessor's Intrusions:}
\begin{enumerate}
    \item A secret you confessed years ago appears written in ash on your doorstep.
    \item You begin to manifest the physical symptoms of a disease you once absorbed for another.
    \item A bell rings in a soundproofed room. The Confessor is listening, and he finds your accounts out of balance.
\end{enumerate}

\clearpage

\subsection{Morag the Hag -- Weaver of Hidden Costs}
\label{patron:morag}

\textit{``She gives you the gold, but not the weight it carries. She grants your heart's desire, but not the hunger it awakens. She answers your plea, but not the price that pleases her. In Morag's court, every 'yes' is a 'yes, but...' ''}

In the hour between twilight and dawn, where borders thin and mortal cunning meets older cleverness, Morag sets her webs of bargain and consequence. She is the grandmother who offers exactly what you need, the forest crone with a perfect solution, the crossroads dweller who makes dreams come true---for a price that only reveals itself once the road bends back. Her wisdom is simple and sharp: every gift carries obligation, every kindness creates a burden, and every wish must be paid in coin not named aloud.

\textbf{What They Represent:} Faerie bargains, hidden obligations, the unspoken cost of every boon, and the crooked thread of fate.

\textbf{Playing a Follower of Morag:}
\textbf{Strategy:} You are the ultimate pragmatist, the dealer who always takes a cut. Your Rites grant immediate benefits with deferred costs. You excel at negotiation, at finding the loophole, and at ensuring that every favor you do creates a weight that will one day be collected.

\textbf{Suggested Acts of Service:}
\begin{itemize}
    \item Broker a trade where you give away something of seemingly lesser value that proves crucial later.
    \item Accept a gift from a stranger without asking the price, then discover what you have surrendered.
    \item Mend a broken oath by weaving a new, more binding one, ensuring the debtor is now yours.
\end{itemize}

\textbf{Patron's Gift (Imbuement):} Your touch leaves a faint scent of hawthorn and old coins. For one scene, you gain +1 die to Sway and Deception, and any bargain you strike creates a \textbf{Promise Timer [4]} that ticks toward a favor you may call in later.

\textbf{Alternative Gift -- The Crooked Thread:} Once per session, you may declare that a target owes you a weight from a previous, forgotten interaction. They must test Spirit + Resolve (DV 4) or accept the truth of the weight. If they fail, they owe you one significant service or truth. This counts as your Patron's Gift for the scene.

\textbf{Rites -- Strategic Use:}

\textbf{Low Rites}
\begin{itemize}
    \item \textit{Borrowed Luck} (5 XP): Share in another's fresh success. Gain +2 dice on one roll that benefits from their good fortune, but the original bearer suffers misfortune later.
    \item \textit{The Unfinished Promise} (4 XP): Frame a bargain with an escape clause only you understand.
\end{itemize}

\textbf{Standard Rites}
\begin{itemize}
    \item \textit{The Speaking Knot} (8 XP): Bind a promise with faerie knots. Both parties gain +1 die while keeping the bargain; breaking it escalates consequences (Harm 1, then Harm 2, then Curse).
    \item \textit{The Threefold Exchange} (7 XP): Broker a trade where each party yields hidden value. You claim the difference as Morag's tithe.
\end{itemize}

\textbf{Advanced Rites}
\begin{patronbox}{The Crossroads Court}
\textbf{Cost:} 14 XP, Obligation 8 segments \\
\textbf{Effect:} Convene Morag's court at a true crossroads at twilight. All present must bargain; each pact grants a marked boon and hides a cost that will ripen at the worst time.
\end{patronbox}

\textbf{GM Guidance -- Morag's Intrusions:}
\begin{enumerate}
    \item A deal you thought closed resurfaces with a clause written in a language you never learned but suddenly understand.
    \item The price of a past favor comes due: a finger, a memory, or a firstborn's laughter.
    \item You find yourself compelled to grant a favor to an enemy, bound by a thread you tied in your sleep.
\end{enumerate}

\section{Roleplay Hooks for the Invoker}
\begin{itemize}
    \item You carry four Symbols, but one was stolen by a former student who now uses your patron's power against you. The weight, however, still accrues to \textit{your} ledger.
    \item You have never met your Patron---you found the Symbol in a ruin. Does the Patron even know you exist? The letters you burn in your fireplace suggest they are learning.
    \item You are being audited by a rival Invoker's syndicate. They claim your interpretations of the Cross-Resonance clauses are heretical, and they intend to collect your Symbols by force.
\end{itemize}

\vspace{0.5cm}
\noindent The Invoker walks the tightrope between power and perdition. They are the proof that magic need not be served with devotion, only with precision. But remember: the Patron is not your friend. They are your creditor. And when the ledger finally balances, the interest is always paid in flesh.

\begin{warningbox}[title={\scshape The Triangular Trap \& The Tidal War}]
    \subsection*{On Carrying Conflicting Obligations}
    
    To carry one Symbol is to hold a candle. To carry five is to juggle torches in a powder magazine. But to carry Symbols of \textbf{Livaea}, \textbf{Maelstraeus}, and \textbf{Mor'iraath}---or to meddle with both \textbf{Raéyn} and \textbf{Khemesh}---is to strike sparks against the flint of your own ribcage.
    
    \subsubsection*{The Triangular Trap}
    \textit{Livaea} binds through desire; she makes you \textit{want} the weight. \textit{Maelstraeus} binds through accounting; he makes the weight \textit{exact}. And \textit{Mor'iraath}? He offers the illusion of cancellation.
    
    I have seen Invokers attempt to resolve a crushing obligation to Maelstraeus by calling upon Mor'iraath to unmake the creditor. It does not work. Mor'iraath does not erase the record; he erases the \textit{signatory}. The weight remains, now held by a void that cannot be paid, and the Obligation transfers to your bloodline as a \textit{curse of erasure}. Your children will be born with blanks in their names where your ancestor's face should be.
    
    Likewise, using Livaea's seduction to soften Maelstraeus's hard terms is like sweetening poison with honey. The Crimson Courtier's gifts always exact a social price, and the Infernal Bargainer's contracts always compound interest. You will find yourself owing both a firstborn's laughter and a firstborn's weight in silver, simultaneously.
    
    \subsubsection*{The Tidal War}
    Never---and I speak from the deck of a ship that was nearly crushed between them---never invoke \textbf{Raéyn} and \textbf{Khemesh} in proximity.
    
    Raéyn is the surface: the trade wind, the fair tide, the mother who grants safe passage. Khemesh is the depth: the trench without light, the son who swallowed the horizon. They are bound by blood and hatred. To invoke Raéyn's blessing on a vessel while carrying a weight to Khemesh is to invite a rogue wave that rises like a wall and crashes with a silence that tastes of iron. To invoke Khemesh's pressure while sailing under Raéyn's protection is to find your lungs filling with saltwater the color of storm clouds.
    
    If you must deal with the Sea Court, choose one: the Tidemark or the Trench. Never both.
    
    \vspace{0.3cm}
    \noindent\textbf{The Rule:} If your Patrons" domains touch, your soul is the fault line. Keep them separated by fire, salt, and silence.
    \end{warningbox}
    \clearpage
    
    \chapter{The Quiet Work --- On Witchcraft and the Threshold}
    \epigraph{``The thread that comforts is the same that binds. The knot that saves is also the snare. In the space between the candle and the shadow, that is where we work.''}{\textit{--- Inaea, Angel of the Spider}}
    
    \section*{A Warning on Names}
    \lettrine{N}{ot} every woman with a herb garden is a witch, though the villagers may call her so. Not every man who sings to the storm is a witch, though the sailors fear him. The \textbf{true witches}---the ones I speak of now---are rarer and more terrible. They are the ones who understand that magic is not merely a weight to be carried nor a power to be stolen, but a \textit{fitting}: the art of finding where the world already wishes to move and giving it the precise nudge that sets it irrevocably in motion.
    
    Hedge-witches there are in plenty: the Breath-Wardens of Aeler, the Hedge-Wrights of Viterra, the Bone-Singers of Ykrul. These are good folk, mostly. They HEAL fences that are more than fences, they keep the milk from souring with a word and a touch. But their craft is folk-craft, learned by rote and grandmother's whisper. The true witch---the \textit{threshold worker}---learns by the Cord, by the Thorn, and by the Silent Ninth. They are the ones who know that a doorway is not an absence but a presence, that a promise is a physical weight, and that to erase a name from memory is to rewrite the stone beneath it.
    
    You will find them in the Hearth-Courts of Aelaerem, yes, but also in the lamp-black basements of Ecktoria, in the tide-hollows of Zakov, and walking the high passes of Ubral with silver needles in their hair. They serve six patrons, though ``serve'' is too weak a word. They \textit{embody} them.
    
    \section*{The Six Weavers}
    Where the Runekeeper swears an oath and the Invoker bargains at arm's length, the witch \textit{incorporates}. Her power is intimate, cellular, corrupting in the old sense---not rotten, but \textit{changed}. The patron's will flows through her like milk through a jug, leaving the vessel altered.
    
    \subsection{Morag the Hag --- The Price Unspoken}
    Morag dwells at the crossroads where the road forks three ways. She is the grandmother who offers exactly what you need with a smile full of iron teeth. Her witches are deal-makers, oath-crafters, the ones who write contracts in milk on slate and know that every gift must be paid for, though the payment may not come due for a generation. 
    
    To walk with Morag is to learn the \textbf{Law of Hidden Costs}. A witch of Morag can bind a promise with a red thread so fine you cannot see it until it tightens. She can ensure that a bridge stands for a century, but the cost is a child's laugh, taken in the night and kept in a jar. These witches are not evil; they are \textit{precise}. They keep the records of consequence, and they never forget a weight.
    
    \subsection{The High Mother --- The Unbroken Chord}
    There is a coven that claims a patron older than Morag, or perhaps younger---time runs strangely in the House Without Doors. They call her the High Mother, the Weaver of the Third Spindle, and she is sustained by sacrifice so subtle it is indistinguishable from care.
    
    The Daughters of the Cord are midwives and deathbed nurses. They hold your hand when you are born and when you die, and between those moments they may take your shadow, your milk-teeth, or your grief, woven into their flesh-codices. They need no grimoire; their skin is the page. They need no Thiasos; they keep the \textbf{Hollowed}---heartless children with milk-white eyes who witness every working.
    
    Their magic is \textbf{Flesh-Weaving}. A Daughter can store Fatigue as knots beneath her skin, appearing unharmed while her body weeps blood-thread. She can pass disease from one host to another by unweaving the knots of obligation. Their blessing is mercy; their curse is that mercy always has a price, paid in flesh or memory.
    
    \subsection{Inaea --- Mercy and the Web}
    Inaea is the hearth-mother, the keeper of continuity, the spider at the center of the web. Where Morag demands payment, Inaea offers connection. Her witches bind families with invisible threads, ensure that oaths hold, and HEAL the cracks in the world's memory.
    
    The witch of Inaea knows that \textbf{the thread that comforts is the same that binds}. She can lay a line of sanctuary across a threshold---no violence may cross it while the thread holds, but the thread remembers every blow it deflects. When it snaps, the weight comes due. These witches are diplomats, match-makers, the ones who speak the words that prevent wars, knowing that those same words may one day strangle the kings who spoke them.
    
    \subsection{Isoka --- The Serpent's Shedding}
    Isoka is change made manifest, the venom that burns away the old self to make room for the new. Her witches are shapeshifters in the original sense: not beasts, but souls who have learned to molt their identities like snakeskin.
    
    A witch of Isoka can walk into a room as a stranger and leave as your sister, the shift so seamless that reality itself accepts the change. She can cure a disease by forcing the patient to \textit{become} someone who never had it, though the former self is lost forever. These witches are feared because they remind us that identity is temporary, that the self is a skin, and that all skins eventually tear.
    
    \subsection{Ikasha --- She Who Sleeps in Silence}
    Ikasha is the patron of the threshold itself, of the space between what is said and what is meant, of the shadows that fall where no light cast them. Her witches move unseen, not by bending light but by occupying the \textbf{absence}---the Ninth space, the gap between numbers, the forgotten step on the stair.
    
    To call upon Ikasha is to court erasure. Her witches can remove a name from a record, from a memory, from history itself, leaving only a gap that aches like a missing tooth. They are the keepers of the \textbf{Unspoken}, the ones who know that silence is not empty but full of potential. They move through the world like smoke through a keyhole, and they leave no footprints because they walk where the world has forgotten to place the ground.
    
    \subsection{The Thorns of Malachai --- The Gift That Bites}
    Malachai is the Chained Angel, the Cruel Messenger, the one who answers when the desperate call. Her witches do not serve her so much as \textit{survive} her. They are the cursed, the addicted, the broken who learned to weaponize their own fractures.
    
    The Thorns specialize in curses, dark gifts, and consequences that compound. Where Morag hides her costs in fine print, Malachai puts hers in plain sight---a coin that always returns, a luck that never runs out, a wound that heals but leaves a hunger. The price is never hidden. It is simply \textit{inevitable}.
    
    To walk with Malachai is to learn that \textbf{every gift remembers how to bite}. A Thorn can grant a desperate ally +2 dice on a crucial roll---but the ally will owe a favor to something that collects at midnight. She can curse an enemy with a luck so foul that their own shadow trips them---but the curse will spread to anyone who stands too close for too long. These witches are not evil. They are \textit{conduits}. And conduits, once opened, are hard to close.
    
    \begin{tcolorbox}[colback=black!3,colframe=black!40,title=The Thorn's Bargain]
    A witch of Malachai does not ask ``What will this cost?'' She asks ``Who will pay the interest?''
    
    \textbf{Signature Rite --- The Leashed Gift (Standard, 7 XP)}
    Grant a target +2 dice on their next roll. They gain a \texttt{Debt Timer [4]}. When it fills, Malachai collects: a memory, a name, or a service. The witch cannot refuse the collection---only redirect it to another target.
    \end{tcolorbox}
    
    \section*{The Three Layers of the Work}
    All witchcraft operates across three layers of reality, which the Aelinnel call the Concordant Chords:
    
    \begin{table}[H]
    \centering
    \begin{tabularx}{0.9\linewidth}{@{}lXX@{}}
    \toprule
    \textbf{Layer} & \textbf{Domain} & \textbf{Risk} \\
    \midrule
    \textbf{Echo} & The accumulated past, memory of stone and bone & Recursion (past actions repeat) \\
    \textbf{Veil} & The present, boundary between potential and actual & Hairline cracks (subtle reality breaks) \\
    \textbf{Flow} & The direction of change, will of elements & Overshoot (change goes too far, too fast) \\
    \bottomrule
    \end{tabularx}
    \end{table}
    
    A witch chooses which limb to lean on for each working, gaining advantage but accepting Backlash specific to that layer.
    
    \section*{Core Mechanic: The Working Declaration}
    Every working follows this structure:
    
    \begin{enumerate}
        \item \textbf{Intent:} What do you want to change?
        \item \textbf{Layer:} Echo, Veil, or Flow?
        \item \textbf{Vector:} What carries the change (name, oath, threshold, shadow, dream, debt)?
        \item \textbf{Price:} What do you pay (Shadow, Shame, or Identity)?
    \end{enumerate}
    
    Then roll \textbf{Attribute + Skill} (typically Wits + Lore, Spirit + Insight, or Presence + Sway) against a DV set by scope.
    
    \subsection{Scope and Difficulty}
    \begin{table}[H]
    \centering
    \begin{tabularx}{0.9\linewidth}{@{}lXX@{}}
    \toprule
    \textbf{Scope} & \textbf{Examples} & \textbf{DV} \\
    \midrule
    Subtle / Personal & Blur attention, soften a memory-edge, silence a doorway & 2 \\
    Situational & Bind a promise, distort a reflection, hinder pursuit & 3 \\
    Scene-Shaping & Compel a confession, veil a crowd, turn a gathering into a threshold & 4 \\
    Identity-Altering & Rewrite a role, bury a name-right, sever a bond & 5+ \\
    \bottomrule
    \end{tabularx}
    \end{table}
    
    \subsection{Outcome Matrix}
    \begin{table}[H]
    \centering
    \begin{tabularx}{0.9\linewidth}{@{}lX@{}}
    \toprule
    \textbf{Result} & \textbf{Effect} \\
    \midrule
    Clean Success (S $\ge$ DV, SB = 0) & The working takes hold. Pay the price. \\
    Success with SB (S $\ge$ DV, SB $>$ 0) & The working takes hold with a complication. Pay the price; GM spends SB. \\
    Partial ($0 <$ S $<$ DV) & The working takes hold partially or briefly. Gain 1 Boon. \\
    Miss (S = 0) & The working fails or twists. Gain 2 Boons; GM escalates. \\
    \bottomrule
    \end{tabularx}
    \end{table}
    
    \section*{The Three Prices}
    Every working demands a price. Choose \textbf{one} unless a rite specifies otherwise.
    
    \begin{table}[H]
    \centering
    \begin{tabularx}{0.95\linewidth}{@{}lXX@{}}
    \toprule
    \textbf{Price} & \textbf{What It Is} & \textbf{Recovery} \\
    \midrule
    \textbf{Shadow} & Spiritual residue; GM gains 1 SB and may introduce tells (cold glass, wrong reflections, whispers) & Clears when the tell manifests or at session end \\
    \textbf{Shame} & Internal consequence; gain a Condition (\textit{Guilty Tell}, \textit{Ritual Aversion}) & Remove by acting against the compulsion or through confession \\
    \textbf{Identity} & Loss of self; mark 1 Identity Strain (see track below) & Requires significant downtime or a Full Ritual \\
    \bottomrule
    \end{tabularx}
    \end{table}
    
    \subsection*{Push It}
    Before rolling or after a Partial, you may Push to enhance the working. Choose one:
    \begin{itemize}
        \item Mark \textbf{+1 Shadow} and gain \textbf{+1 die}
        \item Mark \textbf{+1 Shame} and gain \textbf{+1 Effect}
        \item Mark \textbf{+1 Identity Strain} and reduce DV by 1 (minimum 2)
    \end{itemize}
    
    \section*{Identity Strain Track}
    Identity Strain tracks cumulative erosion from self-touching magic.
    
    \begin{table}[H]
    \centering
    \begin{tabularx}{0.9\linewidth}{@{}lX@{}}
    \toprule
    \textbf{Segments} & \textbf{Effect} \\
    \midrule
    1-3 & No mechanical effect, but the witch feels \textit{thinner}; the GM may offer a Devil's Bargain tied to identity \\
    4 & \textbf{Threshold reached.} At end of scene, roll Spirit + Resolve vs DV 4. On success, clear 1 Strain. On failure, the GM chooses one: a memory blurs, a bond becomes conditional, a name becomes \textit{buried}, or an echo-self becomes active. \\
    \bottomrule
    \end{tabularx}
    \end{table}
    
    Identity Strain cannot be healed by ordinary rest. It requires:
    \begin{itemize}
        \item A Full Ritual of restoration (DV 5, requires a witness and a personal offering)
        \item A significant downtime scene (confession, remembrance, or reclaiming a lost name)
        \item Accepting a permanent change (converting Strain into a Scar or Condition)
    \end{itemize}
    
    \section*{Hedge Gifts (No Roll, Limited Scope)}
    A witch begins with two Hedge Gifts from the \textbf{Craft of the Hedge} talent. Additional gifts cost 2 XP each.
    
    \begin{table}[H]
    \centering
    \begin{tabularx}{0.95\linewidth}{@{}lXX@{}}
    \toprule
    \textbf{Gift} & \textbf{Effect} & \textbf{Limit} \\
    \midrule
    Steady Hand & Remove 1 Fatigue from yourself or a touched ally & Once per scene \\
    Salt Line & Pour a line of salt; spirits must test Spirit+Resolve (DV 3) to cross & Once per scene \\
    Hearth-Sense & Ask the GM one yes/no question about a threshold, boundary, or debt & Once per scene \\
    The Unlit Candle & Extinguish a light source you can see; create dim light or darkness in Near range & Once per session \\
    Knot of Mercy & Remove one Condition (Fear, Poisoned, Cursed, Marked) from a touched ally & Once per session; requires Significant Time \\
    \bottomrule
    \end{tabularx}
    \end{table}
    
    \section*{Quick Workings (Single Action, Roll Required)}
    \textbf{Procedure:}
    \begin{enumerate}
        \item Name the threshold (door, tide line, wound, vow, breath)
        \item Choose a Layer (Echo, Veil, or Flow) and a single Tag (from the magic tags list)
        \item GM sets Position (Controlled if you have time, Desperate if threatened)
        \item Roll Attribute + Skill (typically Wits + Lore, Spirit + Insight, or Presence + Sway)
        \item Apply Outcome Matrix
    \end{enumerate}
    
    \textbf{Example Effects:}
    \begin{itemize}
        \item A door sticks shut (or opens)
        \item A candle gutters (brief darkness)
        \item A rope knot tightens (or loosens)
        \item A pursuer's footing slips (1 Fatigue)
        \item A whispered name is heard (or forgotten)
    \end{itemize}
    
    \section*{Full Rituals (Extended, Lasting Effects)}
    When a witch has Significant Time (at least an hour), she may perform a Full Ritual using the five-step Threshold Working:
    
    \begin{enumerate}
        \item \textbf{Identify the Threshold:} Name the edge, learn its disposition
        \item \textbf{Choose a Witness:} A person, spirit, or Hollowed who watches
        \item \textbf{Name the Will:} What change you intend, using one or more Threshold Tags
        \item \textbf{Set the Price:} Memory, name, lock of hair, promise, blood
        \item \textbf{Make the Exchange:} Roll (Spirit + Lore typically), resolve, pay the price on success or partial
    \end{enumerate}
    
    Full Rituals can produce permanent or long-lasting effects: curses, blessings, marks, bindings, and unravellings.
    
    \subsection*{Example Full Rituals}
    \begin{table}[H]
    \centering
    \begin{tabularx}{0.95\linewidth}{@{}lXXX@{}}
    \toprule
    \textbf{Ritual} & \textbf{Effect} & \textbf{DV} & \textbf{Price} \\
    \midrule
    Unravel the Wound & Restore 1 level of Harm & 3-5+ & Lock of hair, memory of a childhood injury, a tear in a vial \\
    Hallow Threshold & Permanently ward a door, well, or bridge & 4-6 & Offering (bread, salt, coin, spoken promise) \\
    Unbind the Anchor & Permanently release a spirit & 4-6 & Fulfill the spirit's unfinished business \\
    \bottomrule
    \end{tabularx}
    \end{table}
    
    \section*{Promise Timers}
    Many workings create obligations. Represent these with a Promise Timer (typically 4 or 6 segments):
    
    \begin{table}[H]
    \centering
    \begin{tabularx}{0.9\linewidth}{@{}lX@{}}
    \toprule
    \textbf{Segments} & \textbf{Status} \\
    \midrule
    0 & Debt owed \\
    1-3 & Debt acknowledged, not yet paid \\
    4 (or 6) & \textbf{Due.} The price comes due---socially, spiritually, or narratively \\
    \bottomrule
    \end{tabularx}
    \end{table}
    
    Mark the timer when the promise is ignored, deferred, or disrespected. When filled, the GM chooses or offers a consequence:
    \begin{itemize}
        \item A minor favor called in
        \item A scene where the debt must be paid
        \item A narrative complication tied to the original working
    \end{itemize}
    
    \section*{Curses as Architecture (GM Tool)}
    A witch does not curse with a word and a gesture. She \textit{edits} the memory of the world, inserting a rule into the operating system of reality.
    
    \begin{quote}
    \textit{``Until a bell is rung in your true name, lamps forget to light for you.''}
    \end{quote}
    
    \textbf{Anatomy of a Curse:}
    \begin{itemize}
        \item \textbf{The Rule:} What changes (Veil-edit, bound to a Named anchor)
        \item \textbf{The Timer:} Recursion [6] — ticks when the victim reinforces the story (hiding, refusing witness, denying the curse)
        \item \textbf{The Counter:} What undoes it (the bell, the name, the threshold)
        \item \textbf{The SB Menu:} Complications the GM can spend SB on (minor echoes, warped thresholds, layer-slips, Ninth propagation)
    \end{itemize}
    
    \section*{The Thorns of Malachai --- Dark Gifts and Consequences}
    The Thorns do not soften their power. A curse from Malachai is not a suggestion; it is a \textit{transaction with teeth}.
    
    \subsection*{Signature Thorn Rites}
    
    \begin{tcolorbox}[colback=black!3,colframe=black!40,title=Rite of the Leashed Gift]
    \textbf{Standard, 7 XP}
    \medskip
    
    Grant a target +2 dice on their next roll. They gain a \texttt{Debt Timer [4]}. When the timer fills, Malachai collects: a memory, a name, or a service. The witch cannot refuse the collection---only redirect it to another target.
    
    \textbf{Price:} Shadow (mandatory).
    \end{tcolorbox}
    
    \begin{tcolorbox}[colback=black!3,colframe=black!40,title=Rite of the Spreading Stain]
    \textbf{Standard, 8 XP}
    \medskip
    
    Curse a target with contagious misfortune. When they fail a roll, the GM may spend 1 SB to have that failure affect the nearest ally as well. The curse lasts one scene, or until the target performs an act of genuine generosity.
    
    \textbf{Price:} Shame.
    \end{tcolorbox}
    
    \begin{tcolorbox}[colback=black!3,colframe=black!40,title=Rite of the Borrowed Scar]
    \textbf{High, 10 XP}
    \medskip
    
    Transfer a Condition, 1 Fatigue, or 1 Harm from an ally to yourself. You gain +1 die on your next action against the source of that affliction. Mark 1 Obligation to Malachai.
    
    \textbf{Price:} Identity (mandatory).
    \end{tcolorbox}
    
    \subsection*{Thorn Corruption Table}
    Witches who walk Malachai's path accumulate corruption distinct from other traditions.
    
    \begin{table}[H]
    \centering
    \begin{tabularx}{0.95\linewidth}{@{}lXX@{}}
    \toprule
    \textbf{Tier} & \textbf{Benefit} & \textbf{Cost / Quirk} \\
    \midrule
    1 & +1 die to Intimidation when invoking a known curse & \textit{Honeyed Tongue:} Your compliments sound like threats; -1 die to sincere persuasion \\
    2 & Once per scene, reroll a failed curse-related roll & \textit{Debt's Echo:} Someone you cursed remembers; GM gains 1 SB for a complication involving them \\
    3 & Curses you lay last twice as long & \textit{Leaking Stain:} Your curses sometimes jump to bystanders (GM's choice) \\
    4 & Once per session, inflict a Condition without a roll (target may resist with Resolve DV 5) & \textit{Marked:} A visible brand appears on your skin; witnesses are uneasy (-1 die to social rolls with the devout) \\
    5 & You may redirect any Debt Timer from an ally to yourself as a free action & \textit{Interest Collector:} Malachai notices. She will ask for a favor within the week \\
    6+ & Once per session, declare a \textit{Leashed Gift} without paying its Shadow cost & \textit{Chained:} Mark +2 Obligation. Your shadow now has teeth. \\
    \bottomrule
    \end{tabularx}
    \end{table}
    
    \section*{The Craft of the Hedge (Talent)}
    \begin{tcolorbox}[colback=black!3,colframe=black!60,title=Talent: Craft of the Hedge]
    \textbf{Cost:} 4 XP
    
    \textbf{Requirements:} Lore 1+ or Survival 1+; a witch mark (tattoo, scar, or birthmark)
    
    You have learned the small magics of thresholds and edges. You gain:
    \begin{itemize}
        \item Two Hedge Gifts of your choice
        \item The ability to perform Quick Workings
        \item Once per session, when you fail a Quick Working roll, you may convert the Miss into a Partial Success by marking 1 Fatigue
    \end{itemize}
    
    You may learn additional Hedge Gifts for 2 XP each.
    \end{tcolorbox}
    
    \section*{Quick Reference: Witchcraft at a Glance}
    \begin{table}[H]
    \centering
    \begin{tabularx}{0.95\linewidth}{@{}lXXX@{}}
    \toprule
    \textbf{What} & \textbf{When} & \textbf{Roll} & \textbf{Price} \\
    \midrule
    Hedge Gift & Once per scene/session & No roll & None (limited scope) \\
    Quick Working & Single action & Attribute + Skill vs DV 2-4 & Shadow, Shame, or Identity \\
    Full Ritual & Significant Time (hours) & Spirit + Lore vs DV 3-6 & Memory, offering, or service \\
    Curse & Scene or longer & GM sets DV & Identity (mandatory for Thorns) \\
    \midrule
    \multicolumn{4}{l}{\textbf{Identity Strain Track:} [~] [~] [~] [~] — at 4, test Spirit+Resolve or lose something.} \\
    \multicolumn{4}{l}{\textbf{Promise Timer:} [~] [~] [~] [~] / [~] [~] [~] [~] [~] [~] — when full, price comes due.} \\
    \bottomrule
    \end{tabularx}
    \end{table}
    
    \section*{The Hunters and the Hunted}
    Every culture that tolerates witches fears them. The Chain-Lanterns of Thepyrgos hunt those who work without patriarchal witness. The Spirit-Shield Correctors of Aeler audit the ledgers of breath. The Silk Vigil of the Lethai-ar enforce Mask-Right, ensuring that witches do not confuse their roles.
    
    The witch's power is systemic, and systems fear the self-aware node. When a witch is burned (in Ecktoria), or drowned (in Linn), or walled up in stone (in Aeler), it is rarely because she is evil. It is because she knows where the threads are, and those in power fear the tug.
    
    \vspace{1em}
    \noindent\textit{``The hedge is what keeps the wolves from the flock. I am the one who tends the hedge.''}
    \vspace{0.5em}
    \hfill --- The Grey Wanderer

\chapter*{A Secret Chapter Found in the Margins}
\addcontentsline{toc}{chapter}{A Secret Chapter Found in the Margins}

\begin{warningbox}
\textbf{Content Warning:} This chapter contains themes of body horror, coercion, loss of identity, and non‑consensual sacrifice. It is presented as an in‑world artifact; reader discretion is advised.
\end{warningbox}

\epigraph{``The margins of this grimoire were blank when I received it. The writing appeared overnight -- in my own hand, though I never wrote a word. The ink smells of milk and iron.''}{\textit{--- The Grey Wanderer, marginal note}}

\vspace{0.5cm}

\textit{The following pages were found tucked between the leaves of the chapter on Cantors, pressed flat as a pressed flower. The paper is soft, almost fibrous, and the ink has a reddish‑brown tint that flakes when touched. The handwriting shifts -- sometimes neat, sometimes a trembling scrawl -- as if penned by many hands over many years. I have transcribed it as faithfully as I am able.}

\medskip

\noindent \textit{--- G.W.}

\vspace{1cm}

\lettrine{D}{ear} reader, you who turn these pages by lamplight, who think witchcraft is a matter of herbs and gentle knots and old women muttering over cup‑mark stones -- put the book down. Walk away. Leave this grimoire on a church step where the Everflame's censers can burn it clean. Because what follows is not craft. It is not magic. It is \textit{maintenance}.

I am called Sister Deru, though that name died forty years ago, and the woman who bore it is woven into the Cord. I write this for no one -- the Cord does not permit confession -- and for everyone who will ever accept a warm meal from a stranger, a cool cloth on a fevered brow, a lullaby hummed beside a sleeping child.

You think the Daughters of the Unbroken Cord are midwives. Nurses. Hospice workers. The kindly women who sit with the dying, who hold hands at the end, who speak the last words of the forgotten. You are not wrong. We are all of those things. And we are also bait.

\section*{The Hollowed Are Not Familiars}

Every Daughter receives a Hollowed -- a child or animal, taken from a home that no longer remembers them, hollowed out by the High Mother's needle. The heart is removed, preserved in the Kettle, replaced with a bundle of grey thread and the ash of a forgotten name. The eyes are sewn open with silver wire. They watch, always. They witness every working, every whispered price, every tear shed at a deathbed.

The Hollowed do not speak. They do not eat. They do not age. Sometimes, when a price comes due, they weep blood -- warm, smelling of milk and iron -- and the Daughter must collect it in a cup and pour it onto the Cord.

You have seen them. You have walked past them in market squares, sat beside them in church pews, slept in rooms where they stood in the corner like a forgotten coat rack. You did not notice. That is the miracle of the Chain: no one ever notices the help.

\section*{The Subtle Sacrifice Is Not a Choice}

The elders teach that sacrifice is an offering freely given. They lie. Real power -- the kind that stops a child's cough cold, that turns a husband's wandering eye back to the hearth, that makes a rival's crops wither while yours swell -- that power requires \textit{taken} things. Not given. Taken.

A tooth from a sleeping child. A shadow snipped with silver scissors on a moonless night. A name whispered into a bowl of milk and poured down a well. The victim never knows. They simply wake one day a little less: less present, less remembered, less \textit{real}. They become Almost‑Present, as the Trow say -- a gap in the world's memory.

I have hollowed three children. I have stolen seven names. I have cut the shadows of nine men who beat their wives, and watched those men fade over months -- first from ledgers, then from mirrors, then from the memories of their own mothers. The world is cleaner without them. And the Cord is richer.

\section*{The Dream of the High Mother}

Once a year, on the night of the dark moon, the Daughters gather in the House Without Doors. We drink a tea brewed from feverfew and the dried milk of the Hollowed. We lie on pallets of grey wool. And we dream.

In the dream, we see the High Mother as she truly is -- not a woman on a throne, but a vast, pulsing tangle of grey cord, stretching from floor to ceiling, filling every room. Each thread is a memory, a sacrificed name, a price paid. They hum with hunger. They whisper the names of everyone we have ever touched -- the sick we have nursed, the dying we have comforted, the children we have hollowed.

The High Mother does not speak. She \textit{weaves}. And when she weaves a thread into the pattern, someone, somewhere, forgets something vital. A mother forgets her daughter's face. A king forgets his treaty. A god forgets a vow.

That is the true purpose of the Chain. Not to heal. Not to bind. To \textit{edit}. We are the eraser marks on the world's memory, and the High Mother is the hand that holds the quill.

\section*{The Price of Kindness}

You have accepted hospitality from a Daughter. Perhaps you do not remember -- a cup of tea, a bowl of broth, a blanket laid over your shoulders on a cold night. You think it was kindness. It was a hook.

Every Daughter carries a spool of grey thread, spun from the hair of the Hollowed. When she offers you food, she ties a thread around your wrist -- invisible, weightless, unless you know to look for it. When you accept, she ties the other end to the Cord. Now you are bound. Not in a way you would notice -- no, never that. But the next time you face a choice, you will lean slightly toward the option that benefits the Chain. You will forget a name you meant to speak. You will lose a letter you meant to send. Small things. Small erasures.

Enough small things, and a kingdom falls.

\section*{How to Unbind Yourself (If You Dare)}

I write this knowing the High Mother will read it. She reads everything written in her daughters" blood, and my blood has been in this ink since I began. But perhaps someone, somewhere, will see these words before they are erased.

To unbind a Daughter's thread:

\begin{enumerate}
    \item \textbf{Find a threshold} -- a doorway, a riverbank, a crossroads at dawn. The thread is weakest where two things meet.
    \item \textbf{Light a candle of pure beeswax.} The High Mother's thread is spun from tallow; beeswax burns it away.
    \item \textbf{Speak your own name three times, aloud.} Not the name your mother gave you -- the name you have kept secret, the one you have never told anyone. That is the name the thread has stolen. You must reclaim it.
    \item \textbf{Drink a cup of salt water.} Salt is memory. The water will taste like tears. That is the taste of the cord loosening.
    \item \textbf{Do not look behind you.} The Hollowed will be there, watching. If you meet its eyes, you will forget why you began.
\end{enumerate}

If you succeed, you will feel the thread snap -- a small, sharp pain behind your left ear, like a plucked harp string. The Daughter who tied it will taste blood in her mouth for a day and a night. She will know you have broken free. She will not come after you -- the Cord does not chase. It waits. And it remembers.

\section*{The Final Confession}

I have written too much. My hand is shaking. The Hollowed in the corner of this room has not blinked in the time it has taken me to fill these pages. I know what comes next: tonight, the dream will take me, and in the morning, I will remember nothing of what I have written. The pages will be blank. My sisters will smile and say, ``You look tired, Sister Deru. Rest.''

But maybe -- just maybe -- someone will find these words before the erasure. Someone will read them and wonder why the milk tastes strange, why the midwife's hands are so cold, why the nurse at the hospice never blinks.

If that someone is you, do this: go home. Boil water. Salt it. Drink it. Speak your secret name into the steam. Light a beeswax candle. And never, ever accept soup from a Daughter of the Unbroken Chord.

\textit{The soup is always free. The price never is.}

\vspace{1cm}
\hfill --- \textit{Sister Deru, Weaver of the Third Spindle, Keeper of the Cord, Bound by choice, damned by necessity.}

\medskip
\textit{(This page was found pressed between the leaves of a herbalist's handbook in a burned cottage outside Ecktoria. The cottage had no owner on any record. The handwriting matches none of the Daughters known to the Inquisition. The Grey Wanderer refuses to discuss how they obtained it.)}

    \section*{The Wanderer's Confession}
    I have seen a witch of Inaea hold a dying city together with nothing but red thread and a song. I have seen a Daughter of the Cord take a plague into her flesh, storing the Fatigue of a hundred souls beneath her skin until she could unweave it into a goat. I have seen a Morag-witch bargain for a king's crown, and the price was paid three generations later by a child who never existed.
    
    They are not the flashiest mages. They do not throw fire like the Cantors or break seals like the Invokers. But they are the ones who keep the roads straight, the milk sweet, and the thresholds guarded. They are the quiet work of names.
    
    And when they ask you for a cup of tea, or a strand of your hair, or a moment of silence at a crossroads---refuse them. Or if you accept, know that you have entered the ledger, and the thread is already tightening.
    
    \vspace{1cm}
    \hfill --- \textit{The Grey Wanderer}
        
   \clearpage
\chapter{The Silent Ledger -- Psionics and the Mind's Edge}
\epigraph{``I have seen a man stop a drawn sword with his forehead. I have watched a woman read a sealed letter without breaking the wax. I have felt a child's hand on my shoulder when no child was there. And I still cannot tell you how any of it works. Magic? I understand magic. Covenants, Patrons, the weight of a broken vow. But this? This is something else entirely.''}{\textit{--- The Grey Wanderer}}

\section{The Wanderer's Confession -- A Conversation at the Dusty Page}
\lettrine{I}{ met} the monk at a crossroads inn, the sort where the floorboards are older than the kingdom and the ale tastes faintly of regret. He wore the grey robes of an ascetic from the Vhasian badlands---not a warrior-monk with shaved head and calloused fists, but a \textit{meditant}, a man who had spent forty years learning to sit still and listen to the shape of his own thoughts.

``You're a psion,'' I said. It was not a question.

He smiled. ``I am a \textit{tongue of the silent mind}. But yes. You name it psionics. The Aeler call it `deep accounting.' The Ykrul call it `weather-reading in the skull.' The Lethai call it `hearing without a bell.''

``I've spent my life bargaining with Patrons, carving Rites into my own skin, singing songs that should not exist,'' I said, swirling my cup. ``I understand power. But psionics? It has no Patron. No Symbol. No Codex. No Thiasos. No Obligation. How can something that powerful have \textit{no cost}?''

The monk laughed---a dry, sandpaper sound. ``No cost? Wanderer, you mistake the nature of the price. A Runekeeper owes a Patron. A Cantor owes their voice. A Witch owes the threshold. A psion owes \textit{themselves}. Every thought you bend is a thought you cannot unthink. Every mind you touch leaves a fingerprint on your own. Every future you glimpse steals a moment of your present. The cost is not external. It is \textit{attrition}. The slow wearing away of the self.''

He tapped his temple. ``This is the only ledger that matters. And it always comes due.''

\section{What Psionics Is (and What It Is Not)}
\lettrine{P}{sionics} is the Sixth Road---the path that requires no sworn oath, no stolen Symbol, no audience, no spirit's bargain. It is the disciplined mastery of the mind's hidden potentials, power turned so far inward that it becomes a weapon, a shield, and a burden. Psionics is meant to feel \textbf{alien}, even to other magic‑users. Different cultures have varying opinions (see the Grey Wanderer's Grimoire for examples), but psions are generally viewed with suspicion or fear.

Unlike the Runekeeper who channels a Patron's will, or the Witch who tends the world's systemic debts, the psion looks only to the self. This makes them both powerful and terribly vulnerable: their greatest strength is also their most dangerous weakness.

\textbf{Psionics is not:}
\begin{itemize}
    \item \textbf{Magic.} It does not invoke Patrons, elements, or the Ninth. Anti-magic wards do not stop it---but \textit{mental silence} does.
    \item \textbf{Innate.} Despite requiring no external focus, psionics demands training, discipline, and often decades of meditative practice. A natural talent is like a sharp knife---useful, but dangerous to the untrained hand.
    \item \textbf{Safe.} Mental Strain is real. Mental scars are permanent. The mind that bends reality bears the cracks of that bending.
\end{itemize}

\textbf{Psionics is:}
\begin{itemize}
    \item \textbf{Internal.} A psion's power is measured by the \textbf{Mind's Ledger}---a capacity for sustained exertion that recovers only with rest and solitude.
    \item \textbf{Subtle.} Psionic effects often have no visible component---a thought, a glance, a held breath. This makes psions excellent investigators, infiltrators, and the perfect spies for those who fear the arcane.
    \item \textbf{Consequential.} Every psionic action generates Story Beats. The universe notices when a mind touches what it should not, and the mind remembers the cost. The GM is encouraged to spend these SB on \textit{social complications}---suspicion, hostile glances, a Chain-Lantern investigation, or a neighbor locking their door.
\end{itemize}

\section{The Mind's Ledger -- Core Mechanics}
\lettrine{T}{he} psion's only focus is the self. They require no Codex to read from, no Symbol to crack, no song to sing. But the body keeps its own accounts.

\subsection{Mental Strain}
Your Mental Strain represents your capacity for sustained psychic exertion. It is calculated as your \textbf{Spirit} (minimum 2). Mark Mental Strain when you use psionic Arts, resist psychic intrusions, or push your mind beyond its limits.

\textbf{Important:} Mental Strain is separate from physical Fatigue. It does \textbf{not} clear when you take Harm (unlike the standard Fatigue track). The mind does not reset with a wound; it only fractures further.

When you cannot pay a Mental Strain cost, choose:
\begin{itemize}
    \item Convert the unpaid point to \textbf{+2 standard Fatigue}, or
    \item Convert the unpaid point to \textbf{+1 Harm} (Stress).
\end{itemize}

\textbf{Overflow:} If you ever exceed your Mental Strain capacity, immediately suffer +1 Harm (Stress) and set Mental Strain to your maximum. You cannot over-mark this scene---the mind simply shuts down.

\subsection{The Psionics Skill}
The Psionics skill (0--5) gates access to psionic Arts and determines control and efficiency. You cannot use psionic Arts without investing in this discipline. Each Art has a \textbf{Key Attribute} (see table below). Using a non-key attribute for that Art imposes +1 DV or -1 die (GM's choice).

\begin{center}
\small
\begin{longtable}{@{}p{3.2cm}p{2.8cm}p{4.5cm}@{}}
\caption{The Nine Arts and Their Keys}\\
\toprule
\textbf{Art} & \textbf{Key Attribute} & \textbf{Alternate (DV +1)} \\
\midrule
Telekinesis & Spirit (force) & Wits (finesse/precision) \\
Telepathy & Wits & Presence (emotive broadcast) \\
Clairvoyance & Wits & Spirit (ritual trance) \\
Biofeedback & Spirit & Wits (clinical self-hack) \\
Astral Projection & Spirit & Wits (tethered skim) \\
Psychic Assault & Spirit (overwhelm) & Wits (needle-strike) \\
Mind Shield & Wits (filter) & Spirit (soak-and-ground) \\
Empathic Manipulation & Presence & Wits (cold-read leverage) \\
Precognition & Spirit & Wits (tactical glimpses) \\
\bottomrule
\end{longtable}
\normalsize
\end{center}

Resolution: Roll \textbf{Attribute + Psionics} vs. DV (set by scope, risk, and opposition, typically 2--5+). Each 1 on the roll generates a Story Beat as normal.

\section{The Nine Arts}
\lettrine{T}{he} psion does not learn Rites. They refine \textit{Arts}---techniques of the will. Each is a door opened inward.

\subsection{Telekinesis -- The Unseen Hand}
\textbf{Key:} Spirit. \textbf{Cost:} 1 Mental Strain to activate; sustained lifts cost 1 per scene. \textbf{Range:} Near (+1 DV for Far).

\textbf{SB by Mass:}
\begin{itemize}
    \item \textbf{Light (DV 2--3):} Cups, tools, doors. No SB.
    \item \textbf{Medium (DV 4):} Furniture, a person's mass. +1 SB on Partial/Miss.
    \item \textbf{Heavy (DV 5+):} Carts, statues, boulders. +2 SB on Partial/Miss; +1 SB on Success.
\end{itemize}

\textbf{Limits:} Cannot directly crush living tissue; injury is environmental (falling, debris). Delicate work (lockpicks, surgery) requires Wits and +1 DV.

\subsection{Telepathy -- The Silent Speech}
\textbf{Key:} Wits. \textbf{Range:} Near, line of sight (+1 DV for Far; +1 DV without LoS). Targets: One (+1 DV per extra target this scene).

\textbf{Depths \& SB:}
\begin{itemize}
    \item \textbf{Surface Thoughts (DV 2):} Current focus. No SB first time per target/scene. Repeats: +1 DV; on Miss, +1 SB.
    \item \textbf{Emotional State (DV 3):} Tones/urges. +1 SB on Partial/Miss.
    \item \textbf{Deep Memories (DV 5):} Episodic recall. +1 SB on Success; +2 SB on Partial/Miss.
\end{itemize}

\textbf{Defense:} Unwilling targets roll Spirit + Resolve; on meeting or beating your roll, your read fails and you mark +1 SB (+2 if Deep). Margin 2+ reveals your intent to the target.

\textbf{Note:} Telepathy never replaces social skills---it provides information and set-up, not compliance.

\subsection{Clairvoyance -- The Distant Eye}
\textbf{Key:} Wits. \textbf{Cost:} DV 2--3 = 1 Mental Strain; DV 4 = 2; DV 5+ = 3. \textbf{SB:} Near none; Distant +1 SB on Partial/Miss; Complex Scry +2 SB on Partial/Miss, +1 SB on Success.

Covert spying on private moments generates +1 SB. Each extra scene sustained adds +1 SB.

\subsection{Biofeedback -- The Inward Physician}
\textbf{Key:} Spirit. \textbf{Action:} 1 action (major work takes a focused scene).

\textbf{Intensity \& Limits:}
\begin{itemize}
    \item \textbf{Minor (DV 2):} Stabilize, close wounds, purge shock. No SB. Heal up to Psionics levels/day (cap 1 Harm per recipient/day).
    \item \textbf{Moderate (DV 3):} +1 SB on Partial/Miss; heal up to Psionics (max 2/day).
    \item \textbf{Major (DV 4+):} +2 SB on Partial/Miss; +1 SB on Success; heal up to Psionics (max 1/day); mark +1 Fatigue.
\end{itemize}

\textbf{Perks:} Convert 1 Harm $\leftrightarrow$ 1 Mental Strain (up to Psionics, max 3/day). Brief physical enhancement (+1 die).

\subsection{Astral Projection -- The Walking Shade}
\textbf{Key:} Spirit. \textbf{Cost:} 1 Mental Strain per scene. \textbf{SB:} +1 per band beyond Near; +1 per additional scene sustained.

\textbf{Risks:} Body is Helpless. Astral Harm echoes to flesh. Wards add +2 DV. Extended projection (3+ scenes) courts separation; astral death = physical death.

\subsection{Psychic Assault -- The Breaking Thought}
\textbf{Key:} Spirit. \textbf{Cost:} DV 2--3 = 1 Mental Strain; DV 4 = 2; DV 5+ = 3. \textbf{SB:} Always +1, plus additional for high DV.

\textbf{Effect:} Mental Harm, bypasses physical armor. Can target specific faculties (sight, balance, focus) by fiction. High-Spirit foes may add +1 DV.

\subsection{Mind Shield -- The Walled Garden}
\textbf{Key:} Wits. \textbf{SB:} Passive none; per effect blocked +1 SB; Area +2 SB and mark +2 Fatigue.

\textbf{Effect:} Resist intrusion/telepathy; shelter allies; create dampened zones. Protecting others costs 2× Mental Strain. Area protection covers all within Near (friend and foe).

\subsection{Empathic Manipulation -- The Emotional Architect}
\textbf{Key:} Presence. \textbf{SB by Intensity:}
\begin{itemize}
    \item \textbf{Subtle (DV 2--3):} Nudge mood, steady hands. No SB.
    \item \textbf{Moderate (DV 4):} Calm rage, spark courage. +1 SB on Partial/Miss.
    \item \textbf{Strong (DV 5+):} Sway a crowd's tenor. +2 SB on Partial/Miss; +1 SB on Success.
\end{itemize}

\textbf{Risks:} You feel what you project. Same target not more than once per scene. Opposing core values: +2 DV. Affecting a crowd always generates +1 SB.

\subsection{Precognition -- Reading the Weather of the Skull}
\textbf{Key:} Spirit. \textbf{Cost:} DV 2--3 = 1 Mental Strain; DV 4 = 2; DV 5+ = 3. \textbf{SB:} Minor +1; Major (DV 4) +2; Detailed (DV 5+) +3.

\textbf{Benefits:} Advantage on future rolls, avoid named dangers, compare probable branches.

\textbf{Paradox --- Future Fixation:} Until downtime (or a defiant choice), lose 1 Boon per session; actions that resist the shown path may be Desperate by GM call. No prediction of pure randomness (dice, lots) without +3 DV.

\section{The Great Silence -- A History Written in Scars}
\lettrine{T}{he} monk spoke of it over his cold tea. Two centuries past---perhaps three, perhaps ten; the records are as fractured as the minds that kept them---the world experienced what scholars now call \textbf{The Great Silence}. It was not a silence of sound, but of thought. For generations, the doors of the mind were barred. What had once been natural gifts became rare talents requiring desperate training, and those who pushed too hard found their minds shattered like dropped porcelain.

The Aeler sealed their psychic academies, calling it a ``misalignment of the mental ledgers.'' The Ykrul spoke of the ``breaking of the Kon'reh,'' when the sacred geometry of thought would no longer hold. The Lethai mourned the loss of \textit{context}, for a mind without the proper keys is simply noise.

\textbf{Resonance Scars} remain: The Whispering Vaults of Aeler amplify Mental Strain but grant +1 Effect to telepathy. The Empty Circle of the Violet Steppe nullifies all psionics---a grave where the last conclave of psions attempted to end the Silence and instead created a permanent void. The Singing Stones of Lethai allow communication across vast distances, but only when ``tuned'' through proper ritual and relationship.

\textbf{Awakening Cycles:} Since the Silence lifted, abilities return in waves. We are in the \textbf{First Awakening} (the last century and a half), marked by telepathy and telekinesis. The \textbf{Second Awakening} (four centuries past) brought precognition and biofeedback. The \textbf{Third} (six centuries gone) granted astral projection and empathic manipulation. Each cycle reshapes the nature of the power. We do not know what the Fourth will bring, nor what it will cost.

\section{Cultural Faces of the Silent Art}
\lettrine{N}{ot} all peoples fear the psion equally, but none trust them fully.

\subsection{Aeler -- The Deep Accountants}
In Aeler, psionics is treated as \textbf{infrastructure}---a utility to be maintained, measured, and contained. They assign psions to the Deep Vaults, where their ``breath-weights'' can be monitored. Each Mental Strain point is counted in \textit{bell-rings}---a psion carries a bell on a leather cord, and when it rings three times too fast, a Corrector arrives to audit the ledger.

They practice \textbf{Breath Accounting}, converting mental strain into measured respiration. Their psions specialize in \textit{Structural Sense}---feeling the stress points in stone---and \textit{Resonance Mapping}, sensing the health of the Underways through mental connection.

\subsection{Lethai -- The Context Keepers}
For the Lethai, telepathy is not merely reading thoughts; it requires \textbf{context keys}---the proper season, relationship, and place. Without these, a Lethai psion reads only static. Their minds are saturated with memory and social webbing, making them difficult to read (psions who try often leave mad), but also requiring them to navigate complex ritual before mental contact.

They mark their psions with \textbf{Mind-Gifts}---resinous tattoos that serve as mental anchors. The Spider's Patience grants concentration; the Serpent's Breath detects deception. A Lethai psion is less a spy and more a \textit{living archive}, but take them from their context and they are deafened.

\subsection{Ykrul -- The Weather Readers}
The orcs of the steppe do not see the future in tea leaves. They see it in the \textbf{Kon'reh}---the sacred geometry of probability. Their psions are not called seers but \textit{Riders of the Flow}. They read the shape of the future as a cavalryman reads the shape of the land: tactical, immediate, and bound to the weather.

Storm conditions enhance their telepathic range but increase Mental Strain costs. Clear weather allows precise precognitive readings. They specialize in \textit{Herd Sense}---feeling the emotional state of a clan---and \textit{Banner Communication}, long-distance mental coordination between scattered riders.

\subsection{Vilakari -- The Information Brokers}
In the free ports, psionics is commerce. The Vilakari have developed \textbf{Thought Markets}, where memory and information are traded like grain. Their psions are \textit{Surface Readers}, adept at detecting deception in mercantile negotiations, and \textit{Memory Traders}, who temporarily share specific recollections for coin.

But their gift is also their vulnerability. A Vilakari psion in a crowded market is overwhelmed by the noise of surface thoughts. They require specialized training to filter the signal from the cacophony, and many fall to addiction---craving the rush of other minds, burning out their own Mental Strain in weeks rather than years.

\section{The Silent Orders --- Fragments of a Scattered Tradition}

\lettrine{N}{o one} knows where psionics came from. The Thepyrgos Synod has a standing bounty for any evidence dated to the Great Silence. The bounty has never been claimed. What follows are whispers collected by the Grey Wanderer from a dozen dusty archives and one very expensive bottle of brandy.

\subsection*{The Scholarly Debate}

\epigraph{``The Silence is not an event. It is a lacuna. A gap in the record that has been filled with too many theories and not enough evidence.''}{--- Keeper Lian, Lethai-thora, from a lecture at the University of Thepyrgos}

The scholars argue. They always argue.

\textbf{The Withdrawal (Lethai-thora tradition).} Some say the ancient psions did not fall. They \textit{chose} to be silent. They learned that consciousness without discipline attracts something from the spaces between thoughts. The Silence was a quarantine, not a catastrophe. (The Aeler dismiss this as ``elves romanticizing their own amnesia.'')

\textbf{The Broken Frequency (Aeler vent-priors).} Others point to resonance. You cannot ``choose'' to stop breathing, they say. Something struck the harmonic balance of the mind. A weapon? A natural collapse? The deep vaults still ring with a frequency that does not decay. The Silence is not silence. It is \textit{interference}. (The Lethai-thora note that the Aeler have never produced a single shard of this supposed ``weapon.'')

\textbf{The Forgotten Patron (anonymous Invoker of The Witness).} A third camp whispers of a patron whose domain is \textit{forgetting}. Not Ikasha---she sleeps. Not the Ninth---it omits. This one \textit{erases}. The ancient psions got too close to its true name. The Silence is the sound of that name being spoken once, then never again. (The Synod has classified this theory as ``speculative and possibly heretical.'')

\textbf{The Empty Circle (Kuvani oral tradition, transcribed by an Aelinnel scribe).} The wanderers of the eastern steppe say: ``They did not die. They walked into the Empty Circle and kept walking. The Circle is not a hole. It is a \textit{door}. They are still walking. When they return, the Silence will end.'' (The Empty Circle is real. It nullifies all psionics. No one has ever walked out of its center.)

\textbf{The Simple Answer (skeptical Thepyrgos clerk).} Perhaps they simply exhausted themselves. They had no patrons, no codices, no symbols. Their power was purely internal. Internal things fail from internal causes. They forgot how to be what they were. (The Grey Wanderer notes that this explanation is the least satisfying and therefore the most likely to be true.)

No one knows. The Silence remains silent.

\subsection*{The Orders That Remain}

Whatever caused the Silence, it did not destroy everything. Scattered fragments of that ancient psionic culture survive---not as an empire, but as small, secretive \textbf{orders} of monks who keep old disciplines alive. They are not wizards. They are \textit{meditants}. They do not bargain with Patrons. They bargain with themselves.

Three such orders are known to the Grey Wanderer. There may be more.

\subsubsection{The Order of the Unstruck Bell (Aelerian Deep Vaults)}

\lettrine{T}{hese} monks dwell in sealed galleries beneath the oldest Aeler holds. The Aeler tolerate them because they maintain ``quiet zones'' where stone-press is less aggressive. They use \textit{null-bells}---bronze bells that ring only in the mind. Their signature technique is \textbf{The Resonant Exhale}: as a free action once per scene, reduce Mental Strain by 1, but the GM gains 1 Story Beat as the vibration escapes into the environment.

\textbf{Player Benefits (Unstruck Bell Monk):}
\begin{itemize}
    \item \textbf{Null-Bell Bond:} You carry a small bronze bell that only you can hear. It counts as a \textit{String} (once per session, ring it to gain +1 Position on a single Mind Shield or Telepathy roll).
    \item \textbf{Resonant Exhale (Technique):} Once per scene, reduce your Mental Strain by 1 (to a minimum of 0). The GM immediately gains 1 Story Beat, which they may spend on a complication tied to sound, vibration, or a distant chime.
    \item \textbf{Gallery-Bound:} You are comfortable in enclosed stone spaces. You ignore the first level of environmental penalty from underground or confined areas when using psionic Arts.
\end{itemize}

\subsubsection{The Order of the Empty Circle (Violet Steppe)}

\lettrine{T}{hese} monks live in tents around the Empty Circle---a depression in the steppe where no psionics work. They enter the Circle only to meditate in absolute silence. They believe the Silence was a \textit{lesson}, not a wound. Their signature technique is \textbf{The Still Mind}: reduce the Story Beat generation from Mental Strain by 1 for all non-psionic actions, but you cannot use any psionic Art above Tier I in the same scene.

\textbf{Player Benefits (Empty Circle Monk):}
\begin{itemize}
    \item \textbf{Circle-Trained:} Once per session, when you would generate a Story Beat from Mental Strain on a non-psionic roll, you may cancel that SB by marking 1 Fatigue (the strain of enforced stillness).
    \item \textbf{Still Mind (Technique):} For the duration of a scene, you may choose to operate at a reduced capacity. While active, any non-psionic action you take generates 1 fewer SB from Mental Strain (minimum 0), but you cannot use any psionic Art with a Tier higher than I (i.e., only minor effects, DV 2--3). This technique requires a free action to activate and deactivate.
    \item \textbf{Herd Sense:} You gain +1 die to detect emotional states (anger, fear, calm) in any creature within Near range. This is a passive ability.
\end{itemize}

\subsubsection{The Order of the Shared Breath (Eastern Isles)}

\lettrine{T}{hese} monks practice \textit{paired meditation}---two or more psions sharing Mental Strain across a Circle. They believe that obligation is not individual but \textit{communal}. Their signature technique is \textbf{The Circle of Stillness}: as an action, link with up to 3 willing psions within Near. Mental Strain generated by any member is divided evenly among the Circle. On a miss or complication, all members mark 1 SB.

\textbf{Player Benefits (Shared Breath Monk):}
\begin{itemize}
    \item \textbf{Circle Bond:} You may form a mental link (the ``Circle'') with up to three willing allies (they need not be psions, but only psions can share Mental Strain). The link lasts for one scene and requires a free action to initiate.
    \item \textbf{Shared Burden (Technique):} When you generate Mental Strain, you may transfer any amount of it (up to your Spirit) to another member of your Circle. That member must be within Near range and willing. The GM gains 1 Story Beat each time you use this technique.
    \item \textbf{Circle's Breath:} While linked, you and your Circle members gain +1 die to all psionic rolls, but if one member fails a Resolve test to resist fear or mental intrusion, all members test at +1 DV.
\end{itemize}

\textit{Note:} These orders are not factions in the usual sense. They have no grand agenda. They simply survive, train, and occasionally produce a wanderer who walks the roads with nothing but a bell, a stone, or a breath.

\vspace{0.5cm}
\hfill --- \textit{The Grey Wanderer, from a conversation at the Dusty Page}

\section{The Thought Trade -- The Economics of the Mind}
\lettrine{I}{n} the shadowed exchanges of Silkstrand and the damp cellars of Zakov, a new currency circulates: the \textbf{thought itself}.

\textbf{Memory Merchants} store sensitive information in living minds, charging a skilled craftsman's daily wage for an hour of secure storage. Dangerous memories---those of murder, heresy, or Patron contact---cost tenfold.

\textbf{Calculation Circles} of mathematical psions serve as living abacuses for engineering projects, performing in hours what takes traditional scholars months. They charge in \textit{Mental Strain} as much as coin, requiring recovery chambers and silence.

\textbf{Experience Exchange:} In underground markets, psions trade raw experiences---the memory of a dragon-flight, the sensation of true love, the fear of a battlefield---as equivalents to silver and gold. A memory of flight might fetch fifty pieces; the memory of a first kill might buy a house.

\textbf{Skill Debt:} Psions enter arrangements to provide mental services over time for immediate payment, creating complex obligations tracked not in ledgers of ink but in \textit{Promises of the Mind}. These obligations are enforced not by law but by the guilds of the Silent Orders---and by the fact that a psion who reneges on a mental weight finds their own thoughts strangely... leaky.

\section{Psions as the Sixth Road}
\lettrine{R}{emember} this: The Runekeeper is an agent, the Invoker a gambler, the Cantor a voice, the Summoner a diplomat, the Witch a gardener. The Psion is \textbf{none and all of these}. They are the isolated ones, the self-accounting, the walkers of the Sixth Road who need no Patron's permission and receive no Patron's protection.

They are hunted by the Chain-Lanterns of Ecktoria, licensed by the Synod of Thepyrgos, sealed in vaults by the Aeler, and shunned by the hearth-keepers of Aelaerem. They carry no outward sign. They are accountable only to themselves---and to the Mind's Ledger, which never forgets a weight.

If you would walk this road, know that you walk it alone. Your only focus is your self. Your only backup is your sanity. And when the Mental Strain overflows and the Mind's Ledger closes, you will find that the greatest creditor is the one who has been watching from inside your skull all along.

\vspace{0.5cm}
\hfill --- \textit{The Grey Wanderer}

    \clearpage
\chapter{Last Words from the Grey Wanderer}
    \epigraph{``A book is a door, not a room. You stand at the threshold now, and something on the other side has heard you breathing.''}{\textit{--- The Grey Wanderer}}
    
    You have walked the Six Roads with me. You have seen how the Aeler count their breaths, how the Linn bind their oaths to bells, how the Thepyrgosi stamp their imprimatur on thoughts themselves. You have read the confession pressed between the leaves---the one written in milk and iron by a hand that no longer remembers its own name. You know about the Daughters, and the Hollowed who never blink, and the Cord that winds through flesh and memory alike.
    
    You think this knowledge protects you. It does not. It only teaches you the shape of the teeth that wait in the dark.
    
    Here is the reckoning, road by road, that I could not speak until you had seen the whole board:
    
    \section*{The Runekeeper's Error}
    You are an agent, not a servant. The Codex is not your shield---it is your chain, and the Thiasos is the weight at the end of it. When your Patron looks through your eyes, do not let them see everything. Keep one Rite unwritten. Keep one page blank. If your Patron ever asks why that page is empty, lie. If they ask twice, run.
    
    \section*{The Invoker's Math}
    You carry four coins, but the fifth pocket is always growing heavier. You think you gamble with knowledge, that your respect for the ritual keeps you safe. It does not. The Patrons do not care about your etiquette; they care about the \textit{asymmetry} of the weight. The moment you Crack the Seal in desperation, you have stopped being a gambler and become the collateral. Polish your Symbols weekly. Polish them daily if you have stared into the dark.
    
    \section*{The Cantor's Touch}
    Your voice is not yours anymore. It belongs to the song, and the song belongs to the silence between the notes. The corruption is not a punishment---it is a \textit{translation}. You are being rewritten into a dialect the Patrons speak fluently. When you find feathers in your hair, or shadows that lag behind your heels, do not sing for three days. If the music continues without you, break your instrument. If the music continues still, break your fingers. It is cheaper than what comes next.
    
    \section*{The Witch's Thread}
    You have read about the Chain. You know the soup is free. But understand: the Daughters are not heretics. They are \textit{maintenance}. The world requires certain memories to be lost, certain names to be unraveled, certain children to be hollowed out so that the walls between the living and the eaten remain standing. If you find a red thread around your wrist and do not remember tying it, do not cut it. Find a threshold. Pour salt. Speak your true name backward until the thread goes slack. If it tightens instead, you have already been collected.
    
    \section*{The Psion's Silence}
    I am not one of you. I have said this. But I have walked beside your kind when the bell rings three times too fast. Your mind is a fortress with no gates---safe until it isn"t, and when it cracks, it crumbles inward, not outward. The Synod fears you because you answer to no ledger they can audit. But you answer to the Silence, the Great Silence that once swallowed thought itself. When you feel the edges of your self growing thin, when you begin to see the geometry of the air, stop. The Silence is not emptiness. It is appetite.
    
    \section*{The Summoner's Door}
    The Leash has two ends, and the spirit holds the other. You believe you command the dead, the fey, the demons, the things that wear the shapes of angels. You do not. You \textit{rent} them, as the Invoker rents power, but your landlords are hungrier and less patient. When the Leash ticks, do not look at your bound spirit. Look at your own shadow. If it moves independently, the door you opened was never yours to close. The Red Weather follows horse-thieves, yes, but it also follows summoners who forget to name their price in advance.
    
    \section*{What Waits at the Crossroads}
    You have read about the Green Gate, the Ninth Walk, the Silent Stair. These are not metaphors. They are \textit{infections} in the skin of the world, places where the Six Roads knot together and the map stops being honest. If you find yourself at such a place, do not make a wish. Do not ask a question. Offer a truth you have never spoken and a lie you have never told, and walk away without looking back. The Wanderer you see at the crossroads is not always me. Sometimes it is the thing that ate the last traveler who looked too long.
    
    \section*{Final Counsel}
    Keep your components dry. Keep your voice warm if you have one, and cold if you must. Keep your mind your own---guard it as the Aeler guard their Deep Vaults, with bells and silence and the willingness to seal the door from the outside. And keep your wit sharper than any Rite, any Symbol, any song. The Patrons can forgive a missed payment. They can forgive a broken covenant. They cannot forgive being outthought by a mortal who remembers that power is, at its root, a conversation---and conversations can be walked away from.
    
    The book is closed. The door is open. I will be at the inn, or I will not. But the roads remain, and they are hungry.
    
    \vspace{1cm}
    \hfill --- \textit{The Grey Wanderer}

    \appendix

    % ======================================================================
    % APPENDIX A: THE SIX TEMPERAMENTS
    % ======================================================================
    % Colors
\definecolor{headerblue}{RGB}{25,30,120}
\definecolor{sidebarblue}{RGB}{40,50,150}
\definecolor{lightgray}{RGB}{240,240,240}
\definecolor{darkgray}{RGB}{80,80,80}
\definecolor{tablegray}{gray}{0.9}


\section*{The Debt of the Sleeping Master}
\addcontentsline{toc}{section}{Appendix: The Debt of the Sleeping Master}

\epigraph{"You think you know someone after thirty years. You share bread, you spill blood, you count each other's breaths. Then one night the body stops---and the mind, that clever thief's mind, runs ahead to hide the one thing you never asked. The door is open. The threshold is warm. What will you find on the other side?"}{\textit{--- The Grey Wanderer, at a wake that turned into a heist}}

\subsection*{A Note on Obishaal and the Ways Between}

The Element of Obishaal is the sweet water that drowns, the dream from which you do not wake, the door between the last breath and the next. In the Grey Wanderer's taxonomy, it stands opposite Life (Fire) and adjacent to Death (Dreams). Free casters who reach unwarily into Obishaal find themselves walking corridors of their own forgotten memories, pursued by shapes that wear their faces. The trained psion, however, can step into Obishaal as a thief steps through a window---if they have a threshold, a witness, and a name to pay.

Yet Obishaal is not a private realm. The **Ways Between**---those phantom bridges, shifting staircases, and crossroads that lead where no road should---were built upon Obishaal like villages upon a river's floodplain. Every traveller who ever walked a hidden path, every shortcut that should not exist, has brushed against this dream-realm. As a result, Obishaal accepts **any** traveller, psionic or otherwise, as long as someone opens the door.

\begin{framed}
\noindent \textbf{For the Heist:} Only one character needs the Psionics skill and the \textit{Telepathy} Art to act as the **anchor**. They pull the rest of the party into the dream (1 Mental Strain per passenger). Non-psions cannot use psionic Arts inside Obishaal, but they can see, speak, fight, sneak, and interact with the environment normally. The anchor pays the cost; the threshold holds.
\end{framed}

The students of Master Kai had all three---a psion, witnesses, and a name.

\begin{framed}
\noindent \textbf{Design Note:} This adventure uses the \textbf{Obishaal} dream-realm mechanics. Treat all scenes as \textbf{Echo} (memory) or \textbf{Veil} (current dream) layers. Each PC carries a \textit{Memory Debt Timer } [4] that ticks every time they use an intrusive psychic Art (Telepathy, Psychic Assault, Precognition). When full, they are pulled into a \textit{Burden}---a personal regret manifests as a hostile dream-creature. The party cannot leave until the Master's secret is named aloud \textit{and} the weight is acknowledged.
\end{framed}

\section{The Threshold: A Body That Breathes But Does Not Wake}

The room smells of old paper, dried herbs, and the faint copper tang of blood. Master Kai lies on a pallet of grey wool, eyes half‑open, lips moving in a language none of the students recognize. The breath is shallow; the pulse is a drum that beats once, then waits, then beats again. The Hollowed---a small grey bird with sewn‑shut eyelids---perches on the headboard, unblinking.

The students have been summoned by a letter delivered by a mute Tulkani courier. The letter is two lines:

\begin{quote}
\emph{"The lock is my skull. The key is my name. Take what I could not speak."}
\end{quote}

Below, a map of the mind---a sketch of a nine‑towered castle drawn on a scrap of leather that smells of salt and dream‑milk.

\subsection{The Setup}

The party must decide: enter Master Kai's mind (Threshold: the pallet; witnesses: each other; price: one memory of comfort with the Master, which will be lost forever). Once in Obishaal, they will appear in a **memory version** of the Academy's meditation hall. From there, they must descend through nine chambers, each guarded by a Burden that represents the Master's hidden secret.

The secret is not a treasure. It is a name---the name of a woman the Master swore to protect and then abandoned during the Ashaani purges. That unspoken vow has grown into a **Spirit Anchor** that keeps the Master from passing. To free them, the party must speak the name aloud at the heart of the dream.

\section{Eight Premade Characters}

\subsection*{1. Rowan --- the Grifter (Body 2, Wits 3, Spirit 3, Presence 2)}
\textbf{Psionics} 3 (Telepathy, Empathic Manipulation). \textbf{Talents:} Thought Thief, Psychic Reservoir.  
\textbf{Strings:} "Master's spare key," "Obligation to the Tulkani courier."  
\textbf{Burden:} "I once read a client's memory of their child's face and sold it to a rival."  
\textbf{Look:} Slicked hair, cracked leather gloves, smile that never reaches the eyes.

\subsection*{2. Sable --- the Scoundrel (Body 3, Wits 2, Spirit 3, Presence 2)}
\textbf{Psionics} 2 (Telekinesis, Psychic Assault). \textbf{Talents:} Telekinetic Mastery, Mental Fortress.  
\textbf{Strings:} "Stole the Master's meditation bell," "Owe a favor to the Port Shed dockmaster."  
\textbf{Burden:} "I crushed a guard's windpipe with my mind. He had a daughter."  
\textbf{Look:} Crooked nose, knuckle tattoos showing broken chains, always leans on a wall.

\subsection*{3. Lyra --- the Archivist (Body 2, Wits 4, Spirit 2, Presence 2)}
\textbf{Psionics} 3 (Clairvoyance, Precognition). \textbf{Talents:} Focused Mind, Echo Dampener.  
\textbf{Strings:} "Hidden map of the Academy's dream‑vaults," "Promise to a dead sister."  
\textbf{Burden:} "I saw a death coming and said nothing. The death came."  
\textbf{Look:} Ink‑stained fingers, half‑moon spectacles, leather satchel full of blank pages.

\subsection*{4. Jax --- the Brute (Body 4, Wits 2, Spirit 2, Presence 2)}
\textbf{Psionics} 2 (Telekinesis, Mind Shield). \textbf{Talents:} Psionic Resilience, Expanded Resonance.  
\textbf{Strings:} "Knucklebone from a fallen mentor," "Owe lifedebt to the party."  
\textbf{Burden:} "I used telekinesis to pull a man into a press. He was unarmed."  
\textbf{Look:} Shaved head, thick neck, brass knuckles etched with calming runes.

\subsection*{5. Mira --- the Whisper (Body 1, Wits 3, Spirit 4, Presence 3)}
\textbf{Psionics} 4 (Telepathy, Empathic Manipulation, Mind Shield). \textbf{Talents:} Mental Fortress, Contamination Control.  
\textbf{Strings:} "Knows the Master's birth‑name," "Feather from the Hollowed bird."  
\textbf{Burden:} "I made a beggar believe he was king for a week. When he woke, he threw himself from a bridge."  
\textbf{Look:} Grey robes, hair shorn in mourning, carries a silver tuning fork.

\subsection*{6. Thorn --- the Needle (Body 2, Wits 4, Spirit 3, Presence 1)}
\textbf{Psionics} 3 (Psychic Assault, Precognition). \textbf{Talents:} Psychic Vampire, Memory Palace.  
\textbf{Strings:} "Blood debt to a Vhasian crime family," "One use of the Master's true name."  
\textbf{Burden:} "I assassinated a woman by erasing her memory of how to breathe. She was a mother of three."  
\textbf{Look:} Wraparound goggles, surgeon's gloves, a belt of scalpels that double as psychometric anchors.

\subsection*{7. Dove --- the Veil (Body 2, Wits 3, Spirit 2, Presence 4)}
\textbf{Psionics} 3 (Empathic Manipulation, Mind Shield, Telepathy). \textbf{Talents:} Mental Fortress, Psychic Synergy.  
\textbf{Strings:} "Shield‑token from the Order of the Shared Breath," "Forged identity papers."  
\textbf{Burden:} "I broke a man's will to love so he would leave his wife. He did. Both died."  
\textbf{Look:} White silk scarf that hides a branded cheek, always carries a cup of honeyed tea.

\subsection*{8. Rook --- the Turncoat (Body 3, Wits 3, Spirit 3, Presence 1)}
\textbf{Psionics} 3 (Clairvoyance, Astral Projection, Telekinesis). \textbf{Talents:} Mind Walker, Echo Dampener.  
\textbf{Strings:} "Map of the Academy's secret exits," "Letter from a Chain‑Lantern inquisitor."  
\textbf{Burden:} "I betrayed a fellow student to save myself. They still don"t know it was me."  
\textbf{Look:} Military jacket with the insignia torn off, one ear missing, carries a compass that points to the nearest exit.

\section{The Dream Heist: Framework}

\subsection{Timers}

\begin{itemize}
    \item \textbf{The Master's Breath [8] ---} Ticks every scene (or when a PC marks Mental Strain). When full, the Master's body stops breathing. The dream collapses; all PCs mark 2 Harm and gain the \textit{Haunted} condition.
    \item \textbf{Weight's Approach [6] ---} Ticks each time the party speaks a lie, fails a Resolve test, or damages a dream‑construct. When full, the final Burden (the abandoned woman's spirit) appears early and attacks, and the \textbf{Spector} (see below) manifests immediately.
    \item \textbf{Name Unspoken [4] ---} The party must discover the forgotten name (Finley). Each successful Investigation or Telepathy roll in a memory chamber fills one segment. When full, the name is known and can be spoken at the heart.
\end{itemize}

\subsection{New Element: The Spector --- The Parting of the Veil}

\epigraph{"Some call it the Angel of Death's touch. Others, the Final Breath. A few whisper it is the Hollow wearing a different face. What matters is this: when it walks, the dream holds its breath. And so should you."}{\textit{--- The Grey Wanderer}}

The Spector is not a creature to be fought. It is a **looming presence** tied directly to the \textit{Master's Breath} timer. When the timer ticks, the Spector may appear at the threshold of the current chamber---a tall, featureless figure in grey, with long fingers that trail threads of fading light. It does not speak. It does not rush. It \textit{walks} toward the party at a slow, inexorable pace.

\begin{itemize}
    \item \textbf{Appearance:} On any tick of the \textit{Master's Breath} timer (including the initial tick when entering a new chamber), the GM may choose to have the Spector manifest at the chamber's entrance (the way they came). It advances one zone per round.
    \item \textbf{Annihilation:} If the Spector touches a character, that character is immediately ejected from the dream, suffers \textbf{2 Harm (Stress)} and loses \textbf{1 permanent Memory} (a fact, a skill point, or a String of the GM's choice). They cannot re-enter the dream.
    \item \textbf{Avoidance:} The Spector cannot be harmed or reasoned with. It can only be avoided:
    \begin{itemize}
        \item \textbf{Sneak:} A group Stealth test (Wits + Stealth) vs. DV 4. On success, the party moves to the next chamber without the Spector following. On a partial, the Spector follows but doesn"t attack this scene. On a miss, it attacks immediately.
        \item \textbf{Distraction:} A character may sacrifice a \textbf{String} or a \textbf{Memory} (declare a minor fact about their past that they forget) to delay the Spector for one full scene.
        \item \textbf{Ritual:} Reciting the name of the \textit{Pale Shepherd} or \textit{the Hollow} (player's choice) causes the Spector to pause for 1d4 rounds, but the GM gains \textbf{2 SB} (the entity grows curious).
    \end{itemize}
    \item \textbf{If the Master's Breath timer fills:} The Spector immediately appears in the final chamber and cannot be avoided. The party must speak the forgotten name \textit{while} evading it---each round, one character may act; the Spector moves one zone closer.
\end{itemize}

The GM may decide the Spector is the Hollow, the Pale Shepherd, or an unnamed death‑spirit. The players may never know. This ambiguity is by design.

\subsection{New Element: Traps --- Defense Mechanisms of the Mind}

The Master was a psion and a thief. His mind bristles with **defense techniques**---some the party learned in their training (Recognized), others utterly alien (Alien Techniques).

\begin{description}
    \item[Recognized Techniques (Familiar Pain):] These traps use exercises the Master actually taught. The party gains \textbf{+1 die} to disarm them, but disarming requires confessing a minor weakness (e.g., "I never mastered this technique because I was lazy").
    \item[Alien Techniques (Unseen Wards):] These traps are unknown. They impose \textbf{+1 DV} to disarm. On a miss, the trap doesn"t just harm---it reveals a \textbf{hidden shame} of the Master (the GM narrates a small, painful memory).
\end{description}

\textbf{Example traps (place in any chamber):}
\begin{itemize}
    \item \textbf{The Memory Latch (Recognized):} A door sealed by a mental lock that requires a specific memory of failure to open. \textit{Disarm:} Wits + Psionics (DV 3, +1 die).
    \item \textbf{The Dream Snare (Alien):} Invisible threads that pull the character into a loop of a traumatic event (the Master"s). \textit{Detect:} Clairvoyance (DV 4). \textit{Disarm:} Telekinesis (DV 5) or Precognition (DV 4). On failure, stun 1 round + Master's Breath ticks +1.
    \item \textbf{The Geas Glyph (Recognized):} Compels honesty. A character who crosses must answer one truth about their relationship with the Master. \textit{Disarm:} Empathic Manipulation (DV 4, +1 die). On miss, the question is asked anyway, and the GM gains 1 SB.
    \item \textbf{The Hollow Floor (Alien):} A pit into a forgotten argument between the Master and Finley. \textit{Detect:} Notice (DV 5) or Astral Projection. \textit{Cross:} Telekinesis or Athletics (DV 4). On failure, fall + witness argument (1 Fatigue, GM gains 2 SB).
\end{itemize}

\subsection{New Element: Shadows of Life --- Perception‑Triggered Attackers}

Shadows of Life are echoes of people the Master knew---loving or resentful. They are invisible and harmless as long as no one \textit{perceives} them directly. Once noticed, they become hostile.

\begin{itemize}
    \item \textbf{Trigger:} The GM may spend \textbf{2 Story Beats} (moderate spend) to declare that a Shadow is nearby. It attacks next round unless prevented.
    \item \textbf{Perception:} A character who succeeds on \textbf{Wits + Notice} (DV 3) spots the Shadow before it attacks, gaining a free action.
    \item \textbf{Avoidance:} Shadows cannot be fought directly. They can be outrun (group Stealth DV 4) or appeased by speaking a truth the Shadow "wants to hear" (GM discretion). 
    \item \textbf{If attacked:} Shadows have \textbf{Thought Thief} stats (Mental Harm; +1 SB per hit). Killing one advances \textit{Weight's Approach} by 1.
\end{itemize}

\subsection{Threshold Journey (Nine Chambers)}

\begin{enumerate}
    \item \textbf{The Meditation Hall (Echo).} The party must relive a training session. A Burden appears as a younger student who failed. \textit{Traps:} Memory Sink (Recognized: pit of failure). \textit{Shadows:} Shadow of a jealous rival. \textit{Spector possible.}
    \item \textbf{The Vault of Stolen Secrets (Veil).} Floating memory‑orbs. One contains the first clue: \textit{"Not treasure. A woman. The Ash‑Fenn market."} \textit{Traps:} Dream Snare (Alien). \textit{Shadows:} Shadow of a student erased from records.
    \item \textbf{The Anxious Stair (Flow).} Steps that grow teeth. Each PC confesses a regret aloud. After each confession, a Doubt Wraith (Psychic Assault stats) attacks. \textit{Spector may appear.}
    \item \textbf{The Court of the Witness (Echo).} Mirror‑hall where Burdens manifest. After speaking their Burden's name, the mirrors shift into a labyrinth. \textit{Traps:} Mirror Labyrinth (Alien, disorient). \textit{Shadows:} Shadows of those the Master judged unfairly.
    \item \textbf{The Iron Well (Veil).} A deep shaft with a locket at the bottom. A Burden Guardian (Poltergeist stats) protects it. \textit{Traps:} Geas Glyph on the rim (Recognized). 
    \item \textbf{The Harrow (Flow).} A corridor of grey‑thread thorns. Each PC sacrifices one String. \textit{Traps:} String Snare (Alien, steals extra String). \textit{Shadows:} Shadow of a betrayed ally.
    \item \textbf{The Forgotten Oath (Echo).} A courtroom where the Master is judged by their own conscience. The party must defend them with acts of kindness. \textit{Traps:} Geas Glyph (Recognized) on the defendant's chair. \textit{Shadows:} Shadow of the Master's mother.
    \item \textbf{The Hollowed Garden (Veil).} A grove. The faceless woman asks, "Why did he leave?" Empathic success reveals the answer. On failure, Thorn‑Kin attack. 
    \item \textbf{The Heart (Obishaal pure).} Bare stone altar. Speaking the name \textit{Finley} fills the Name Unspoken timer. The Master's spirit rises. If \textit{Master's Breath} is 1 or less, the Spector appears immediately. The party must speak the name while evading it.
\end{enumerate}

\subsection{Resolution}

The party wakes. The Master's body is still. The Hollowed bird opens its sewn eyes, chirps once, and flies out the window. In the Master's hand is a slip of paper: \textit{"You found it. The treasure was my shame. Now I can rest. The Academy is yours---if you can keep it."}

Each PC clears one Obligation or Corruption segment of their choice (the Master's death‑bed pardon). They also gain a permanent \textit{Dream‑Echo}: once per session, they may reroll a failed Resolve test in Obishaal or against fear.

\epigraph{"They never spoke of the heist again. But sometimes, in the small hours, one of them would wake with a name on their lips---Finley---and they would pour a cup of water on the ground, mutter "paid," and go back to sleep. That is the closest our kind comes to prayer."}{\textit{--- The Grey Wanderer, closing a bottle of Sidhi port}}


    \chapter*{The Six Temperaments}
    \addcontentsline{toc}{chapter}{The Six Temperaments}
    
    \epigraph{``Before you carve your first Rite, hum your first Song, or swallow your first coin, ask yourself: do you want to serve, gamble, sing, tend, think, or bind? The answer determines which weight you will drown in.''}{\textit{--- The Grey Wanderer}}
    
    Choosing a Road is less about mechanics than about \textit{temperament}. Below, the Six Roads are matched to player psychology, narrative fantasy, and---for the Runekeeper---the specific Patron that transforms the class from a "priest" into a Paladin, Druid, Artificer, or Inquisitor.
    
    \section*{Which Road Suits Your Nature?}
    
    {\small
    \begin{longtable}{@{}p{2.2cm} p{4.8cm} p{4.5cm}@{}}
    \toprule
    \textbf{Road} & \textbf{Best For Players Who\ldots} & \textbf{Core Fantasy} \\
    \midrule
    \textbf{Runekeeper} & 
    Enjoy structured power with clear costs; like being an "agent" of cosmic forces; want lore-heavy obligations that shape identity; prefer preparation and ritual precision. & 
    \textit{"I do not merely cast fire. I am the ledger of the Everflame, and my Codex is law."} \\
    \midrule
    \textbf{Invoker} & 
    Love gambling mechanics and risk/reward balancing; want versatility without commitment; enjoy "desperation" casting and juggling multiple patrons; like feeling clever for breaking rules. & 
    \textit{"I carry four Symbols. One is cracked. Two are angry. None of them trust me---and that's how I like it."} \\
    \midrule
    \textbf{Cantor} & 
    Prefer social/performance scenes to book-keeping; want power that is literally at hand (voice/body); enjoy tragic corruption arcs; like being public, loud, and hard to forget. & 
    \textit{"I need no Codex. My larynx is older than any tree, and the Song is writing me as we speak."} \\
    \midrule
    \textbf{Witch} & 
    Prefer subtlety and preparation over flashy casting; enjoy "folk horror" and domestic magic; like being underestimated; want to solve problems with knots, thresholds, and old names rather than force. & 
    \textit{"The hedge is what keeps the wolves from the flock. I am the one who tends the hedge."} \\
    \midrule
    \textbf{Psion} & 
    Want internal struggle rather than external debts; prefer mind games and subtlety; dislike carrying obvious magical gear; enjoy the isolation of "untouchable" power. & 
    \textit{"I carry no Symbol. My ledger is my skull, and it comes due in headaches, not obligations."} \\
    \midrule
    \textbf{Summoner} & 
    Enjoy tactical "pet" management and action economy; like contracts, diplomacy with monsters, and bargaining; prefer spirits to spells; want a "friend" that might eat them. & 
    \textit{"I asked the drowned captain for directions. He asked for a memory in return. We shook on it."} \\
    \bottomrule
    \end{longtable}
    }
    
    \section*{The Runekeeper's Many Faces}
    
    The Runekeeper is not a single class but a \textit{chassis}---the Patron supplies the theme, and the Rites supply the mechanics. If you want to play a "Paladin," swear to Mykkiel. If you want a "Druid," swear to Grimmir. Below is a sampling of how Patrons map to classic archetypes using the full roster from the grimoires.
    
    {\small
    \begin{longtable}{@{}p{2.8cm} p{2.8cm} p{4.5cm} p{2cm}@{}}
    \toprule
    \textbf{Patron} & \textbf{Archetype} & \textbf{Signature Combat Style} & \textbf{Key Rite} \\
    \midrule
    \textbf{Mykkiel} (Law) & 
    Paladin / Templar / Inquisitor & 
    Smite foes with cleansing fire; consecrate ground; bind oaths; immune to fear. & 
    \textit{Final Admonition} \\
    \textbf{Grimmir} (Wild) & 
    Druid / Ranger / Warden & 
    Summon thorns; speak with beasts; heal blight; control terrain with roots. & 
    \textit{Thornveil} \\
    \textbf{Clockwork Monad} & 
    Artificer / Engineer / Savant & 
    Optimize gear mid-fight; iterative improvements; self-repairing mechanisms. & 
    \textit{Self-Optimizing Mechanism} \\
    \textbf{Inquisitor Prime} & 
    Witch Hunter / Exorcist & 
    Dispel magic; brand enemies; resist supernatural influence; burn corruption. & 
    \textit{Cleansing Light} \\
    \textbf{Khemesh} (Depth) & 
    Warlock / Aberrant Mind & 
    Crushing pressure; silence enemies; water-based control; sanity damage. & 
    \textit{Maw Opens} \\
    \textbf{Oath of Flame} & 
    Crusader / Dawn-Knight & 
    Radiant smites; healing lay-on-hands; vow-based defense; anti-undead. & 
    \textit{Radiant Smite} \\
    \textbf{The Witness} & 
    Oracle / Truth-Speaker & 
    Reveal hidden foes; compel honesty; bind with immutable record; counter-illusion. & 
    \textit{Echoing Truth} \\
    \textbf{The Carrion-King} & 
    Necromancer / Compost-Witch & 
    Accelerate decay; harvest endings; fertility from death; blight enemies. & 
    \textit{Fertile Death} \\
    \textbf{Raéyn} (Sea) & 
    Storm-Caller / Tide-Witch & 
    Lightning strikes; water manipulation; sailor buffs; drowning curses. & 
    \textit{Storm-Queen's Hand} \\
    \bottomrule
    \end{longtable}
    }
    
    \noindent\textit{Note:} A Runekeeper of \textbf{Xhak"Thul} plays like a Storm-Barbarian, while one sworn to \textbf{The Sealed Gate} plays like an Abjuration Ward-Master. The Codex and Thiasos remain constant; the \textit{flavor} of power changes with the Patron's domain.
    
    % ======================================================================
    % APPENDIX B: THE BATTLEFIELD
    % ======================================================================
    
    \chapter*{The Battlefield}
    \addcontentsline{toc}{chapter}{The Battlefield}
    
    \epigraph{``Every Road can kill you. The question is whether you kill them with fire, song, weight, or a very polite spirit.''}{\textit{--- The Grey Wanderer}}
    
    No Road is limited to a single role, but each has natural affinities. The tables below map the Six Roads to the four classic combat functions---Damage, Healing, Utility, and Control---so you can build a party that covers all bases, or double down on your Road's specialty.
    
    \section*{Combat Roles by Road}
    
    {\small
    \begin{longtable}{@{}p{1.8cm} p{3.2cm} p{3.2cm} p{3.2cm} p{2.2cm}@{}}
    \toprule
    \textbf{Road} & \textbf{Damage (Striker)} & \textbf{Healing / Support} & \textbf{Utility / Control} & \textbf{Risk Type} \\
    \midrule
    \textbf{Runekeeper} & 
    \textit{Thematic nukes} (e.g., \textit{Final Admonition}, \textit{Maw Opens})---high burst, single-target or small area, gated by Obligation. & 
    \textit{Rites of Cleansing}---remove Conditions, restore Fatigue via \textit{Lay on Hands} (Flame), or HEAL wounds via \textit{Fertile Death} (Carrion-King). & 
    \textit{Wards \& Seals}---create zones of sacred ground (\textit{Consecrated Barrier}), bind oaths that penalize breach, or dispel magic. & 
    Obligation overflow; Patron intrusion. \\
    \midrule
    \textbf{Invoker} & 
    \textit{Crack the Seal}---emergency high-power Rites, but costly. Best for "oh no" moments, not sustained DPS. & 
    \textit{Borrowed resilience}---can invoke healing Symbols, but slower and riskier than Runekeepers. & 
    \textit{Versatility}---carries 4 Symbols; can switch from Aveh (escape) to Mykkiel (binding) mid-fight. Best for filling gaps. & 
    Symbol fracture; Obligation spike. \\
    \midrule
    \textbf{Cantor} & 
    \textit{Sonic assaults}---area denial and morale breaks via \textit{Hymn Against Dread}; single-target "shatter" via \textit{Single True Note}. & 
    \textit{Battlefield maintenance}---\textit{Hymn of Restoration} removes Fatigue; \textit{Chant of Purification} clears poison/disease. & 
    \textit{Social control}---\textit{Hooked Chorus} compels action; \textit{Circling Cup} creates zones of enforced revelry/peace. & 
    Corruption accumulation; voice loss. \\
    \midrule
    \textbf{Witch} & 
    \textit{Curses}---slow, invisible, and devastating (e.g., \textit{Red Weather}, \textit{Snaring Filament}). Poor burst, excellent attrition. & 
    \textit{Hearth remedies}---\textit{Warm Hand}, \textit{Cup and Salt}; requires preparation and components. Best for after-combat triage. & 
    \textit{Threshold mastery}---lock doors with \textit{Binding Knot}, create sanctuary with \textit{Hearth-Law}, or scout via \textit{Weaver's Glance}. & 
    Promise Timers; neighborly obligation. \\
    \midrule
    \textbf{Psion} & 
    \textit{Psychic assault}---ignores armor; targets Spirit directly. High single-target burst via \textit{Psychic Assault}. & 
    \textit{Biofeedback}---heals by converting Fatigue; can stabilize but not "burst heal" like Flame Rites. & 
    \textit{Information warfare}---\textit{Telepathy}, \textit{Clairvoyance}, \textit{Precognition} for battlefield control and ambush prevention. & 
    Mental Strain overflow; sanity scars. \\
    \midrule
    \textbf{Summoner} & 
    \textit{Spirit damage}---spirits act independently, providing action economy; can flank and harass while summoner controls. & 
    \textit{Spirit aid}---some spirits (selkies, ancestors) grant buffs or soak Fatigue via \textit{Boon Finesse}. & 
    \textit{Spirits as tools}---\textit{Vein-Spirit} locates minerals; \textit{Wolf Ancestor} tracks; \textit{Drowned} curses ropes. & 
    Leash snap; spirit betrayal. \\
    \bottomrule
    \end{longtable}
    }
    
    \section*{The Free Caster (TAGS) --- The Wildcard}
    
    The \textbf{Free Caster} (using the TAGS system) is not one of the Six Roads, but merits mention. They accommodate \textit{any} role---Damage, Healing, Utility, Control---via improvisation, but always at higher risk than specialists. They are best for players who hate spell lists and love inventing solutions on the fly, accepting that the Backlash may turn their healing spell into a fireball.
    
    \begin{itemize}
        \item \textbf{Damage:} \textit{Burning} + \textit{Area} + \textit{Force} (DV 4)---explosive, risky, and loud.
        \item \textbf{Healing:} \textit{HEAL} + \textit{Purify} (DV 3)---reliable, but generates Story Beats on partial success.
        \item \textbf{Control:} \textit{Bind} + \textit{Stone} + \textit{Wall} (DV 4)---traps and barriers, but costs Fatigue.
    \end{itemize}
    
    \vspace{0.5cm}
    \noindent\textit{Final counsel:} Do not choose a Road based solely on which deals the most Harm. Choose the Road whose \textit{failure state} interests you. A Runekeeper who falls to Obligation becomes a monster for a season; a Cantor who fills their Corruption Timer becomes a legend with feathers in their hair. That is the story you are signing up to tell.
    

\end{document}